% Paper 34: GeoVac as a Two-Layer Framework
% Projection Taxonomy and the Empirical Matches Catalog
% GeoVac Project, May 2026
% v0.1 (skeleton) -- Standalone observation paper.  Names the
% framework's two-layer architecture (bare graph + projection
% family), gives a per-projection dictionary tagging each by
% physical variables introduced, dimensions introduced, and
% transcendental signature, and consolidates the running list of
% machine-precision matches as a tagged catalog.

\documentclass[11pt]{article}
\usepackage{amsmath,amssymb,amsthm}
\usepackage{booktabs}
\usepackage{array}
\usepackage{longtable}
\usepackage[margin=1in]{geometry}
\usepackage{hyperref}

\newtheorem{observation}{Observation}
\newtheorem{definition}{Definition}
\newtheorem{remark}{Remark}
\newtheorem{prediction}{Prediction}
\newtheorem{theorem}{Theorem}

\title{GeoVac as a Two-Layer Framework:\\
Projection Taxonomy and the Empirical Matches Catalog}
\author{GeoVac Project}
\date{May 2026}

\begin{document}
\maketitle

% ===================================================================
\begin{abstract}
% ===================================================================
The GeoVac framework decomposes into two structurally separate
layers.  \textbf{Layer~1} is a bare combinatorial graph: quantum-number
labels $(n,\ell,m,\ldots)$, integer eigenvalues, rational matrix
elements, no physics.  \textbf{Layer~2} is a finite family of named
projections from the bare graph onto physical settings.  Every
physical variable, every dimensional unit, and every transcendental
constant that appears in any GeoVac calculation enters through
exactly one projection, and is fully traceable to it.

We catalogue nineteen projection mechanisms and tag each one along
three independent axes: (a)~which physical variables it introduces
(charge $Z$, coupling $\alpha$, energy scale $\hbar\omega$, internuclear
distance $R$, gauge coupling $\beta$, cutoff $\Lambda$, $\ldots$);
(b)~which physical dimension or unit it attaches to graph data
(energy, length, dimensionless, solid angle); and (c)~its
transcendental signature in the sense of Paper~18 (rational, $\pi^{2k}$,
half-integer Hurwitz, odd~$\zeta$, Dirichlet~$L$, $\ldots$).  The
projection axis is a structural sibling of Paper~18's transcendental
axis: each transcendental tier corresponds to a specific projection
type, and reading the dictionary in either direction recovers the
other.

We then assemble the empirical matches catalogue: a tagged table of
GeoVac quantities that match standard known physics values to
machine precision (or stated uncertainty).  Each row is annotated
by projection(s), variable(s), dimension, transcendental class, and
match precision.  Reading the catalogue surfaces a falsifiable
structural prediction:
\begin{prediction}[Error compounds with projection depth.]
GeoVac matches that involve a single projection are exact at
machine precision (e.g.\ Pauli counts, hydrogen Bohr energies,
Dirac fine-structure formula, $D^2$ heat-kernel coefficients).
Matches that chain multiple projections accumulate residual
error roughly proportional to the number of projections involved.
\end{prediction}
The Lamb shift example (three chained projections, 2--3\% residual
error) is consistent with the prediction; the catalogue identifies
several other multi-projection chains for future test.

This paper introduces no new computation.  It is a structural
consolidation of the framework's self-presentation: language for
discussing where physics enters, what carries dimensions, and which
transcendentals are forced by which projection.
\end{abstract}

% ===================================================================
\section{Introduction: The Two-Layer Thesis}
\label{sec:intro}
% ===================================================================

The recurring difficulty in discussing GeoVac is that the framework
sits between two languages.  In one language, GeoVac is a sparse
graph Hamiltonian framework for quantum chemistry: nodes labeled by
quantum numbers, edges from selection rules, eigenvalues that
predict atomic and molecular spectra (Paper~14, Paper~17, Paper~20).
In the other language, GeoVac is a discrete spectral triple in the
Marcolli--van~Suijlekom lineage of graph-spectral-triple
constructions (Paper~32), with a heat-kernel expansion that
reproduces the QED $\beta$-function at one loop (Paper~28) and
recovers seven of eight QED selection rules with a single Hopf-base
calibration constant $1/(4\pi)$ per loop (Paper~33).

These two languages are not in conflict, and we have not been able
to choose between them, because they are descriptions of two
different layers of the same construction.  This paper names the
layers and the bridge between them.

\begin{definition}[Layer~1: the bare graph]
\label{def:layer1}
The Layer~1 object is the dimensionless combinatorial graph.  Nodes
are quantum-number labels $(n, \ell, m, \kappa, m_j, \ldots)$.
Eigenvalues of the bare graph Laplacian are integers ($n^2 - 1$ on
the Coulomb $S^3$ graph, $n + 3/2$ on Dirac $S^3$, $N$ on the
Bargmann--Segal $S^5$ graph).  Matrix elements of bare graph
operators (Casimir, ladder, adjacency) are rational, integer, or
algebraic over $\mathbb{Q}[\sqrt{2}, \sqrt{3}, \sqrt{6},\ldots]$
(from Wigner~$3j$ algebra); $\pi$ does not appear in any Layer~1
quantity (Paper~24's $\pi$-free certificate; the same applies on
the Coulomb side).  Layer~1 carries no physical units, no
coupling constants, no observable energies.
\end{definition}

\begin{definition}[Layer~2: the projection family]
\label{def:layer2}
The Layer~2 object is a finite family of named maps from the bare
graph to physical settings.  Each projection is a structural
construction with a name in the literature (Fock 1935 conformal
projection; Hopf bundle; Bargmann--Segal transform; Connes--Chamseddine
spectral action; Sturmian basis at exponent $\lambda = Z/n$;
Camporesi--Higuchi spinor lift; Wigner~$3j$ angular coupling;
Wigner~$D$-matrix rotation between molecular centers; Wilson
plaquette; molecule-frame hyperspherical separation; vector-photon
promotion; Mellin/heat-kernel regularization).  Each projection
adds (i)~specific physical variables, (ii)~specific physical
dimensions, and (iii)~specific transcendental content to the bare
graph data on which it acts.
\end{definition}

\begin{observation}[Single-projection origin of physical content]
\label{obs:single_origin}
Every coupling constant, energy scale, length scale, and
transcendental constant that appears in any GeoVac computation can
be traced to exactly one projection in Layer~2.  The bare graph
does not contain $\alpha$, $Z$, $\hbar\omega$, $R$, $\Lambda$, $\pi$,
or $\zeta(2k)$.  These quantities are introduced by, and only by,
the projection that the calculation invokes.
\end{observation}

The remainder of this paper is the dictionary.

% ===================================================================
\section{Layer 1: The Bare Graph}
\label{sec:layer1}
% ===================================================================

% TODO: This section should cite Papers 0, 7, 24, 22 in detail.
% Brief here; expand from existing material.

\subsection{Quantum-number labelling}
\label{sec:layer1_labels}
Nodes of the bare GeoVac graph are tuples of quantum numbers.  The
specific tuple depends on which natural geometry is in play:
\begin{itemize}
\setlength\itemsep{0pt}
\item Atomic Coulomb $S^3$: $(n, \ell, m_\ell)$ for the scalar sector,
$(n, \kappa, m_j)$ for the Dirac sector.
\item Bargmann--Segal $S^5$ (3D HO): $(N, \ell, m_\ell)$ in $\mathrm{SU}(3)$
$(N, 0)$ irreps.
\item Composed molecular: per-block tuples of the above, with
inter-block edges labeled by the multipole order $L$ of the
cross-center expansion.
\end{itemize}
No physical quantities are required to define these labels.  They
are pure representation-theoretic data of compact Lie groups.

\subsection{$\pi$-free certificate}
\label{sec:layer1_pi_free}
On both the Coulomb $S^3$ Fock graph and the Bargmann--Segal $S^5$
graph, every adjacency matrix entry, every diagonal eigenvalue, and
every CG-coupling matrix element is exact rational or algebraic in
$\mathbb{Q}[\sqrt{2}, \sqrt{3}, \sqrt{5}, \sqrt{6}, \sqrt{10},
\sqrt{15}, \ldots]$ (Paper~24~\cite{paper24}, Paper~28's
graph-native QED certificate).  No $\pi$ appears in Layer~1 by
exact computation.

\subsection{Layer-1 observables}
\label{sec:layer1_obs}
Some observables are computable entirely within Layer~1.  These are
the ``zero-projection'' matches in the catalogue
(Section~\ref{sec:catalog}).  Examples:
\begin{itemize}
\setlength\itemsep{0pt}
\item Pauli term count $N_\text{Pauli} = 11.11 \cdot Q$ (across 28
composed molecules; Paper~14).
\item Magic numbers $2, 8, 20, 28, 50, 82, 126$ (HO + spin-orbit
graph eigenvalue counting; Paper~23).
\item Angular sparsity density $1.44\%$ at $\ell_\text{max}=3$
(Paper~22, potential-independent).
\item All Gaunt selection rules (CG sparsity).
\end{itemize}
These exist on Layer~1 alone.  They have no $\pi$, no $\alpha$, no
energy scale, no transcendentals of any kind, and they match
experiment exactly because there is nothing to project---they are
consequences of compact-Lie-group representation theory acting on
discrete quantum-number labels.

% ===================================================================
\section{Layer 2: The Projection Family}
\label{sec:layer2}
% ===================================================================

We enumerate nineteen projection mechanisms.  For each: source and
target spaces, physical variables introduced, dimensions
introduced, transcendental signature, and existing GeoVac papers
that exercise the projection.  Projections~14 (rest-mass) and~15
(observation / temporal-window) were added in sprint KG / Paper~35
(May~2026); projection~16 (two-body Dirac / Breit retardation) was
added in the precision-catalogue sprint of 2026-05-08, motivated by
the positronium 1S--2S framework-native residual that surfaced as
the sixth (and structurally cleanest) instance of the
multi-focal-composition wall named in CLAUDE.md~\S~2.
Projections~17--19 (nuclear charge-density / Foldy--Friar,
magnetization-density / Zemach, and tensor multipole) were added in
the intra-particle hunting sprint of 2026-05-09, after the audit
memo \texttt{debug/paper34\_v\_base\_unit\_audit\_memo.md} flagged
magnetization-density as a single 17th-projection candidate and the
follow-up hunt established that the three mechanisms are
categorically distinct (different Hamiltonian operators, different
selection rules; documented in
\texttt{debug/intra\_particle\_L\_hunting\_memo.md}).  A boundary
subsection (\S~\ref{sec:layer2_boundary}) names the
internal-graph-derived vs.\ Layer-2-input distinction, with the
external-field hunting memo
(\texttt{debug/external\_field\_hunting\_memo.md}, same date)
verifying that uniform external classical fields (Zeeman, Stark) are
Layer-2 inputs rather than \S~\ref{sec:layer2} projections.

\subsection{Fock conformal projection: $\mathbb{R}^3 \to S^3$}
\label{sec:proj_fock}
\textbf{Source:} continuum Coulomb Hamiltonian $H = -\tfrac{1}{2}\nabla^2 - Z/r$
on $\mathbb{R}^3$.\\
\textbf{Target:} unit $S^3$ via stereographic-like map at focal
length $p_0^2 = -2E$.\\
\textbf{Variables introduced:} $Z$ (nuclear charge), $E$ (energy
eigenvalue).\\
\textbf{Dimension introduced:} energy (in atomic units, Hartree).
The dimensional content is carried by the choice
$m_e = \hbar = e = 1$.\\
\textbf{Transcendental signature:} rational prefactor
$\kappa = -1/16$ (the Rydberg-to-graph-eigenvalue conversion).
$\pi$ enters through $\mathrm{Vol}(S^3) = 2\pi^2$ when subsequent
spectral integrals are performed.\\
\textbf{Used in:} Paper~7 (S$^3$ proof), Papers 13/15/17
(hyperspherical hierarchy), Paper~18 \S~III (Peter--Weyl
realization), all atomic energy computations.

\subsection{Hopf bundle: $S^3 \to S^2 \times S^1$ (locally)}
\label{sec:proj_hopf}
\textbf{Source:} $S^3 \cong \mathrm{SU}(2)$.\\
\textbf{Target:} $S^2$ base $\times S^1$ fibre.\\
\textbf{Variables introduced:} $\alpha$ (the fine-structure constant
treated as a candidate observable in Paper~2's conjecture).\\
\textbf{Dimension introduced:} dimensionless.\\
\textbf{Transcendental signature:} $\pi = \mathrm{Vol}(S^2)/4$
(Hopf measure factor; Sprint~A $\alpha$-PI identification).\\
\textbf{Used in:} Paper~2 (fine-structure conjecture), Paper~25
(Wilson--Hodge gauge structure), Paper~30 (SU(2) Wilson sibling).

\subsection{Bargmann--Segal transform: $\mathbb{R}^3$ HO $\to S^5$ Hardy space}
\label{sec:proj_bargmann}
\textbf{Source:} 3D isotropic harmonic oscillator
$H_{\text{HO}} = \tfrac{1}{2}p^2 + \tfrac{1}{2}\omega^2 r^2$.\\
\textbf{Target:} holomorphic Hardy-space sector of $S^5$, restricted
to $\mathrm{SU}(3)$ $(N,0)$ symmetric irreps.\\
\textbf{Variables introduced:} $\hbar\omega$ (HO frequency).\\
\textbf{Dimension introduced:} energy.\\
\textbf{Transcendental signature:} \emph{none}.  The Bargmann--Segal
graph is bit-exactly $\pi$-free in rational arithmetic at every
finite $N_\text{max}$ (Paper~24).  $\pi$ appears only as
$\mathrm{Vol}(S^5) = \pi^3$ in continuum integration measures, never
in lattice data.\\
\textbf{Used in:} Paper~24 (HO discretization), Paper~31 (universal
vs.\ Coulomb partition).

\subsection{Stereographic / conformal coordinate change}
\label{sec:proj_stereo}
\textbf{Source:} compact sphere ($S^2$, $S^3$, $S^5$).\\
\textbf{Target:} flat-space coordinates ($\mathbb{R}^2, \mathbb{R}^3,
\mathbb{R}^5$).\\
\textbf{Variables introduced:} length scale (radial coordinate $r$).\\
\textbf{Dimension introduced:} length.\\
\textbf{Transcendental signature:} conformal factor producing the
$1/r$ Coulomb potential as coordinate distortion (Paper~7).\\
\textbf{Used in:} Paper~7 (Coulomb-as-coordinate), all calculations
that quote results in physical bohr/Hartree.

\subsection{Sturmian reparameterization at $\lambda = Z/n$}
\label{sec:proj_sturmian}
\textbf{Source:} bare Coulomb $S^3$ graph.\\
\textbf{Target:} same graph relabelled as a Coulomb Sturmian basis
spanning bound states plus a discretized continuum.\\
\textbf{Variables introduced:} $Z$, $n$ (re-used; no new variables).\\
\textbf{Dimension introduced:} preserves existing.\\
\textbf{Transcendental signature:} preserves rationality of
Layer~1.  No new $\pi$ or other transcendentals enter at this step;
calibration-tier transcendentals (e.g.\ Bethe logarithm) enter
through subsequent spectral mode sums.\\
\textbf{Used in:} Papers 8--9, Paper~19, sprint LS-2 (Bethe logarithm
from Sturmian sum at 2--3\%).

\subsection{Connes--Chamseddine spectral action / heat kernel}
\label{sec:proj_spectral_action}
\textbf{Source:} Dirac operator $D$ on bare $S^3$ graph (or
extension).\\
\textbf{Target:} heat-kernel expansion
$\mathrm{Tr}\, f(D^2/\Lambda^2)$ with cutoff function $f$ and scale
$\Lambda^2$.\\
\textbf{Variables introduced:} $\Lambda$ (UV cutoff), $\alpha$ (when
gauge sector included).\\
\textbf{Dimension introduced:} energy via $\Lambda$.\\
\textbf{Transcendental signature:} $\sqrt{\pi} \times \mathbb{Q}$
Seeley--DeWitt coefficients on unit $S^3$ ($a_0 = a_1 = \sqrt{\pi}$,
$a_2 = \sqrt{\pi}/8$); resulting observable coefficients are
$\pi^{2k} \times \mathbb{Q}$ at one loop (T9 theorem, Paper~28).\\
\textbf{Used in:} Paper~28 (QED on $S^3$), Paper~32 (spectral triple
synthesis), sprint LS-1 (Uehling shift $-27.13$ MHz from
$\Pi = 1/(48\pi^2)$).

\subsection{Camporesi--Higuchi spinor lift}
\label{sec:proj_spinor}
\textbf{Source:} scalar bare $S^3$ graph.\\
\textbf{Target:} Dirac-on-$S^3$ graph with eigenvalues
$|\lambda_n| = n + 3/2$, degeneracies $g_n = 2(n{+}1)(n{+}2)$.\\
\textbf{Variables introduced:} $\alpha$ enters through $(Z\alpha)^2$
relativistic content.\\
\textbf{Dimension introduced:} dimensionless (the spinor lift itself
is unit-preserving).\\
\textbf{Transcendental signature:} half-integer Hurwitz zetas
$\zeta(s, n + 3/2)$, producing odd-$\zeta$ content $\zeta(3),
\zeta(5), \ldots$ at integer $s$ (Paper~28 \S~IV); Catalan~$G =
\beta(2)$ and Dirichlet~$\beta(4)$ at the vertex level via
quarter-integer shifts $\zeta(s, 3/4), \zeta(s, 5/4)$.  Lives in
the spinor-intrinsic algebraic ring $R_\text{sp} =
\mathbb{Q}(\alpha^2)[\gamma]/(\gamma^2 + (Z\alpha)^2 - 1)$
(Paper~18 \S~IV).\\
\textbf{Used in:} Paper~28 \S~IV, Paper~33, Tier-2 spin-ful composed
qubit encoding (Paper~14 \S~V).

\subsection{Wigner $3j$ angular coupling}
\label{sec:proj_3j}
\textbf{Source:} pair of angular-momentum bare-graph nodes.\\
\textbf{Target:} coupled angular-momentum eigenstate.\\
\textbf{Variables introduced:} none.\\
\textbf{Dimension introduced:} dimensionless.\\
\textbf{Transcendental signature:} pure rational
($\mathbb{Q}[\sqrt{2k+1}]_{k}$).\\
\textbf{Used in:} all Gaunt integral evaluations, Paper~22 angular
sparsity, every Slater $F^k$ computation.

\subsection{Wigner $D$-matrix rotation between molecular centers}
\label{sec:proj_wignerD}
\textbf{Source:} composed orbital block at center $A$.\\
\textbf{Target:} same block in molecular coordinate frame at center
$B$.\\
\textbf{Variables introduced:} nuclear position vector
$\vec{R}_{AB}$.\\
\textbf{Dimension introduced:} length.\\
\textbf{Transcendental signature:} preserves rationality with
$\mathbb{Q}[\sqrt{2}, \sqrt{3}, \sqrt{6}]$ algebraic content from
the Wigner $d$-matrix at non-collinear angles.\\
\textbf{Used in:} multi-center molecules in Paper~17 and Paper~14
(LiF, CO, N$_2$, NaCl, CH$_2$O, C$_2$H$_2$, C$_2$H$_6$, F$_2$).

\subsection{Wilson plaquette: edge $\to$ 2-cell}
\label{sec:proj_wilson}
\textbf{Source:} Hopf graph 1-cochains.\\
\textbf{Target:} 2-cells (plaquettes) carrying $\mathrm{SU}(2)$ link
variables.\\
\textbf{Variables introduced:} $\beta = 1/g^2$ (gauge coupling).\\
\textbf{Dimension introduced:} dimensionless.\\
\textbf{Transcendental signature:} non-abelian $\mathrm{SU}(2)$ Haar
content; in the maximal-torus reduction yields the abelian
$U(1)$ content of Paper~25.\\
\textbf{Used in:} Paper~30 (SU(2) Wilson lattice gauge), Paper~25
(Hodge-1 Laplacian as gauge propagator), transverse photon
construction (Paper~28 \S~transverse\_photon).

\subsection{Vector-photon promotion}
\label{sec:proj_vector_photon}
\textbf{Source:} scalar 1-cochain photon on bare graph.\\
\textbf{Target:} vector-harmonic photon with explicit $(q, m_q)$
quantum numbers and Wigner~$3j$ vertex coupling.\\
\textbf{Variables introduced:} none (combinatorial labels).\\
\textbf{Dimension introduced:} dimensionless; specifically a
solid-angle measure on the Hopf $S^2$ base.\\
\textbf{Transcendental signature:} $1/(4\pi)$ per loop, identified
as the $S^2$ Weyl exchange constant of the Hopf base (Paper~33).\\
\textbf{Used in:} Paper~33 (recovers six angular-momentum selection
rules), sprint VP-1 (vertex correction calibration).

\subsection{Molecule-frame hyperspherical separation}
\label{sec:proj_molframe}
\textbf{Source:} composed atomic blocks.\\
\textbf{Target:} molecule-frame hyperspherical coordinates
$(\rho, \alpha, \theta_{12}, \ldots)$ with internuclear distance $R$
appearing in the charge function.\\
\textbf{Variables introduced:} $R$ (internuclear distance); for
polyatomics, additional bond angles.\\
\textbf{Dimension introduced:} length.\\
\textbf{Transcendental signature:} piecewise-smooth in $R$ at the
adiabatic potential level (Track~S, Paper~15); Gaunt integrals
preserve rationality at the angular level.\\
\textbf{Used in:} Paper~15 (Level~4), Paper~17 (composed geometry),
all H$_2$/HeH$^+$/H$_2$O calculations.

\subsection{Drake--Swainson asymptotic subtraction}
\label{sec:proj_drake_swainson}
\textbf{Source:} divergent finite-basis spectral sum
$J_v(nl) / I_v(nl)$ for $\ell > 0$ Bethe-logarithm states, where the
velocity-form closure $I_v(nl) = (Z^4/n^3)\,\delta_{\ell,0}$ vanishes
exactly for $\ell > 0$ on the bare Sturmian-projected graph.\\
\textbf{Target:} K-independent regularized value
$\ln k_0(n,\ell) = J_v(n,\ell) / D_{\text{drake}}(n,\ell)$, with the
structural denominator
\[
  D_{\text{drake}}(n,\ell) = \frac{2(2\ell + 1)\,Z^4}{n^3}.
\]
This recovers $D_{\text{drake}}(nS) = I_v(nS) = 2Z^4/n^3$ for $\ell = 0$
(matching the closure value), and provides a non-vanishing structural
normalization for $\ell > 0$.\\
\textbf{Variables introduced:} subtraction scale $K$ (transient; the
final result is $K$-independent over a plateau spanning $\sim 3.6$
orders of magnitude in $K$, as verified at sprint LS-4).\\
\textbf{Dimension introduced:} energy via $K$ (transient,
cancelled in $\beta_{\text{low}}(K) + \beta_{\text{high}}(K) = J_v$);
the resulting $\ln k_0$ is dimensionless.\\
\textbf{Transcendental signature:} flow tier (Paper~18 \S~IV).
$D_{\text{drake}}$ itself is rational ($Z^4 \cdot \mathbb{Q}$); the
transcendental content sits in $J_v$, where the spectral-sum
logarithms $\ln |2 \Delta E_m / \mathrm{Ry}|$ are inherited from the
Sturmian projection and carry calibration $\alpha$-dependent pieces
through the matching to the standard Bethe formula.  The structural
denominator
$D_{\text{drake}}(n,\ell) = (\text{spin})\cdot(\text{angular degeneracy})\cdot Z^4/n^3$
is a purely combinatorial quantity (spin~$\cdot$~$2\ell+1$~$\cdot$~hydrogenic
density), making the projection's structural content
combinatorial-rational while its operational content is calibration.\\
\textbf{Used in:} sprint LS-4 (Bethe logs for 2P at $+2.4\%$, 3P at
$+6.3\%$, 3D at $-0.24\%$ versus Drake-Swainson 1990 tabulated
values).  The combined hydrogen Lamb shift via this projection
chained with Sturmian and spectral-action gives $-0.39\%$ residual at
the sweet-spot $N=40$ (transient cancellation; converges to LS-1
baseline $-3.10\%$ at $N \to \infty$, reflecting the residual
$\alpha^5$ multi-loop ceiling).  The projection is structurally a
sibling of the Schwartz~(1961) integral form, simplified to act on
the discrete Sturmian pseudostate basis directly.

\subsection{Rest-mass projection}
\label{sec:proj_restmass}
\textbf{Source:} bare relativistic spectrum on $S^3 \times \mathbb{R}$
(or any natural-geometry spectrum), in the massless limit.\\
\textbf{Target:} massive spectrum
$\omega_n^2 = \omega_n^{(0)\,2} + m^2$, where $\omega_n^{(0)}$ is the
massless eigenvalue from the ambient natural geometry.\\
\textbf{Variables introduced:} $m$ (rest mass; equivalently rest-energy
$mc^2$ at $c=1$).\\
\textbf{Dimension introduced:} mass $=$ energy at $c=1$.\\
\textbf{Transcendental signature:} \textit{trivial}.  The bare-graph
algebraic-extension ring is preserved when $m^2 \in \mathbb{Q}$ (or
any ring extension of $\mathbb{Q}$ already present in the ambient
spectrum).  Verified explicitly in sprint KG-1 for the Klein-Gordon
spectrum on $S^3 \times \mathbb{R}$: $\omega_n = c_n \sqrt{d_n}$ with
$c_n \in \mathbb{Q}$ and $d_n$ a positive square-free integer for all
$n \in [1, 50]$ across the rational $m^2$ panel
$\{0, 1, 1/4, 2\}$ (Paper~35).  The mass projection is parameter-only;
it does not introduce $\pi$ or any other transcendental.\\
\textbf{Used in:} muonic / electronic / positronium reduced-mass
corrections; Paper~23 nuclear-electronic composed Hamiltonian
($I \cdot S$ hyperfine coupling); any future Klein-Gordon or
Dirac-on-$S^3$ relativistic mass calculation; Paper~35
Proposition~1.A.

\subsection{Observation / temporal-window projection}
\label{sec:proj_observation}
\textbf{Source:} bare relativistic spectrum on a spatial natural
geometry crossed with open temporal direction $\mathbb{R}$.\\
\textbf{Target:} same spectrum on the spatial geometry crossed with
periodic Euclidean time $S^1_\beta$ of period $\beta$ (equivalently:
finite-time observation window of duration $\beta$ in real time;
finite-temperature partition function at temperature $T = 1/\beta$).
For bosons, temporal Matsubara modes $\omega^t_k = 2\pi k / \beta$,
$k \in \mathbb{Z}$; for fermions with antiperiodic time,
$\omega^t_k = (2k+1)\pi / \beta$.\\
\textbf{Variables introduced:} $\beta$ (or $T = 1/\beta$).\\
\textbf{Dimension introduced:} time $=$ inverse energy.\\
\textbf{Transcendental signature:} $2\pi \cdot \mathbb{Q}$ per
Matsubara mode in the spectrum directly; in spectral integrals over
the compactified direction, $\pi^{2k} \cdot \mathbb{Q}$ from
$\zeta_R(2k)$ values, including the Stefan-Boltzmann
constant $\pi^2/90$ in the high-$T$ limit of $S^3 \times S^1_\beta$.
The first $\pi$-bearing eigenvalue of the compactified Klein-Gordon
spectrum is the $(n{=}0, k{=}1)$ Matsubara mode at
$\omega^2 = 4\pi^2/\beta^2$, with no earlier step injecting $\pi$
(verified in sprint KG-2; Paper~35).\\
\textbf{Used in:} all finite-temperature observables (Casimir at
$T \neq 0$, free-energy density, Stefan-Boltzmann constants); all
QED loop integrals that involve continuous time integration (Schwinger
proper-time, Mellin-Barnes, etc.); spectral-action evaluation
$\mathrm{Tr}\, f(D^2/\Lambda^2)$ when the trace runs over a
compact-time direction; Paper~35 Proposition~1.B.  This projection is
structurally what an observer does when integrating an instantaneous
spectrum over a finite measurement window.

\subsection{Two-body Dirac / Breit retardation}
\label{sec:proj_breit_retardation}
\textbf{Source:} framework-native single-particle Eides~\S~3.2
self-energy / Lamb-shift bracket on the Sturmian-reparameterized
Fock graph, valid in the $\lambda_l \ll \lambda_n$ hierarchy where the
$1/M_\text{nucleus}$ recoil expansion is regular.\\
\textbf{Target:} full two-body Breit / Bethe--Salpeter Hamiltonian for
the bound state at order $\alpha^4$, retaining all retardation and
crossed-photon terms that the $1/M$ expansion drops as subleading.\\
\textbf{Variables introduced:} the rest-mass ratio $m_l/m_n$ (and not
the absolute mass~$m$ already introduced by the rest-mass projection
of \S~\ref{sec:proj_restmass}).  This ratio does \emph{not} vanish in
the equal-mass leptonic limit $m_l = m_n$ where the $1/M$ expansion
is invalid.  Distinct from \S~\ref{sec:proj_restmass}: rest-mass
projection rescales single-particle observables under
$m_e \to m_\text{red}$; this projection adds the genuine two-body
Breit kinematic content that has no single-particle analog.\\
\textbf{Dimension introduced:} energy.\\
\textbf{Transcendental signature:} $\alpha^4 \cdot \mathbb{Q}$.  Pure
rational coefficient $\times \alpha^4$, no transcendental injected
beyond $\alpha$ already carried by the spectral-action projection.
Structurally the cleanest tier: same calibration class as the leading
magnetic dipole-dipole structure, ring-preserving over
$\mathbb{Q}(\alpha)$.  Distinct from the magnetization-density /
Zemach projection class (Paper~34 catalogue Layer-2 inputs), which
couples to nuclear spatial structure; this projection is a relativistic
two-body kinematic effect at the leading $\alpha^4$ order.\\
\textbf{Used in:} positronium 1S--2S two-photon transition (sprint
precision-catalogue 2026-05-08, equal-mass leptonic limit
$\lambda_a = \lambda_b = 0.5$); positronium 1$^3$S$_1$ HFS (annihilation
channel separated as a distinct Layer-2 mechanism); future positronium
fine-structure tests; future tests of the equal-mass $\mu^+ e^-$
limit if measured.\\
\textbf{Empirical anchor:} Ps 1S-2S framework-native at $+64.75$ ppm
vs Fee~1993 $1{,}233{,}607{,}216.4(3.2)$ MHz; Layer-2 $\alpha^4$
Breit input from Penin--Pivovarov~1998 (PRL~80, 2101) at the
$\sim 80$ GHz scale closes the residual to sub-MHz.\\
\textbf{Honest scope.}  This projection has one empirical anchor as
of 2026-05-08 (Ps 1S-2S).  Its definition is structural --- the
construction (single-particle Eides bracket $\to$ full two-body
Breit Hamiltonian) is established physics (Bethe--Salpeter~1957,
Karplus--Klein~1952, Penin--Pivovarov~1998, Czarnecki--Melnikov--Yelkhovsky~1999,
Karshenboim~2005~\S~4) --- but the framework's catalogue
verification rests on this single observable until additional
equal-mass or near-equal-mass two-body precision data are tested.\\
\textbf{Structural reading.}  This is the \emph{structurally cleanest
known instance} of the multi-focal-composition wall named in
CLAUDE.md~\S~2.  The wall's other five instances (Sprint H1
Yukawa, LS-8a $Z_2/\delta m$, HF-3 recoil, HF-4 Zemach, HF-5
multi-loop $a_e$) all surface in the multi-loop / renormalization
sector at $\alpha^5$ or higher, where two-body-projection failure
is entangled with renormalization failure.  The Breit retardation
content surfaces at $\alpha^4$ --- one full order \emph{before}
the LS-8a renormalization wall --- so the two-body-projection
failure is isolated and visible without renormalization confounds.
This is also what places the projection in the Layer-2 input class
(structurally a sibling of magnetization-density / Zemach in
\S~\ref{sec:proj_magnetization_density}): both take the framework's
single-particle output and add a structural correction the bare
spectral action does not autonomously generate.

\subsection{Nuclear charge-density correction (Foldy/Friar)}
\label{sec:proj_charge_density}
\textbf{Source:} bare Coulomb $V_\text{ne}(\vec{r}) = -Z e^2 / |\vec{r}|$
on the Fock-projected $S^3$ graph (point-source nucleus).\\
\textbf{Target:} convolved Coulomb $V_\text{ne}(\vec{r}) \to \int
d^3r' \, \rho_E^\text{nuc}(\vec{r}')\,(-Ze^2/|\vec{r}-\vec{r}'|)$,
with $\rho_E^\text{nuc}$ the nuclear electric charge distribution; at
leading order in the small-$r$ expansion the correction reduces to
the Foldy / Friar contact term
$\Delta V = +\tfrac{2\pi}{3} Z\alpha\,\langle r^2\rangle_E\,\delta^3(\vec{r})$.\\
\textbf{Variables introduced:} $\langle r^2\rangle_E$ (or
equivalently the charge radius $r_E$ for a given species).\\
\textbf{Dimension introduced:} length (one new $[L]$ scale per
nuclear species: proton $r_p$, deuteron $r_d$, etc.; structurally
distinct from atomic / internuclear length scales).\\
\textbf{Transcendental signature:} ring-preserving over
$\mathbb{Q}(\alpha)$.  The Foldy / Friar correction is rational in
the charge form-factor moments at every order of the small-$r$
expansion; the only transcendental content is the foundational
$\alpha$ already present in the spectral-action chain.\\
\textbf{Used in:} hydrogen Lamb shift $r_p$ contribution (currently
consumed externally via Eides Tab.~7.3 / Friar 1979); $\mu$H Lamb
shift Friar moment leading-order (Sprint MH Track A; $4.3\%$ at
leading); future framework-native operator-level Foldy module
(sibling of \texttt{geovac/magnetization\_density.py}, not yet
implemented).\\
\textbf{Empirical anchor:} proton charge radius
$r_p = 0.84075(64)$ fm (CODATA 2022); deuteron charge radius
$r_d = 2.1415(45)$ fm (Pohl et al.~2017).  Hydrogen 2$S$ Lamb shift
contribution from $r_p$ at the $\sim +1.18$ MHz level (Sprint LS-7
reframing; Eides Tab.~7.3).\\
\textbf{Honest scope.}  The framework currently consumes $r_p, r_d$
as external scalars (\texttt{R\_PROTON\_BOHR} and analogs) at the
point-source level.  Operator-level construction of $V_\text{ne}
\to V_\text{ne} \star \rho_E$ is the natural next step but has not
been implemented; the catalogue rows that depend on this projection
treat $r_p$ as a Layer-2 input.\\
\textbf{Structural reading.}  Categorically distinct from
\S~\ref{sec:proj_magnetization_density}: the Hamiltonian operator
modified is $V_\text{ne}$ (Coulomb), not $A_\text{hf}^\text{contact}$
(Fermi contact).  Empirical separation: hydrogen Lamb shift depends
on $r_p$ but only trivially on $r_Z$, while 21~cm hyperfine depends
on $r_Z$ but only trivially on $r_p$.

\subsection{Nuclear magnetization-density correction (Zemach)}
\label{sec:proj_magnetization_density}
\textbf{Source:} bare Fermi-contact hyperfine operator
$A_\text{hf}^\text{contact} = (8\pi/3)\,\mu_e \cdot \mu_N\,\delta^3(\vec{r})$
on the Fock-projected $S^3$ graph (point-source nuclear magnetic
moment).\\
\textbf{Target:} extended-magnetization Fermi-contact operator
$A_\text{hf}^\text{contact}(1 - 2 Z\alpha m_e r_Z + O(r_Z^2))$,
where $r_Z = \int d^3r\, d^3r'\,\rho_E^\text{nuc}(\vec{r}')\,
\rho_M^\text{nuc}(|\vec{r}-\vec{r}'|)\,|\vec{r}|$ is the Zemach
radius defined as the convolution moment of the nuclear charge and
magnetization profiles.\\
\textbf{Variables introduced:} $r_Z$ per nuclear species
(sub-features: nuclear magnetic radius $r_M$ and Friar moment
$\langle r^3\rangle_{(2)}$ enter as profile-shape parameters).\\
\textbf{Dimension introduced:} length (a new $[L]$ scale,
structurally distinct from \S~\ref{sec:proj_charge_density}).\\
\textbf{Transcendental signature:} ring-preserving over
$\mathbb{Q}(\alpha)$ at leading order; profile-dependence at
sub-leading order admits Gaussian, exponential, and dipole forms
with structurally identical leading $-2 Z\alpha m_e r_Z$ behaviour
(profile independence at leading order verified in Sprint HF
Track 4 / Sprint MH Track C).\\
\textbf{Used in:} hydrogen 21~cm hyperfine HF-4 row (Sprint HF;
$r_Z = 1.045$ fm yields $-39.5$ ppm framework-native shift); $\mu$H
1S Zemach mass-enhancement (Sprint MH Track B/C; factor $185.94
\times$ ep enhancement at leading order); deuteron 1S hyperfine
(Sprint precision-catalogue; $r_Z(D) = 2.593$ fm yields $-98.0$ ppm
framework-native shift).\\
\textbf{Empirical anchors:} proton Zemach radius
$r_Z = 1.045(20)$ fm (Eides 2024 atomic compilation) or $r_Z =
1.013(10)(12)$ fm (lattice QCD 2024); deuteron Zemach radius
$r_Z(D) = 2.593(16)$ fm (Friar--Payne 2005).\\
\textbf{Honest scope.}  Operator-level construction is mature:
\texttt{geovac/magnetization\_density.py} ($\sim 600$ lines, 57
tests; Sprint MH Track C + W1b recoil-mixing extension, May 2026)
implements the convolved Fermi-contact operator with the
\texttt{lepton\_mass} parameter eliminating manual mass-scaling for
muonic systems and an opt-in \texttt{include\_recoil\_mixing} flag
extending the kernel to next-to-leading order with the recoil-mixing
prefactor $m_l/(m_l + m_n)$ (per arXiv:2604.06930 eq.~95) and the
Friar moment correction $(1/2)(Z\alpha m_\text{red})^2 \langle r^2
\rangle_{(2)}$ (per Friar~1979 / Eides~\S6.2).  At leading order the
kernel matches Eides analytical to floating-point precision; with
recoil-mixing on, the framework muH 1S HFS prediction matches
Krauth~2017's full-theory Zemach line to $\sim 1\%$ (vs $\sim 3\%$
overshoot at LO-only).  $r_Z$ values themselves are Layer-2 inputs
(cannot be predicted from framework structure; consumed from QCD
lattice or atomic spectroscopy).  The W1b recoil-mixing extension
sprint (May 2026, see \texttt{debug/w1b\_recoil\_mixing\_extension\_memo.md})
showed that this kernel improvement does NOT close the $\sim 0.22$~fm
muH-alone vs Eides offset diagnosed by Sprint Calc-rZG-extended;
that offset is now attributable to Layer-2 itemization-convention
mismatch between Krauth~2017, Eides~2024, and Antognini-Krauth-Pohl~2015,
not to kernel approximation.\\
\textbf{Structural reading.}  Sibling of
\S~\ref{sec:proj_charge_density} (different Hamiltonian operator)
and of \S~\ref{sec:proj_breit_retardation} (Breit retardation:
both take the framework's single-particle output and add a
structural correction the bare spectral action does not
autonomously generate).  Categorical distinction from
charge-density is empirical: 21~cm hyperfine depends on $r_Z$ but
trivially on $r_p$, electronic Lamb shift depends on $r_p$ but
trivially on $r_Z$.

\subsection{Nuclear tensor multipole coupling}
\label{sec:proj_tensor_multipole}
\textbf{Source:} multipole expansion of the nuclear charge potential
$V_\text{ne}(\vec{r}) = -Ze^2/|\vec{r}|$, truncated at the
monopole on the bare Fock-projected $S^3$ graph (spherical-symmetric
nucleus).\\
\textbf{Target:} same potential with $L \geq 2$ multipole
corrections retained.  At leading rank-2: quadrupole coupling
$H_Q = -e\,Q_N\,T^{(2)}_{ij}\,(\partial_i E_j)/6$, with $Q_N$ the
nuclear electric quadrupole moment and $T^{(2)}_{ij}$ a rank-2
spin tensor on the nuclear-spin Hilbert space.  Higher ranks
(octupole, hexadecapole) follow the same pattern at higher angular
momentum.\\
\textbf{Variables introduced:} $Q_N$ per nuclear species (rank-2
leading); higher-rank moments enter at higher $L$.\\
\textbf{Dimension introduced:} length-squared (the rank-2
quadrupole has dimension $[L]^2$; rank-$L$ multipoles have dimension
$[L]^L$).  This is the only projection in \S~\ref{sec:layer2} whose
dimensional injection is not strictly $[L]^1$ or $[E]$; the
squared-length structure reflects the tensor character of the
operator.\\
\textbf{Transcendental signature:} ring-preserving over $\mathbb{Q}$
at the angular level (multipole expansion uses Wigner $3j$ at every
step).  Sub-leading corrections involve the same algebraic ring as
\S~\ref{sec:proj_charge_density}.\\
\textbf{Used in:} deuteron 1S hyperfine s-state contribution
(sub-ppm; \emph{not} load-bearing in current catalogue); future
molecular HD / D$_2$ rotational hyperfine where electronic
gradients at the nucleus are nonzero and $Q_d$ couples at the
percent level (Komasa--Pachucki 2020 sub-percent precision regime);
future heavy-nucleus hyperfine spectra where $L \geq 1$ valence
orbitals couple to $Q_N$ at leading order.\\
\textbf{Empirical anchors:} deuteron quadrupole moment
$Q_d = 0.285699(15)(18)$ fm$^2$ (Komasa--Pachucki 2020); $^{17}$O
and $^{14}$N quadrupole moments for molecular spectroscopy
benchmarks; sub-percent precision spectroscopy in HD / D$_2$
rotational HFS as the load-bearing test.\\
\textbf{Honest scope.}  Forward-looking projection slot.  No
catalogue row currently uses this projection at leading order
(deuteron 1S HFS s-state sees only sub-ppm contribution; the
load-bearing tests are in molecular rather than atomic
spectroscopy).  Operator-level construction has not been
implemented; this projection occupies a structural slot, with the
empirical anchor in molecular HD / D$_2$ HFS not currently
catalogued.\\
\textbf{Structural reading.}  Categorically distinct from
\S~\ref{sec:proj_charge_density} and
\S~\ref{sec:proj_magnetization_density}: tensor coupling to a
gradient is structurally different from scalar coupling to a contact
density.  Different selection rules: $Q_N$ requires $l\geq 1$ on the
electronic side (does not contribute to s-states at leading order),
while charge-density and magnetization-density Foldy / Friar / Zemach
contributions are pure $l=0$ contact corrections.  The
Wigner--Eckart theorem ensures the tensor and scalar projections
are mutually orthogonal: they cannot be unified by a single
profile-broadening prescription.

\subsection{Boundary of \S\ref{sec:layer2}: internal graph-derived variables}
\label{sec:layer2_boundary}

The nineteen projections enumerated above share a structural
property worth naming explicitly: each introduces a variable that
is \emph{internal to the framework's graph and projection structure}
--- a graph-derived integer or rational invariant (such as $Z$, $n$,
the foundational $\alpha$, or a Wigner $3j$ coefficient), a
continuous parameter that controls a coordinate or focal-length
choice (such as $\rho$, $R$, $\Lambda$, $\beta$, $m$, $r_E$, $r_Z$,
$Q_N$), or a function of prior projection variables (compositional
projections per Observation~\ref{obs:composition}).

This is structurally distinct from \emph{Layer-2 calibration inputs}
that the framework consumes from external sources but does not
autonomously generate.  Examples: nuclear g-factors $g_p, g_d$
(extracted from precision spectroscopy or QCD lattice), multi-loop
QED coefficients (Eides Tab.~7.3, Karshenboim 2005, Penin--Pivovarov
1998 $\alpha^4$ Breit input), Yukawa couplings (per CLAUDE.md
\S~1.7 WH5; Sprint H1 verdict), and external classical-field
strengths (Zeeman $B$, Stark $\vec{\mathcal{E}}$ --- verified by the
2026-05-09 external-field hunting sprint to be Layer-2 inputs, not
\S~\ref{sec:layer2} projections; the bare Landau spectrum is HO-class
and $\pi$-free, structurally a Bargmann--Segal sibling on the
cyclotron plane, with $\pi$ entering only through the 2D phase-space
measure as a calibration~$\pi$ in the M1 Mellin sub-mechanism, not
through Matsubara compactification).

The structural test that distinguishes \S~\ref{sec:layer2}
projection from Layer-2 input is whether the variable's content is
internal graph-derived or external classical-field /
QCD-lattice-derived.  Zemach and charge radii are intrinsic
focal-lengths once the framework's spectral content (not external
calibration) determines the convolved-potential structure given the
radius; nuclear g-factors and Yukawas are external because no
spectral structure of the framework determines their value.  This
boundary is not absolute: promotion of a Layer-2 input to a
projection slot is the structural commitment that the variable
acquires a focal-length role beyond bare scalar consumption, as
Zemach did when \texttt{magnetization\_density.py} was implemented
at the operator level (\S~\ref{sec:proj_magnetization_density}).
Future Layer-2 inputs may follow if the framework develops
operator-level constructions that give them focal-length status.

% ===================================================================
\section{The Three-Axis Dictionary}
\label{sec:dictionary}
% ===================================================================

Table~\ref{tab:projection_axes} consolidates Section~\ref{sec:layer2}
into the three-axis dictionary that this paper proposes as a
companion to Paper~18's transcendental classification.

\begin{table}[h]
\centering\small
\begin{tabular}{p{3.5cm} p{3cm} p{2cm} p{4cm}}
\toprule
Projection & Variables added & Dimension & Transcendental signature \\
\midrule
Fock conformal & $Z$, $E$ & energy & $\kappa = -1/16$;
$\pi$ via $\mathrm{Vol}(S^3)$ \\
Hopf bundle & $\alpha$ & dimensionless & $\pi = \mathrm{Vol}(S^2)/4$ \\
Bargmann--Segal & $\hbar\omega$ & energy & \textit{none} (graph
$\pi$-free) \\
Stereographic & $r$ (length) & length & conformal factor (Coulomb
$1/r$) \\
Sturmian reparam & $Z, n$ (re-use) & preserves & rational
preserved \\
Spectral action & $\Lambda$, $\alpha$ & energy via $\Lambda$ &
$\sqrt{\pi}\times \mathbb{Q}$ SD coefficients;
$\pi^{2k}\times \mathbb{Q}$ observables \\
Camporesi--Higuchi spinor & $\alpha$ via $(Z\alpha)^2$ &
dimensionless & half-integer Hurwitz $\to$ odd~$\zeta$;
Catalan~$G$, $\beta(4)$ at vertex \\
Wigner $3j$ & none & dimensionless &
$\mathbb{Q}[\sqrt{2k{+}1}]_k$ \\
Wigner $D$ rotation & $\vec{R}_{AB}$ & length &
$\mathbb{Q}[\sqrt{2}, \sqrt{3}, \sqrt{6}]$ \\
Wilson plaquette & $\beta = 1/g^2$ & dimensionless & SU(2) Haar \\
Vector-photon promotion & none & dimensionless & $1/(4\pi)$ per loop \\
Molecule-frame hypersph. & $R$ & length & piecewise-smooth in $R$ \\
Drake--Swainson asymptotic subtraction & $K$ (cancels) & dimensionless &
flow tier; $D_{\text{drake}} = 2(2\ell+1)Z^4/n^3$ rational \\
Rest-mass projection & $m$ & mass &
\textit{trivial} (ring-preserving for $m^2 \in \mathbb{Q}$) \\
Observation / temporal-window & $\beta$ (or $T = 1/\beta$) & time &
$2\pi \cdot \mathbb{Q}$ per Matsubara mode; $\pi^{2k}\cdot\mathbb{Q}$ in
integrated quantities; Stefan-Boltzmann $\pi^2/90$ in high-$T$ limit \\
Two-body Dirac / Breit retardation & $m_l/m_n$ & energy &
$\alpha^4 \cdot \mathbb{Q}$ (ring-preserving over $\mathbb{Q}(\alpha)$) \\
Charge-density (Foldy/Friar) & $r_E$ per nucleus & length &
$\mathbb{Q}(\alpha)$ ring-preserving \\
Magnetization-density (Zemach) & $r_Z$ per nucleus & length &
$\mathbb{Q}(\alpha)$ ring-preserving (leading); profile-shape
sub-features at sub-leading \\
Tensor multipole (quadrupole, etc.) & $Q_N$ per nucleus &
length-squared & $\mathbb{Q}$ via Wigner $3j$ at angular level \\
\bottomrule
\end{tabular}
\caption{The projection dictionary.  Each row is a structural
construction that introduces specific physical content into Layer~1
graph data.  The fourth column is the transcendental signature in
the sense of Paper~18; the second and third columns are the new
axes proposed by this paper.  The thirteenth row (Drake--Swainson
asymptotic subtraction) was added in sprint LS-4 to close the 2P
Bethe-log structural floor identified in LS-3; the fourteenth and
fifteenth rows (rest-mass and observation/temporal-window) were
added in sprint KG / Paper~35 (May~2026) to resolve the
mass-vs-observation ambiguity that had been bundled with other
parameter projections; the sixteenth row (two-body Dirac / Breit
retardation) was added in sprint precision-catalogue 2026-05-08
motivated by the Ps 1S-2S sixth-multi-focal-wall instance, where
the $1/M_\text{nucleus}$ expansion fails at equal mass and the
$\alpha^4$ Breit content surfaces one order before the LS-8a
renormalization wall.}
\label{tab:projection_axes}
\end{table}

\begin{observation}[Two-axis duality]
\label{obs:duality}
The transcendental classification of Paper~18 and the projection
classification of Table~\ref{tab:projection_axes} are not independent.
Each Paper~18 transcendental tier corresponds to specific projections,
and each projection's transcendental signature is a specific Paper~18
tier:
\begin{itemize}\setlength\itemsep{0pt}
\item Calibration $\pi^{2k}$ tier $\leftrightarrow$ spectral action /
Hopf bundle / vector-photon promotion projections.
\item Spinor-intrinsic $(Z\alpha)$ tier $\leftrightarrow$
Camporesi--Higuchi spinor lift.
\item Odd-$\zeta$ tier $\leftrightarrow$ Camporesi--Higuchi spinor
lift in combination with Mellin/heat-kernel evaluation at integer~$s$.
\item Dirichlet~$L$ tier $\leftrightarrow$ vertex topology composed
with Camporesi--Higuchi at quarter-integer Hurwitz shift.
\item Embedding tier $\leftrightarrow$ molecule-frame hyperspherical
or non-trivial bound-state coupling.
\item Combinatorial-rational tier $\leftrightarrow$
Wigner $3j$ / $D$ alone (no projection beyond Layer~1); also the
structural denominator $D_{\text{drake}}(n,\ell) = 2(2\ell+1)Z^4/n^3$
in the Drake--Swainson projection.
\item Flow tier (Paper~18 \S~IV) $\leftrightarrow$ Drake--Swainson
asymptotic subtraction.  This is the projection that takes a
divergent finite-basis ratio (where closure makes the denominator
vanish) and produces a $K$-independent regularized value through
asymptotic subtraction.  The flow content sits in the
$\alpha$-dependent logarithms of the regulated spectral sum, which
cancel against the $K$-cutoff terms.
\item \emph{Trivial / ring-preserving tier} $\leftrightarrow$
rest-mass projection.  The mass projection introduces a parameter
without introducing a transcendental: the bare-graph algebraic
ring is closed under the addition of $m^2$ when $m^2 \in \mathbb{Q}$.
This is the only projection in the dictionary that introduces a
physical variable while preserving the transcendental signature of
the source spectrum.  Paper~35 verifies this explicitly for the
Klein-Gordon spectrum on $S^3 \times \mathbb{R}$.
\item \emph{Calibration $2\pi \cdot \mathbb{Q}$ tier (temporal subtype)}
$\leftrightarrow$ observation / temporal-window projection.  This is
the categorical home of $\pi$ in the framework: every $\pi$ that
appears in a GeoVac observable can be traced to either (a)~a Hopf-base
calibration~$\pi$ (Hopf bundle / vector-photon promotion), (b)~a
spectral-action $\sqrt{\pi}\cdot\mathbb{Q}$ (heat-kernel coefficient),
or (c)~a temporal Matsubara $2\pi/\beta$ (this projection).  The
Casimir constant on $S^3 \times S^1_\beta$ is the cleanest
illustration: spatial-only Casimir is $1/240$ (rational, no $\pi$);
adding $S^1_\beta$ injects $\pi$ exactly through the Matsubara
mechanism.
\item \emph{Layer-2 input tier (two-body kinematic subtype)}
$\leftrightarrow$ two-body Dirac / Breit retardation projection
(\S~\ref{sec:proj_breit_retardation}).  Same ring-preserving
character as the rest-mass projection (no transcendental injected
beyond the $\alpha$ already present in the spectral-action chain),
but distinct mechanism: rest-mass projection rescales single-particle
observables by $m_e \to m_\text{red}$, whereas Breit retardation
introduces the rest-mass ratio $m_l/m_n$ as a two-body kinematic
control parameter.  In the equal-mass leptonic limit
$m_l/m_n \to 1$ the standard $1/M_\text{nucleus}$ recoil expansion
is invalid, and the projection becomes load-bearing for level
energies at $\alpha^4$.  Empirical anchor: Ps 1S--2S framework-native
at $+64.75$ ppm in \S~\ref{sec:offprecision}; closure to sub-MHz
requires the Penin--Pivovarov 1998 $\alpha^4$ Breit input as
Layer-2~content.
\end{itemize}
\end{observation}

\begin{remark}[variable--dimension correlation]
\label{rem:vd_correlation}
The variable axis and dimension axis are not independent: every
projection in Section~\ref{sec:layer2} that introduces a physical
dimension also introduces a physical variable that carries that
dimension.  Of the nineteen projections, twelve add a dimension and
all twelve also add the variable that carries it (perfect correlation
on the V/D side); zero projections add a dimension without a
variable.  The three intra-nuclear projections of
\S~\ref{sec:proj_charge_density}, \S~\ref{sec:proj_magnetization_density},
and \S~\ref{sec:proj_tensor_multipole} preserve this V/D
correlation: each adds an intra-nuclear focal length together with
its associated dimension ($[L]$, $[L]$, $[L]^2$ respectively).  The transcendental axis is the only one that admits an
independence witness --- vector-photon promotion (\S~\ref{sec:layer2},
projection~11) introduces $1/(4\pi)$ without introducing any new
variable or dimension.  This is the structural content of
Section~\ref{sec:axiomatization} (Theorem~3 composition table):
of the three axes, only the transcendental axis is genuinely
independent of the other two.  The sixteenth projection
(two-body Dirac / Breit retardation,
\S~\ref{sec:proj_breit_retardation}) preserves the V/D correlation:
it adds the variable $m_l/m_n$ and the dimension energy together,
and lives in the ring-preserving $\alpha^4 \cdot \mathbb{Q}$ tier.
\end{remark}

% ===================================================================
\subsection{Base-unit decomposition of the projection family}
\label{sec:base_units}
% ===================================================================

Remark~\ref{rem:vd_correlation} establishes that every dimensioned
projection is paired with the variable carrying that dimension, at
the level of compound dimensions (energy, length, time, mass).  This
subsection decomposes the compound dimensions into base-unit
constituents and asks: how often is each base unit injected across
the nineteen projections, and which projections inject which units?

We work in atomic units throughout: $\hbar = m_e = e = 1$, so the
mechanical base units are $[L]$, $[M]$, $[T]$, with charge $[Q]$
carried by the integer label $Z$.  Energy is the compound
$[E] = [M][L]^2[T]^{-2}$, but in atomic units acts as an effective
base for the framework's primary observables; we track $[E]$
alongside the strictly-base units below.  We follow the
\S~\ref{sec:layer2} order of projections for cross-reference with
Table~\ref{tab:projection_axes}.

\begin{table}[h]
\centering\small
\begin{tabular}{p{3cm} p{2.2cm} p{0.9cm} p{4.5cm} p{1.6cm}}
\toprule
Projection & Variable & Dim & Base-unit injection & Role \\
\midrule
Fock conformal & $Z$, $E$ & $[E]$ &
foundational anchor for atomic units; sets $a_0$ ($[L]$) and Ry
($[E]$) jointly via $p_0 = Z/n$; $[Q]$ as integer label $Z$ &
foundational \\
Hopf bundle & $\alpha$ (Paper~2) & dim'less & none & pure \\
Bargmann--Segal & $\hbar\omega$ & $[E]$ &
foundational anchor (HO sibling of Fock); sets oscillator length
($[L]$) and $\hbar\omega$ ($[E]$) jointly &
foundational (sibling) \\
Stereographic & $r$ & $[L]$ & $[L]$ as continuous coordinate; no new
scale (coordinatizes the foundational $[L]$) &
coordinate-explicit \\
Sturmian reparam & $Z, n$ (re-use) & --- & none & reparam \\
Spectral action & $\Lambda$, $\alpha$ & $[E]$ & new $[E]$ scale (UV
cutoff $\Lambda$); $\alpha$ dimensionless & scale-extension \\
Camporesi--Higuchi spinor & $\alpha$ via $(Z\alpha)^2$ & dim'less &
none & pure \\
Wigner $3j$ & none & dim'less & none & pure \\
Wigner $D$ rotation & $\vec{R}_{AB}$ & $[L]$ & new $[L]$ scale
(internuclear position vector) & scale-extension \\
Wilson plaquette & $\beta = 1/g^2$ & dim'less & none & pure \\
Vector-photon promotion & none & dim'less & none & pure \\
Mol-frame hyperspherical & $R$ & $[L]$ & new $[L]$ scale
(internuclear distance; same physical length as $|\vec{R}_{AB}|$) &
scale-extension \\
Drake--Swainson & $K$ (cancels) & $[E]$ & transient $[E]$ scale
(cancels in $\beta_\text{low}(K) + \beta_\text{high}(K)$) &
scale-extension (transient) \\
Rest-mass & $m$ & $[M]$ & first $[M]$ scale beyond foundational
$m_e = 1$; one projection per particle species &
scale-extension (new base) \\
Observation / temporal-window & $\beta$ (or $T$) & $[T]$ & first
$[T]$ scale beyond foundational $\tau_0 = \hbar/\mathrm{Ry}$; one
projection per measurement window &
scale-extension (new base) \\
Two-body Dirac / Breit & $m_l / m_n$ & dim'less &
ratio of two prior $[M]$ scales (lepton, nucleon); no new base
content & composition \\
Charge-density (Foldy/Friar) & $r_E$ per nucleus & $[L]$ &
new $[L]$ scale (intra-nuclear, charge distribution); modifies
$V_\text{ne}$ via convolution & scale-extension \\
Magnetization-density (Zemach) & $r_Z$ per nucleus & $[L]$ &
new $[L]$ scale (intra-nuclear, magnetization distribution);
modifies $A_\text{hf}^\text{contact}$ & scale-extension \\
Tensor multipole (rank-$\geq 2$) & $Q_N$ per nucleus & $[L]^2$ &
rank-$L$ multipole moment, dim $[L]^L$; couples to electronic
gradients ($l \geq 1$ on electronic side) &
scale-extension (rank-2) \\
\bottomrule
\end{tabular}
\caption{Base-unit decomposition of the nineteen projections.
``Variable'' is the new physical content; ``Dim'' repeats the
compound dimension from Table~\ref{tab:projection_axes};
``Base-unit injection'' shows what enters in $[L]$, $[M]$, $[T]$,
$[Q]$, with $[E]$ tracked as the framework's effective base
observable; ``Role'' classifies the projection by its action on the
dimensional content of a projection chain.  Six roles emerge:
\emph{foundational} (sets the atomic-units anchor),
\emph{coordinate-explicit} (re-expresses an anchored unit as a
continuous coordinate), \emph{reparam} (relabels without adding
content), \emph{scale-extension} (adds a new variable in a base
unit, possibly a new base), \emph{pure} (adds transcendental and/or
combinatorial content with no dimensional injection), and
\emph{composition} (variable is a function of prior projection
variables).}
\label{tab:base_units}
\end{table}

The inverse view --- which projections inject each base unit ---
makes the structure visible:

\begin{itemize}\setlength\itemsep{0pt}
\item $[L]$ (scalar): \emph{six} projections.  Fock or
Bargmann--Segal anchors the foundational length ($a_0$ or oscillator
length); stereographic re-expresses it as a continuous coordinate
$r$; Wigner~$D$ rotation introduces nuclear position
$\vec{R}_{AB}$; molecule-frame hyperspherical separation introduces
internuclear distance $R$ (Wigner~$D$ and mol-frame inject the same
physical length scale through different projections);
charge-density (\S~\ref{sec:proj_charge_density}) introduces the
intra-nuclear charge radius $r_E$; magnetization-density
(\S~\ref{sec:proj_magnetization_density}) introduces the
intra-nuclear Zemach radius $r_Z$.  These last three intra-nuclear
length scales are each \emph{categorically distinct} from the
atomic / internuclear scales: they modify different Hamiltonian
operators (Coulomb $V_\text{ne}$, Fermi-contact
$A_\text{hf}^\text{contact}$) with different selection rules, with
empirical separation across hydrogen Lamb shift / 21~cm hyperfine /
deuteron molecular HFS.

\item $[L]^2$ (rank-2 tensor): \emph{one} projection.  Tensor
multipole (\S~\ref{sec:proj_tensor_multipole}) introduces the
nuclear quadrupole moment $Q_N$ at rank-2 leading; higher-rank
multipoles ($L \geq 3$) follow at higher $[L]^L$.  Categorically
distinct from scalar $[L]$ injections: the operator is tensorial
and requires $l \geq 1$ on the electronic side.
\item $[E]$: \emph{four} projections.  Fock or Bargmann--Segal
anchors the foundational energy (Ry or $\hbar\omega$); spectral
action introduces UV cutoff $\Lambda$; Drake--Swainson introduces a
transient subtraction scale $K$ (which cancels in the regularized
result); two-body Breit retardation introduces $\alpha^4 \cdot
\mathrm{Ry}$ as a correction-energy scale, where $\alpha^4$ is
dimensionless and Ry is foundationally anchored.
\item $[M]$: \emph{one} projection (rest-mass).  Each particle
species in the system requires a separate rest-mass projection;
two-particle systems compose two such projections.
\item $[T]$: \emph{one} projection (observation /
temporal-window).  Each measurement window or temperature requires
a separate observation projection.
\item $[Q]$: integer-counter label $Z$ at the Fock projection only.
$[Q]$ does not appear as a continuous scale anywhere in the
projection family --- charge is always counted, never
measured-against-a-scale.
\end{itemize}

\begin{observation}[$L$--$E$ asymmetry]
\label{obs:LE_asymmetry}
Of the mechanical base units (and the charge unit), only $[L]$ and
$[E]$ are injected by multiple projections; $[M]$, $[T]$, and $[Q]$
are each instantiated by exactly one projection beyond the
foundational atomic-units anchor.  After the addition of the three
intra-nuclear projections (charge-density, magnetization-density,
tensor multipole), the asymmetry is no longer symmetric: $[L]$
admits \emph{six} scalar instances plus one rank-2 tensor instance,
while $[E]$ admits four.  The framework records more length scales
(atomic, internuclear, polyatomic, and three intra-nuclear
distribution scales) than energy scales (Rydberg, UV cutoff,
multi-loop subtraction, two-body Breit correction).  The structural
reason: every distinct spatial structure introduces an independent
focal length because different intra-nuclear and intra-molecular
geometries enter different Hamiltonian operators with different
selection rules; spectral content has fewer distinct roles
(foundational, regulator, subtraction, correction).  The conjugate
pair $[L] \leftrightarrow [E]$ remains the only multi-instance
axis, but $[L]$ now dominates: physics has more places to put
length scales than to put energy scales.
\end{observation}

\begin{observation}[Compositional projections]
\label{obs:composition}
The two-body Dirac / Breit retardation projection
(\S~\ref{sec:proj_breit_retardation}) is the first identified
instance of a structural family of \emph{compositional projections}:
projections whose variable is a function (typically a ratio or
product) of variables introduced by prior projections.  All other
projections in Table~\ref{tab:base_units} introduce variables
independent of prior content; compositional projections are
structurally distinct because they are \emph{gated} by the prior
instantiation of two or more distinct values of a foundational
variable in the same projection chain.

We catalogue two confirmed instances of compositional projections:

\begin{enumerate}\setlength\itemsep{0pt}
\item \emph{Kinematic composition.}  Variable: rest-mass ratio
$m_l / m_n$.  Gating: prior instantiation of two distinct $[M]$
scales by separate rest-mass projections (\S~\ref{sec:proj_restmass}).
Manifests at $\alpha^4$ in two-body bound-state energies, one full
order before the $\alpha^5$ multi-loop renormalization wall.
Empirical anchor: positronium 1S--2S framework-native at
$+64.75$~ppm (\S~\ref{sec:catalog}; sixth and structurally cleanest
instance of the multi-focal-composition wall named in
CLAUDE.md~\S~2).

\item \emph{Angular composition.}  Variable: ratio of effective
charges $Z_a / Z_b$ within a single atom (e.g., He 1s
$Z_\text{eff} = 2$ vs.\ 2p $Z_\text{eff} = 1$).  Gating: prior
instantiation of two distinct foundational Fock anchors at different
effective charges within the same electronic Hilbert space (the
``internal multi-focal'' pattern of \S~\ref{sec:catalog}).
Manifests at $\alpha^2$ in fine-structure splittings, where the
bipolar harmonic decomposition $(k_1{=}0, k_2{=}2)$ direct $+$
$(k_1{=}1, k_2{=}1)$ exchange and the Drake combining coefficients
$(\tfrac{3}{50}, -\tfrac{2}{5}, \tfrac{3}{2}, -1)$ are the
structural carrier.  Empirical anchor: He $2{}^3P_0$--$2{}^3P_1$
at $-0.014\%$ and $2{}^3P_0$--$2{}^3P_2$ span at $-0.20\%$
(\S~\ref{sec:catalog}).
\end{enumerate}

The two instances differ in their gating mechanism (kinematic at
the two-body level vs.\ angular at the two-orbital level) but share
the compositional character: each requires the prior projection
chain to have instantiated two distinct values of a foundational
variable, so that the compositional projection's variable (ratio,
product, or function) is well-defined.

The structural prediction the generalization makes: any future
GeoVac observable whose computation requires a function of two or
more foundationally-anchored variables in the same projection chain
belongs to this family, and the order at which it manifests in
$\alpha$-perturbation theory follows the projection at which the
function is gated.  Candidate compositional projections not yet
observed:

\begin{itemize}\setlength\itemsep{0pt}
\item Ratio of two length scales within a single nucleus (e.g.,
$r_\text{charge} / r_Z$, the Friar/Foldy-versus-Zemach distinction
inside one proton; would gate sub-Zemach corrections to hyperfine).
\item Product of three rest-mass scales (three-body kinematic
content beyond Breit retardation; would gate three-body bound-state
corrections at $\alpha^4$ or higher).
\item Ratio of two temperature scales $\beta_a / \beta_b$ in
multi-window observation (e.g., simultaneous interrogation at two
distinct $T$).
\end{itemize}

The compositional-projection family is open-ended: confirmed
instances locate the gating-order pattern (Breit at $\alpha^4$,
internal multi-focal at $\alpha^2$), and additional empirically
anchored instances would refine the structural prediction.
\end{observation}

This decomposition sharpens
Remark~\ref{rem:vd_correlation}: V--D correlation holds at the
compound dimension level, but at the base-unit level $[L]$ and
$[E]$ each admit four instances while $[M]$, $[T]$ admit one each
and $[Q]$ enters only as an integer label.  The base-unit
decomposition is a finer falsifiable test than the compound
dimension alone: an observable that depends on multiple length
scales must engage the foundational projection plus the relevant
length-injecting projections; an observable that depends on
temperature must engage the observation projection; an observable
depending on a species mass ratio must engage two rest-mass
projections plus Breit retardation.  The empirical matches
catalogue (\S~\ref{sec:catalog}) provides a forward test:
each row's projection chain should be traceable to a base-unit
signature consistent with the dimensions of the matched
observable.

% ===================================================================
\section{The Empirical Matches Catalogue}
\label{sec:catalog}
% ===================================================================

We list GeoVac quantities that match standard known physics values
to machine precision or stated uncertainty.  Each row is tagged by
projection(s), variables involved, dimension, transcendental class,
and match precision.  This is a structural index, not an exhaustive
benchmark report; the underlying numerical evidence lives in
\texttt{tests/} and the cited papers.

\begin{longtable}{p{3.5cm} p{2.8cm} p{1.3cm} p{2.2cm} p{1.8cm} p{1.4cm}}
\toprule
Match & Projection(s) & Vars & Dim & Trans.\ class & Match \\
\midrule
\endhead
\textbf{Zero-projection (Layer~1 only)} & & & & & \\
Pauli count $N = 11.11{\cdot}Q$ & --- & none & dimensionless &
integer/rational & exact \\
Magic numbers $2,8,20,28,\ldots$ & HO + LS counting & none &
dimensionless & integer & exact \\
Gaunt selection rules & 3$j$ algebra & none & dimensionless &
rational & exact \\
$\pi$-free certificate (Coulomb) & --- & none & dimensionless &
rational & exact \\
$\pi$-free certificate (Bargmann) & --- & none & dimensionless &
rational & exact \\
Angular sparsity 1.44\% at $\ell_\text{max}=3$ & 3$j$ algebra &
none & dimensionless & rational & exact \\
Discrete $c_1 = 0$ on Hopf-$S^3$ Fock graph at $n_\text{max} = 2, 3$
(Condon--Shortley $+$ $Z_2$ alternating) & --- &
none & dimensionless & integer & exact (Fraction) \\
\midrule
\textbf{Single-projection} & & & & & \\
Hydrogen $E_n = -Z^2/(2n^2)$ Ha & Fock + $\kappa$ & $Z, E$ &
energy & rational & exact symbolic \\
2$p$ doublet $\alpha^2/32$ Ha (Z=1) & spinor lift $\circ$ Fock &
$Z, \alpha$ & energy & $(Z\alpha)^2{\cdot}\mathbb{Q}$ & exact symbolic \\
Dirac fine-structure formula ($n\le 4$) & spinor lift $\circ$ Fock &
$Z, \alpha$ & energy & $(Z\alpha)^2{\cdot}\mathbb{Q}$ & exact \\
$\Pi = 1/(48\pi^2)$ VP coefficient & spectral action &
$\Lambda, \alpha$ & dimensionless & $\pi^2{\cdot}\mathbb{Q}$ &
exact symbolic \\
$\beta(\alpha) = 2\alpha^2/(3\pi)$ one-loop & spectral action &
$\alpha$ & dimensionless & $\pi{\cdot}\mathbb{Q}$ & exact symbolic \\
$D^2$ heat kernel SD coefficients & spectral action & $\Lambda$ &
energy & $\sqrt{\pi}{\cdot}\mathbb{Q}$ & exact symbolic \\
$F = \zeta_R(2) = D_{n^2}(d_\text{max})$ & Fock + Dirichlet at
$s=4$ & graph data & dimensionless & $\pi^2{\cdot}\mathbb{Q}$ &
exact symbolic \\
$\Delta^{-1} = g_3^\text{Dirac} = 40$ & spinor lift & graph data &
dimensionless & integer & exact \\
$B = 42$ Casimir trace at $m=3$ & --- & graph data &
dimensionless & integer & exact \\
$D_\text{even}(4)$ closed form & spinor lift (vertex selection) &
$\alpha$ & dimensionless & $\pi^2, \pi^4, G, \beta(4)$ & PSLQ 80~dps \\
$E_\text{Cas}^{S^3} = 1/240$ (conformal scalar) & Fock + rest-mass
projection (conformal value $m^2 = 1$) & none & dimensionless &
rational & exact symbolic \\
$E_\text{Cas}^{S^3, \text{Dirac}} = 17/480$ (full Dirac) & Fock +
Camporesi--Higuchi spinor lift & none & dimensionless & rational &
exact symbolic \\
$E_\text{Cas}^{\text{HO}, S^5} = -17/3840$ (3D HO, Bargmann--Segal) &
Bargmann--Segal $+$ spectral action ($\zeta$-regularization) & none &
dimensionless & rational & exact symbolic \\
Hydrogen Bohr levels $E_n = -Z^2/(2n^2)$ Ha & Fock $+$ energy-shell
rescaling & $Z, n$ & energy & rational & exact symbolic \\
Static Coulomb-like Green's fn $(L+I)^{-1}$ on Hopf $S^3$ at
$n_\text{max} = 2, 3$ & Fock (Layer-1 matrix inverse, no continuous
integration) & none & dimensionless & rational & exact symbolic at
finite $n_\text{max}$ \\
\midrule
\textbf{Two-projection chain} & & & & & \\
Stefan-Boltzmann $\pi^2/90$ on $S^3 \times S^1_\beta$ (high-$T$) &
Fock $+$ observation/temporal-window $+$ spectral action
($\zeta$-regularization) & $T$ or $\beta$ & energy density &
$\pi^2 \cdot \mathbb{Q}$ & exact symbolic \\
Uehling shift $-27.13$ MHz at 2$S$ & spectral action $\circ$ Fock &
$Z, \alpha$ & energy & $\pi^2{\cdot}\mathbb{Q}$ & $<1\%$ \\
21~cm hyperfine $A_\text{HF} = 1418.84$ MHz Bohr-Fermi (HF-1, $+$ external
$g_p$) & spinor lift (Fermi-contact vertex) $\circ$ Fock $\circ$ rest-mass
projection ($m_p$, proton) & $\alpha, g_p, m_e/m_p$ & energy &
$\alpha^2{\cdot}\mathbb{Q}$ & $-1102$ ppm (= $-a_e$) \\
21~cm hyperfine $A_\text{HF} = 1420.488$ MHz $+$ Parker--Toms $a_e$
(HF-2, $+$ external $g_p$) & spinor lift (Fermi-contact vertex) $\circ$
Fock $\circ$ rest-mass projection ($m_p$, proton) $\circ$ spectral
action (curvature expansion) & $\alpha, g_p, m_e/m_p$ & energy &
$\alpha^2{\cdot}\mathbb{Q}, c_1 \cdot \mathbb{Q}$ & $+58$ ppm \\
21~cm hyperfine $A_\text{HF} = 1420.432$ MHz $+$ external Zemach $r_Z$
(HF-4) & HF-2 chain $\circ$ rest-mass projection $\circ$
magnetization-density (\S~\ref{sec:proj_magnetization_density}) &
$r_Z$ & energy & calibration & $+18$ ppm (Eides Tab.~7.3 consistent) \\
$\mu$H 1S Bohr--Fermi $\nu_F(\mu p) = 182.4433$ meV (Sprint MH Track B,
$+$ external $g_p$) & spinor lift (Fermi-contact vertex) $\circ$ Fock
$\circ$ rest-mass projection ($m_\mu$, muon, lepton swap) $\circ$
rest-mass projection ($m_p$, proton) & $\alpha, g_p, m_\mu/m_e, m_e/m_p$
& energy & $\alpha^2 \cdot \mathbb{Q}$ &
$+2$ ppm vs Eides pure-QED $\nu_F = 182.443$ meV \\
$\mu$H 1S Zemach mass-enhancement factor $185.94 \times$ ep
(Sprint MH Track B; operator-level Track C, May 2026) & rest-mass
projection $\circ$ magnetization-density
(\S~\ref{sec:proj_magnetization_density}) $\circ$ Fock; \texttt{lepton\_mass}
parameter on \texttt{MagnetizationDensitySpec} eliminates the manual
mass-scaling step Track B used in test scripts (Pauli-string sum now
carries correct lepton-mass scaling natively for downstream muonic
VQE) & $r_Z, m_\text{red}$ & dimensionless & calibration &
$0.55\%$ vs Eides muonic target $-7300$ ppm; bit-identical to Track B
manual scaling (intrinsic gap is sub-leading Eides corrections, not
operator construction) \\
$\mu$H 2$S$--2$P$ Uehling $+205.0074$ meV full kernel
(Sprint MH Track A) & spectral action $\circ$ Fock $\circ$ rest-mass
projection ($m_e \to m_\mu$) & $Z, \alpha, m_\mu/m_e$ & energy &
$\pi^2{\cdot}\mathbb{Q}$ & $<1$ ppm vs Antognini~2013 \\
Mu 1S-2S transition $2{,}455{,}528{,}721$~MHz (rest-mass rescaling
from H experiment, sprint precision-catalogue 2026-05-08) & Fock $\circ$
rest-mass projection (lepton-unchanged, nucleus $m_p \to m_\mu$) &
$\alpha, m_e, m_\mu/m_e$ & frequency &
$\alpha^2{\cdot}\mathbb{Q}$ & $-0.11$ ppm vs Mu-MASS Crivelli~2018 \\
Mu 1S-2S transition $2{,}455{,}528{,}394$~MHz (framework Bohr$+$SE $+$
cross-register recoil) & Fock $\circ$ spectral action $\circ$ rest-mass
projection $\circ$ cross-register $V_{eN}$ (Roothaan) & $\alpha, m_e,
m_\mu/m_e$ & frequency & $\alpha^2{\cdot}\mathbb{Q}$ &
$-0.25$ ppm vs Mu-MASS \\
Mu 2$S_{1/2}$-2$P_{1/2}$ Lamb shift $1047.63$ MHz (framework one-loop
QED with full Uehling, sprint precision-catalogue 2026-05-08) &
Fock $\circ$ spectral action $\circ$ Sturmian $\circ$ rest-mass projection
& $\alpha, m_e, m_\mu/m_e$ & frequency & $\alpha^2{\cdot}\mathbb{Q}$ &
$+0.013\%$ vs Karshenboim 2005 theory $1047.49(5)$ MHz \\
Bethe log $\ln k_0(2S) = 2.786$ (LS-3) & Sturmian (acc.) $\circ$ Fock & $Z$ &
dimensionless & calibration & $-0.92\%$ \\
Bethe log $\ln k_0(1S) = 3.002$ (LS-3) & Sturmian (acc.) $\circ$ Fock & $Z$ &
dimensionless & calibration & $+0.60\%$ \\
D 1S hyperfine $\nu_\text{HFS}(D) = 327.3975$ MHz Bohr--Fermi strict
(sprint precision-catalogue 2026-05-08, $+$ external $g_d$) & spinor
lift (Fermi-contact vertex) $\circ$ Fock $\circ$ rest-mass projection
($m_d$, deuteron); $I=1$ multiplicity $3/2$ on
$A_\text{hf}\,\hat{I}\cdot\hat{S}$ & $\alpha, g_d, m_e/m_d$ & frequency &
$\alpha^2{\cdot}\mathbb{Q}$ & $+40$ ppm vs Wineland--Ramsey 1972
$327.384352522(2)$ MHz \\
Mu 1S hyperfine $\nu_\text{HFS}(\text{Mu}) = 4464.1905$ MHz framework
Bohr--Fermi $+$ Schwinger one-loop (sprint precision-catalogue 2026-05-08,
muonium triple closure) & spinor lift (Fermi-contact vertex) $\circ$ Fock
$\circ$ rest-mass projection (lepton-unchanged, nucleus $m_p \to m_\mu$);
leptonic nucleus, no QCD content & $\alpha, m_e, m_\mu/m_e$ & frequency &
$\alpha^2{\cdot}\mathbb{Q}$ & $+199$ ppm vs Liu 1999 $4463.302765(53)$ MHz;
residual cleanly LS-8a multi-loop QED $+$ recoil (no QCD competing budget) \\
He $2{}^3P_0$--$2{}^3P_1$ fine structure $+29{,}612.91$ MHz (Sprint~3
BF-D + Sprint~4 DD; Drake combining $(\tfrac{3}{50},-\tfrac{2}{5},
\tfrac{3}{2},-1)$, J-pattern from rank-$k$ 6j; sprint precision-catalogue
2026-05-08) & spinor lift $\circ$ Wigner $3j$ ($6j$-coupled) $\circ$
Wigner $3j$ (Gaunt) $\circ$ Fock; internal multi-focal (1s, 2p at
distinct $Z_\text{eff}$), bipolar $(k_1{=}0,k_2{=}2)$ direct +
$(k_1{=}1,k_2{=}1)$ exchange & $\alpha, Z$ & frequency &
$\alpha^2{\cdot}\mathbb{Q}$ & $-0.014\%$ vs NIST $+29{,}616.951(6)$ MHz;
below LS-8a $\alpha^3/\pi$ budget \\
He $2{}^3P_0$--$2{}^3P_2$ span $+31{,}844.07$ MHz (Sprint~3 BF-D
+ Sprint~4 DD; sprint precision-catalogue 2026-05-08) & same chain &
$\alpha, Z$ & frequency & $\alpha^2{\cdot}\mathbb{Q}$ &
$-0.20\%$ vs NIST $+31{,}908.131(6)$ MHz; within LS-8a one-loop QED budget \\
Ps 1S-2S transition $1{,}233{,}687{,}096.1$~MHz (framework-native: Bohr at
$m_\text{red}(ee) = 0.5$ $+$ Eides~\S\,3.2 SE Lamb-shift differential;
sprint precision-catalogue 2026-05-08, equal-mass leptonic 1S-2S; $+$
external Penin--Pivovarov $\alpha^4$ Breit) & Fock $\circ$ rest-mass
projection (equal-mass limit $\lambda_a = \lambda_b = 0.5$) $\circ$
spectral action (Eides) & $\alpha, m_e$ & frequency &
$\alpha^2{\cdot}\mathbb{Q}$ & $+64.75$ ppm vs Fee~1993
$1{,}233{,}607{,}216.4(3.2)$~MHz \\
\midrule
\textbf{Three-projection chain} & & & & & \\
Bethe log $\ln k_0(2P) = -0.0307$ (LS-4) & Drake--Swainson $\circ$ Sturmian
$\circ$ Fock & $Z$ & dimensionless & flow & $+2.4\%$ at $N=40$ \\
Bethe log $\ln k_0(3P) = -0.0406$ (LS-4) & Drake--Swainson $\circ$ Sturmian
$\circ$ Fock & $Z$ & dimensionless & flow & $+6.3\%$ at $N=40$ \\
Bethe log $\ln k_0(3D) = -0.005236$ (LS-4) & Drake--Swainson $\circ$ Sturmian
$\circ$ Fock & $Z$ & dimensionless & flow & $-0.24\%$ at $N=40$ \\
Lamb shift $1054$ MHz (LS-4 N=40 sweet spot) & Fock $\circ$ spectral action
$\circ$ Sturmian $\circ$ Drake--Swainson & $Z, \alpha$ & energy & calibration
$+$ flow & $-0.39\%$ (T$+$A cancellation) \\
Lamb shift $1029$ MHz (LS-3) & Fock $\circ$ spectral action $\circ$
Sturmian (acc.) & $Z, \alpha$ & energy & calibration & $-2.77\%$ (T$+$A)\\
$\mu$H 2$S$--2$P$ Lamb shift $202.17$ meV (framework $+$ literature,
Sprint MH Track A) & Fock $\circ$ spectral action $\circ$ Sturmian
$\circ$ Drake--Swainson $\circ$ rest-mass projection ($m_e \to m_\mu$)
& $Z, \alpha, m_\mu/m_e, r_p$ & energy & calibration $+$ flow &
$-0.10\%$ vs CREMA $202.3706(23)$ meV \\
Heisenberg--Euler 1-loop $\alpha^2/(45\pi m_e^4)$ &
spectral action (Schwinger Mellin) $+$ vector-photon (one loop)
& $\alpha, m_e$ & length${}^4 \cdot$ energy${}^{-1}$ &
$\alpha^2 \cdot \mathbb{Q} / \pi$ & structural (S) \\
Lyman $\alpha$ emission $A_{21}(2P{\to}1S) = 6.2683 \times 10^8$~s$^{-1}$
(Sprint Calc-L, 2026-05-09) & Fock $\circ$ Wigner $3j$ $\circ$
vector-photon promotion & $Z, \alpha, \omega$ & 1/time &
$\alpha^3 \cdot \mathbb{Q}[\sqrt{6}]$ & $+0.055\%$ vs NIST
$6.2649 \times 10^8$~s$^{-1}$; residual is non-relativistic
(Schr\"odinger vs Dirac at $(Z\alpha)^2 \sim 5\times 10^{-5}$),
closes by adding Camporesi--Higuchi spinor lift to the chain
(depth-prediction consistent) \\
H 1S static dipole polarizability
$\alpha_p(\text{H}, 1S) = \tfrac{9}{2}\, a_0^3$
(Sprint Calc-P, 2026-05-09) & Fock $\circ$ Wigner $3j$ $\circ$
Sturmian reparam.\ at $\lambda = Z$ & $Z$ & length${}^3$ &
rational & exact (Dalgarno--Lewis solution
$F|1s\rangle = -\tfrac{1}{2} z(r{+}2)|1s\rangle$ has $p$-wave radial
part lying exactly in span$\{S_2, S_3\}$ at $N_\text{basis} = 2$);
$Z^{-4}$ scaling exact at $Z = 1, 2, 3, 4$.  Bound-states-only sum
on bare Fock graph plateaus at $\sim 81.25\%$ of $\tfrac{9}{2}$
(Bethe--Salpeter 1957 partition); the missing 18.75\% continuum
content is recovered \emph{exactly} by the Sturmian projection.
\textbf{Counter-example to a strong reading of
Prediction~\ref{pred:error_depth}}: depth-3 chain at exact match
because the analytic answer is rational with finite-Sturmian-basis
representation. \\
\midrule
\textbf{Conjectural / Observation status} & & & & & \\
$K = \pi(B + F - \Delta) \approx 1/\alpha$ & Fock $\circ$ Hopf
$\circ$ spinor & $\alpha$ & dimensionless & three-sector mix &
$8.8\!\times\!10^{-8}$ \\
\bottomrule
\caption{The empirical matches catalogue.  Rows are grouped by
projection depth.  ``Trans.\ class'' uses Paper~18's tier names.
``Match'' is the precision against standard known values:
``exact'' or ``exact symbolic'' means equality at machine
precision; numerical percentages are residuals.  The bottom row
documents Paper~2's $\alpha$ conjecture explicitly as a
multi-projection coincidence; the combination rule itself is a
numerical observation without first-principles derivation
(Paper~18 \S~III.5; Paper~2 status v2.26.1).}
\label{tab:catalog}
\end{longtable}

\begin{remark}[signature $\ne$ value]
\label{rem:signature_value}
Theorem~\ref{thm:dimensional_consistency}'s signature consistency is
a type-level statement: chains that produce the same signature can
produce wildly different numerical values.  For example, the
hydrogen Bohr level $E_n = -Z^2/(2n^2)$ (Fock alone, signature
$\{$energy, no~$\pi\}$) and the hydrogen Lamb shift (Fock $+$
spectral action $+$ Sturmian $+$ Drake--Swainson, signature
$\{$energy, $\pi\}$) both have dimension $\{$energy$\}$, yet differ
by 13 orders of magnitude.  The depth-prediction of
\S~\ref{sec:prediction} addresses the value-level question;
signature arithmetic addresses only the type-level question.
\end{remark}

% -------------------------------------------------------------------
\subsection{Off-precision matches and error-source classification}
\label{sec:offprecision}
% -------------------------------------------------------------------

The catalogue above records matches at machine precision or stated
sub-percent agreement.  Equally informative is the catalogue of
matches that are \emph{off} precision, classified by the source of
residual error.  We propose five distinct error sources, each with
its own scaling behaviour and remediation pathway:

\begin{itemize}
\setlength\itemsep{0pt}
\item \textbf{Truncation (T):} finite $n_\text{max}$, $\ell_\text{max}$,
or basis size $N$.  Residual scales as a power law in the cutoff;
remediated by extending the cutoff.
\item \textbf{Basis quality (B):} structural limit of the chosen
angular or radial basis at fixed cutoff.  E.g.\ PK $\ell_\text{max}=2$
optimal operating point for composed LiH; higher $\ell$ degrades.
Remediated only by changing basis architecture.
\item \textbf{Approximation order (A):} leading-order treatment
omitting higher-order perturbative terms.  Residual is a known
fraction of the standard physics (e.g.\ one-loop QED carries
$\sim 5\%$ residual from $\alpha^5$ multi-loop).
\item \textbf{Calibration mismatch (C):} chained projections where
the intermediate calibration constant for the wrong diagram
topology is applied.  Confirmed in the Schwinger $\alpha/(2\pi)$
graph-native sprint (VP-1, VP-2): the vacuum-polarization projection
constant $C_{\text{VP}}$ does not transfer to vertex correction.
\item \textbf{Structural floor (S):} the framework cannot represent
the relevant physics at the chosen layer.  E.g.\ scalar Fock graph
cannot represent vector-photon QED selection rules without
projection promotion.  Remediated by adding the missing projection
to Layer~2.
\end{itemize}

\begin{table}[h]
\centering\small
\begin{tabular}{p{3.8cm} p{2.5cm} p{1.5cm} p{1.4cm} p{4.5cm}}
\toprule
Match & Projection(s) & Residual & Source & Notes \\
\midrule
He ground state (graph-native CI, $Z=2$) & Fock + Sturmian &
$0.20\%$ at $n_\text{max}=9$ & S & Small-$Z$ graph-validity-boundary
artifact (Sec.~2 of CLAUDE.md, v2.9.2 CUSP-2). \\
He ground state (2D variational + cusp) & Fock + spectral cusp &
$0.019\%$ at $\ell_\text{max}=7$ & T & Track DI Sprint 1; cusp-corrected
$0.004\%$ at $\ell_\text{max}=4$. \\
LiH $R_{\text{eq}}$ (composed, PK) & Fock $\circ$ molecule-frame
& $5.3\%$ at $\ell_\text{max}=2$ & B & PK $\ell_\text{max}=2$ optimal
operating point; higher $\ell$ degrades structurally
(Track BQ, Paper~17). \\
BeH$_2$ $R_{\text{eq}}$ & composed + exchange & $11.7\%$ & B &
Same PK structural ceiling as LiH; full 1-RDM exchange.
\\
H$_2$O $R_{\text{eq}}$ & composed + lone pairs & $26\%$ & B &
6{:}1 charge asymmetry overwhelms angular basis at $\ell_\text{max}=2$.
\\
H$_2$ $D_e$ (Level 4) & molecule-frame hyperspherical & $4.0\%$ & T
& $96.0\%$ of experimental $D_e$ at $\ell_\text{max}=6$; CBS
$\sim 97\%$. \\
LiH 4-electron FCI energy & balanced coupled $\circ$ Fock &
$0.20\%$ at $n_\text{max}=3$ & T & Improvement $1.8\% \to 0.20\%$
from $n_\text{max}=2 \to 3$; structural drift in $R_{\text{eq}}$
remains. \\
Schwinger $\alpha/(2\pi)$ from
graph-native vertex & vector-photon $\circ$ graph-native & $33\%$
level overshoot & C & Wrong projection constant; $C_{\text{VP}} \times
F_2$ diverges (sprint VP-1; CLAUDE.md graph-native QED bullet). \\
Anomalous moment $C \times F_2$ asymptotic & graph-native &
$12.7\times$ overshoot at $n_\text{max}=3$ & C & Vector promotion
gets $3.6\times$ closer to scalar baseline but does not close
projection mismatch (sprint VP-1). \\
Bethe logarithm $\ln k_0(2S) = 2.812$ & Sturmian (acc.) $\circ$ Fock &
$-0.92\%$ at $N=40$ & T & Acceleration form $3.3\times$ tighter than
LS-2 velocity at $N=50$; sprint LS-3. \\
Bethe logarithm $\ln k_0(1S) = 2.984$ & Sturmian (acc.) $\circ$ Fock &
$+0.60\%$ at $N=20$ & T & Acceleration form crosses Drake at smaller $N$
than velocity; sprint LS-3. \\
Bethe logarithm $\ln k_0(2P) = -0.0307$ & Drake--Swainson $\circ$
Sturmian $\circ$ Fock & $+2.4\%$ at $N=40$ & T & \textbf{LS-4 closes
the LS-3 structural floor.} New Drake--Swainson projection
(\S~\ref{sec:proj_drake_swainson}) produces $K$-independent value via
$\ln k_0 = J_v / D_{\text{drake}}$ with structural denominator
$D_{\text{drake}}(n,\ell) = 2(2\ell+1)Z^4/n^3$.  Convergence
monotonic: $+2.4\%$ at $N=40 \to +1.24\%$ at $N=55$. \\
Bethe logarithm $\ln k_0(3P) = -0.0406$ & Drake--Swainson $\circ$
Sturmian $\circ$ Fock & $+6.3\%$ at $N=40$ & T & Slowest-converging
$\ell>0$ case; sprint LS-4. \\
Lamb shift $1057.85$ MHz & Fock $\circ$ spectral action $\circ$
Sturmian (acc.) $\circ$ (Drake--Swainson, l=1) & $-0.39\%$ at LS-4
sweet spot $N=40$ & T cancels A & At $N=40$, $T_{2S} = +28.8$~MHz
truncation cancels $A = -32.78$~MHz one-loop ceiling; converges to
LS-1 baseline $-3.10\%$ at $N\to\infty$.  Sprint LS-4. \\
Lamb shift $1057.85$ MHz (LS-3 baseline) & Fock $\circ$ spectral action
$\circ$ Sturmian (acc.) & $-2.77\%$ (LS-3 prediction $1029$ MHz) &
\textbf{C} (resolved by LS-6a) $+$ A & Sprint LS-6a (May 2026)
re-derived the same one-loop physics in the Eides \S~3.2 canonical
convention and found the LS-1 implementation had inadvertently
subtracted the Uehling kernel constant from the canonical $10/9$ SE
coefficient (Uehling already enters separately as a VP contribution).
The convention fix supersedes this row: see the LS-6a row below for
the clean result. The $\sim +24.7$~MHz residual that had been
classified A is structurally a $\sim +27.13$~MHz one-loop convention
correction (\textbf{C} resolved), leaving the genuine A-tier
multi-loop residual of $\sim +5.65$~MHz. \\
Lamb shift $1057.85$ MHz (LS-6a, Eides \S~3.2 convention) & Fock
$\circ$ spectral action $\circ$ Sturmian (acc.) $\circ$ (Drake--Swainson,
$\ell=1$) & $-0.534\%$ (LS-6a prediction $1052.19$~MHz) & A (small) +
non-loop sectors (reframed by LS-7) & Eides \S~3.2 canonical one-loop
SE convention.  Convention shift from LS-1: $+27.13$~MHz, exactly
$(4/15) \cdot \alpha^3 Z^4 / (\pi n^3) \cdot \text{Ha-to-MHz}$ at
$n=2$, $Z=1$ (the Uehling kernel constant $10/9 - 38/45 = 4/15$).
Components: SE 2$S_{1/2}$ $+1066.44$~MHz (Bethe $+953.40$~MHz +
canonical $10/9$ piece $+113.04$~MHz); VP 2$S_{1/2}$ $-27.13$~MHz;
SE 2$P_{1/2}$ $-12.88$~MHz; VP 2$P_{1/2}$ $0$.  $5.8\times$
reduction in absolute error vs LS-3.  \textbf{Sprint LS-7 reframed
the residual decomposition}: the LS-5 / LS-6a attribution to
``multi-loop QED in Eides Tab.~7.4 band $[+3.5, +10.7]$~MHz'' was a
misreading of which Eides table is the proper multi-loop reference.
Eides Tab.~7.3 ($\alpha^5$ multi-loop QED proper) totals only
$\sim +1.20$~MHz; the remaining $\sim +4.4$~MHz of the $-5.65$~MHz
residual is non-loop physics (recoil $-2.40$, finite nuclear size
$+1.18$, hyperfine averaging $\sim +5.0$, sum $\sim +4.98$).  The
framework's one-loop QED machinery on $S^3$ is therefore even more
complete than originally framed: the multi-loop test target for
Paper~35 Prediction~1 is sharpened to $\sim +1.20$~MHz, with the
non-loop sectors covered by the Paper~34~\S\,III.14 rest-mass
projection and Paper~23's nuclear-electronic composed Hamiltonian
(not a test of Paper~35).  Sprint provenance:
\texttt{debug/ls6a\_eides\_convention\_memo.md},
\texttt{debug/ls7\_two\_loop\_se\_memo.md}. \\
Two-loop SE prefactor $(\alpha/\pi)^2 (Z\alpha)^4 / n^3$ on
Dirac-$S^3$ (LS-7, structural) & Fock $\circ$ spectral action
$\circ$ Sturmian (acc.) $\circ$ Camporesi--Higuchi $\circ$ iterated
proper time & structural prefactor confirmed to symbolic exactness;
dimensionless bracket $C_{2S} = +3.63$ used from Eides literature
(LS-8a target for native derivation) & A (sharpened) & Iterated CC
spectral action on Dirac-$S^3$ produces the standard two-loop SE
prefactor with the $1/\pi^2$ tracing to two iterated Schwinger
proper-time integrations, confirming Paper~35
Prediction~\ref{pred:error_depth} \emph{at the structural level} for
two-loop content.  Weak test (the same prefactor appears in
flat-space two-loop QED); the strong test is LS-8a's native
derivation of $C_{2S}$.  No structural obstruction encountered;
\texttt{geovac/qed\_two\_loop.py} infrastructure is adequate.
Sprint provenance: \texttt{debug/ls7\_two\_loop\_se\_memo.md}. \\
He $2{}^3P_1$--$2{}^3P_2$ small interval $+2{,}231.16$~MHz (Drake-pattern
Breit--Pauli + spinor; Sprint~3 BF-D + Sprint~4 DD + Sprint~5 CP) &
spinor lift $\circ$ Wigner $3j$ ($6j$-coupled, Drake) $\circ$ Wigner $3j$
(Gaunt) $\circ$ Fock &
$-2.62\%$ vs NIST $+2{,}291.176(15)$~MHz & A & Partial-cancellation
amplification: absolute residual $-60$~MHz on a $2.3$~GHz interval.
The dominant intervals (P$_0$-P$_1$ at $-0.014\%$, P$_0$-P$_2$ span at
$-0.20\%$) are in \S~\ref{sec:matches}; this small interval has the
same absolute residual but $\sim$10$\times$ larger fractional residual
because of the structural cancellation SO ($-58$~GHz) $+$ SS ($-9.5$~GHz)
$+$ SOO ($+70$~GHz) $\approx +2.2$~GHz.  Residual is below the LS-8a
$\alpha^3/\pi$ one-loop QED budget on the absolute scale; multi-electron
SS/SOO higher-order corrections (Drake 1971 \S~IV) are the standard
remediation.  Sprint precision-catalogue 2026-05-08. \\
Li 2$^2$P doublet (Sprint 5 CP) & Breit--Pauli + Slater $Z_{\text{eff}}$
& $+8.9\%$ & A & Standard convention with $Z_\text{val}=1$;
core polarization worsens accuracy (parity-forbidden
$2s2p \leftrightarrow 2p^2$ mixing confirmed). \\
21~cm $A_\text{HF} = 1418.84$ MHz (HF-1 reduced-mass, $+$ external $g_p$)
& spinor lift (Fermi-contact vertex) $\circ$ Fock $\circ$ rest-mass
projection ($m_p$, proton) & $-1102$ ppm (= $-a_e$) & A & Bohr-Fermi level;
$g_e = 2$ (no anomalous moment).  Residual is exactly $a_e \approx
\alpha/(2\pi)$, addressed by HF-2.  Sprint HF, 2026-05-07. \\
21~cm $A_\text{HF} = 1420.488$ MHz (HF-2, $+a_e$ Schwinger, $+$ external
$g_p$) & HF-1 chain $+$ graph-native $F_2$ on Dirac-$S^3$ & $+58$ ppm & A
& Schwinger
asymptote $a_e = \alpha/(2\pi)$ derived structurally on Dirac graph;
\textbf{Parker--Toms first-order curvature correction $c_1 = R/12 = 1/2$
verified at $0.5\%$ precision against the spectral sum}, no fits.  See
Paper~28 \S\ref{sec:curvature_coefficients}.  Residual $+58$ ppm is
recoil + Zemach + multi-loop. \\
21~cm $A_\text{HF} = 1420.432$ MHz (HF-4, $+r_Z$ external) & HF-2 chain
$+$ Zemach Layer-2 input ($r_Z = 1.045$ fm) & $+18$ ppm & A + S(Zemach)
& Convention catch: standard form is $\Delta\nu/\nu = -2 Z\alpha m_e r_Z$
(natural units), $-2 Z r_Z[\text{bohr}]$ (atomic units), giving
$-39.5$ ppm at $r_Z = 1.045$ fm.  Framework has no
magnetization-distribution operator on the proton register
(structural floor S); $r_Z$ is external focal length.  Final $+18$ ppm
matches Eides Tab.~7.3 attribution to Bodwin--Yennie recoil-beyond-
reduced-mass + multi-loop QED + nuclear polarizability. \\
Recoil from cross-register cross-spatial coupling (HF-3) & Track NI
nuclear-electronic & not derivable & S & \texttt{finite\_size\_correction}
returns scalar at classical \texttt{R\_PROTON\_BOHR}; 0/71 Pauli strings
couple spatial qubits across registers; the 12 cross-register terms are
all spin-spin hyperfine.  Framework couples discrete spin labels but
not spatial coordinates.  Sprint HF, 2026-05-07. \\
Multi-loop $a_2$ on Dirac-$S^3$ (HF-5) & iterated CC spectral action
$\circ$ Camporesi--Higuchi & $a_2$ extraction blocked & S (renormalization)
& Bare two-loop topologies (VP-on-photon, SE-on-electron, crossed) UV
diverge as $\sim N^{3.79}$.  Framework reproduces $(\alpha/\pi)^2$
prefactor, three-topology decomposition, SO(4) selection at all four
vertices, UV degree count, but does not autonomously generate $Z_2/Z_3/
\delta m$ counterterms.  Same structural wall as LS-8a (Lamb shift),
vertex sector. \\
$\mu$H 1S HFS $\nu_\text{HFS}(\mu p) = 181.32$ meV (BF + Schwinger +
leading Zemach, Sprint MH Track B) & rest-mass projection $\circ$
spinor $\circ$ Fermi contact $\circ$ Fock; magnetization-density
Layer-2 input $r_Z = 1.045$~fm & $-7710$ ppm vs Krauth full theory
182.725 meV & A (one-loop closure only) + S (LS-8a wall, muonic regime)
& Dominant residual is electron-VP (Uehling) in the muonic potential,
the leading muonic QED correction (Eides Tab.~7.4 / Karshenboim 2005:
$\Delta_\text{el-VP} \approx +1.5$ meV).  Same LS-8a wall as HF-5; in
muonic atoms the electron loop is enhanced because the muonic Bohr
radius probes the electron Compton wavelength.  Inner-factor input
data tier per Paper~18 \S~IV.6.  Sprint MH Track B, 2026-05-08. \\
Ps 1S HFS $116.79$~GHz (Fermi-contact only, framework-native) & spinor
lift (Fermi-contact vertex) $\circ$ Fock $\circ$ rest-mass projection
($m_p \to m_e$, equal-mass limit) & $-42.58\%$ vs Ishida 2014
$203.389$~GHz & A & Framework-native part captures only the magnetic
dipole-dipole contribution ($4/12$ of the LO HFS).  The
remaining $3/12$ is the annihilation channel $e^+e^- \to \gamma \to e^+e^-$,
a virtual-loop second-quantized effect not generated by the framework's
bare action -- structurally analogous to the K\"all\'en--Sabry two-loop
VP input for muonic H. Multi-focal architecture at $\lambda_a = \lambda_b$
verified: Roothaan multipole termination at $L_\text{max} = 2 \ell_\text{max}$
preserved (Gaunt selection is angular, not small-parameter).
Sprint precision-catalogue, 2026-05-08. \\
Ps 1S HFS $204.39$~GHz (LO total $= 7/12 \cdot m_e c^2 \alpha^4 / h$,
framework $+$ Layer-2 annihilation) & spinor lift (Fermi-contact vertex)
$\circ$ Fock $\circ$ rest-mass projection $+$ Layer-2 annihilation &
$+0.49\%$ vs Ishida 2014 & A & Residual is multi-loop QED
(Czarnecki--Melnikov--Yelkhovsky 2000, $\alpha^5 m_e c^2/h$ corrections).
LS-8a wall confirmed in the equal-mass leptonic regime: same
renormalization gap as $\mu$H Track A and HF-5.  Sprint precision-catalogue,
2026-05-08. \\
D 1S hyperfine $\nu_\text{HFS}(D) = 327.4779$ MHz (BF $+$ recoil
$+$ Schwinger $a_e$ $+$ leading Zemach with $r_Z(D) = 2.593$~fm
external; sprint precision-catalogue 2026-05-08) & spinor $\circ$ Fermi
contact $\circ$ Fock $+$ reduced-mass $+$ vertex correction
$+$ magnetization-density Layer-2 input
& $+286$ ppm vs Wineland--Ramsey 1972 $327.384352522(2)$ MHz & A
(framework-native) $+$ S (LS-8a wall, multi-loop QED) $+$ S (deuteron
polarizability, QCD-internal NN dynamics)
& Strict Bohr-Fermi at $+40$ ppm verifies $I=1$ cross-register angular
coupling structure: $A_{hf}\,\hat{I}\cdot\hat{S}$ Hamiltonian unchanged,
$3/2$ multiplicity is Clebsch--Gordan, no new mechanism vs $I=1/2$.
Cumulative residual $+286$ ppm budgeted to deuteron polarizability
($+44$ ppm Pachucki--Yerokhin 2010, QCD-internal), multi-loop QED
($\sim$few ppm, LS-8a wall), recoil NLO, finite-size charge.
Operator-level Zemach $-98.0$ ppm with $r_Z(D) = 2.593$~fm matches
analytic $-2 Z m_e r_Z$ to $10^{-14}$ ppm; profile (Gaussian vs
exponential) independence preserved at leading order, same as Sprint MH
Track C / hydrogen 21cm (Sprint HF Track 4).  Multi-focal architecture
verified at I=1: deuteron quadrupole moment couples only to gradients,
sub-ppm for s-states (negligible).  Sprint precision-catalogue,
2026-05-08. \\
Mu 1S-2S transition (test of rest-mass projection at unchanged lepton)
& Fock $\circ$ rest-mass projection (m$_p \to m_\mu$) & $-0.11$ ppm vs
Mu-MASS & A & Pure $m_\text{red}$ rescaling from H experimental;
sub-ppm verification of the rest-mass projection theorem (Paper~34
\S~III.14) at varied nucleus mass.  Differential-recoil contribution
(framework + cross-register $V_{eN}$ Roothaan) closes to $-0.25$ ppm.
Sprint precision-catalogue, 2026-05-08.  \\
$\mu$H 2$S$--2$P$ Lamb framework-native $200.50$ meV (no literature
inputs, Sprint MH Track A) & Fock $\circ$ spectral action $\circ$
Sturmian $\circ$ Drake--Swainson $\circ$ rest-mass projection
($m_e \to m_\mu$) & $-0.92\%$ vs CREMA $202.3706$ meV & A + S
(LS-8a wall, muonic regime)
& Framework reproduces dominant Uehling at $<1$ ppm vs Antognini and
Friar moment leading-order at $4.3\%$.  $-0.92\%$ residual decomposes
exactly into missing K\"all\'en--Sabry two-loop VP ($+1.51$ meV),
multi-loop ($+0.19$ meV), recoil next-to-leading ($-0.045$ meV), and
nuclear polarizability ($+0.013$ meV).  Same LS-8a renormalization
wall as HF-5 and Sprint MH Track B, vertex sector with electron loop
in muonic potential.  Plus literature-input total: $-0.10\%$ vs CREMA.
\textbf{Headline architectural result}: full Uehling kernel integration
(rather than contact-density approximation) handles the muonic
$\beta = 2 m_e a_\mu = 1.475$ regime where Bohr radius $\sim$ electron
Compton wavelength.  Sprint MH Track A, 2026-05-08. \\
Mu 2$S_{1/2}$-2$P_{1/2}$ Lamb $1047.63$ MHz vs Mariam~1982 experiment
$1054(22)$ MHz & Fock $\circ$ spectral action $\circ$ Sturmian $\circ$
rest-mass projection & $-0.6\%$ (within $\pm 22$ MHz exp uncertainty)
& A & Framework-native at $+0.013\%$ vs Karshenboim~2005 theory
$1047.49(5)$~MHz; $-0.6\%$ residual against Mariam dominated by 1982-era
experimental uncertainty, not framework limitation.  Karshenboim's
``Recoil NLO $-11.30$ MHz'' is absorbed into the framework's
$m_\text{red}$ Eides bracket (the $+1074.13$ MHz framework SE differs
from Karshenboim's $+1085.84$ MHz SE itemization by exactly the recoil
NLO).  Naively adding Karshenboim's recoil to the framework total
double-counts.  See \S~\ref{sec:convention_exposures} for catalogue
context (entry \S~\ref{sec:conv_karshenboim_recoil}).
Sprint precision-catalogue, 2026-05-08. \\
Mu 1S hyperfine framework Bohr--Fermi $+$ Schwinger one-loop
$4464.1905$ MHz vs Liu~1999 $4463.302765(53)$ MHz & spinor lift
(Fermi-contact vertex) $\circ$ Fock $\circ$ rest-mass projection
(lepton-unchanged, nucleus $m_p \to m_\mu$) & $+199$ ppm & A &
\textbf{Cleanest LS-8a isolation in the catalogue}: leptonic nucleus
(point antimuon, no QCD) means residual is purely multi-loop QED $+$
cross-register recoil NLO, no Zemach / polarizability / hadronic VP
budget interfering.  Karshenboim 2005 itemized total of corrections
($+4.703$ MHz, $\sim +1054$ ppm relative to BF strict) decomposes as
$\alpha^2(Z\alpha) \sim +0.17$ MHz (two-loop QED), $\alpha(Z\alpha)^2(m/M)^2
\sim +0.7$ MHz (recoil $+$ QED), hadronic VP $\sim +0.2$ MHz, plus the
dominant Schwinger one-loop already captured by the framework.  The
$+199$ ppm framework residual is the structural depth of the LS-8a wall
at $\alpha^2(Z\alpha)$ for Bohr--Fermi-class observables in clean
isolation: no QCD competing budget, both leptons receiving full
$\alpha/(2\pi)$ at one loop.  Adding Karshenboim's full Layer-2
itemization closes the gap to $<1$~ppm.  Sprint precision-catalogue,
2026-05-08. \\
Ps 1S-2S framework-native $+64.75$ ppm vs Fee~1993
$1{,}233{,}607{,}216.4(3.2)$~MHz (sprint precision-catalogue 2026-05-08,
equal-mass leptonic 1S-2S) & Fock $\circ$ rest-mass projection $\circ$
spectral action (Eides) & $+64.75$ ppm & A & \textbf{Sixth instance of
the multi-focal-composition wall, and the cleanest known.}  At equal
mass $\lambda_a = \lambda_b$ the standard $1/M_\text{nucleus}$ recoil
expansion is invalid; the dominant residual is the $\alpha^4$ Breit /
two-body-Dirac correction at the $\sim 80$~GHz scale, NOT multi-loop
QED ($\sim$10 MHz, three orders of magnitude smaller).  Same structural
wall as Sprint H1 Yukawa, LS-8a $Z_2/\delta m$, HF-3 recoil, HF-4
Zemach, HF-5 multi-loop $a_e$: framework couples discrete labels via
Wigner symbols cleanly, does not natively compose two Fock-style
projections at equal mass.  Cleanest known instance because the failure
surfaces at $\alpha^4$ -- one full order before the LS-8a renormalization
wall -- isolating two-body-projection failure from renormalization
failure.  This residual is the empirical anchor for the projection family
in \S~\ref{sec:proj_breit_retardation} (two-body Dirac / Breit
retardation, the sixteenth projection, added 2026-05-08 motivated by
this match).  Layer-2 $\alpha^4$ Breit input from Penin--Pivovarov 1998
(PRL~80, 2101) closes the chain to sub-MHz; cumulative match is partly
explained-by-construction (Breit value inferred as residual).  Honest
verdict: framework-native is at $+65$~ppm; all closure beyond Bohr$+$SE
Lamb is Layer-2.  Sprint precision-catalogue, 2026-05-08. \\
He $2{}^1S_0$--$2{}^3S_1$ exchange splitting $7183.3$~cm$^{-1}$ vs NIST
$6421.47$~cm$^{-1}$ (sprint precision-catalogue 2026-05-08, ss-only
sub-block at $n_\text{max}=11$; refreshed 2026-05-08 with fixed
\texttt{compute\_rk\_float} evaluator, trajectory bit-identical to
displayed 0.1~cm$^{-1}$ precision, largest shift 0.05~cm$^{-1}$) &
graph-native CI $\circ$ Fock &
$+11.86\%$ & B & \textbf{First multi-electron correlation-driven catalogue
entry.}  No relativistic, no QED, no cross-register content; pure $V_{ee}$
exchange in the angular momentum eigenbasis.  Splitting is \emph{less}
accurate than absolute energies (34$\times$ ratio: $\sim 12\%$ on splitting
vs $\sim 2.7$--$2.8\%$ on individual states) -- common-mode errors on
$2^3S$ and $2^1S$ have \emph{opposite} phase impact, no cancellation.
\textbf{Variational bound violated for both excited states} (not for
$1^1S$): graph $\kappa = -1/16$ adjacency adds constant inter-shell
coupling regardless of orbital extent, overbinding loosely-bound
Rydberg-like states.  Same mechanism as the small-$Z$
graph-validity-boundary artifact ($Z_c \approx 1.84$, He at $Z=2$).
Single-exponent hydrogenic basis at $k = Z = 2$ cannot represent
$Z_\text{eff} \approx 1$ screening of 2s; multi-focal composition
(Phase~C-W1c, May~2026) is the structural extension.  Sprint
precision-catalogue, 2026-05-08. \\
Self-consistent $r_Z$ global extraction (Sprint Calc-rZG, 2026-05-09;
Layer-2 corrected 2026-05-09; extended-fit diagnostic 2026-05-09)
& §III.18 magnetization-density global fit across H 21cm $+$
$\mu$H 1S HFS $+$ D 1S HFS & \emph{Original 3-obs fit} (Layer-2 corrected,
$\sigma_{\mu H}$=1500 ppm, ``LS-8a wall scale''): $r_Z(p) = 1.180(166)$ fm
[$+0.80\sigma$ above Eides $1.045(20)$, $+1.00\sigma$ above lattice
$1.013(10)(12)$], $r_Z(D) = 2.583(1588)$ fm [$-0.01\sigma$ from
Friar--Payne $2.593(16)$], $\chi^2/\text{dof} = 0.43/1$.
\emph{Extended-fit diagnostic} (sprint Calc-rZG-extended, $\sigma_{\mu H}$
hardened to literature kernel scale ${\sim}367$ ppm): $r_Z(p) =
1.257(51)$~fm, $\sigma$ closes to $51$~mfm but extracted central
value sits $+3.85\sigma$ above Eides and $+4.55\sigma$ above lattice.
& C & \textbf{Framework $r_Z(p)$ extraction is currently kernel-limited,
not Layer-2-itemization-limited}.  Original 3-obs fit at $\sigma=166$~mfm
(driven by $\sigma_{\mu H}=1500$~ppm conservative) lands within $1\sigma$
of all three literature determinations after Layer-2 budget correction
(2026-05-09 bug fix; original rZG memo specified Layer-2$_D$ = $-150$ ppm
based on a mis-itemization, corrected to $-286$ ppm per Pachucki--Yerokhin
2010 + Friar--Payne 2005 closure analysis; Layer-2$_H$ tightened from
$0$ to $-18.5$ ppm per Eides Tab.~7.3); $\chi^2/\text{dof} = 0.43/1$
confirms internal fit consistency.  \textbf{The original $r_Z(D) = 6.18$
fm artifact was traced to a Layer-2 budget misspecification, NOT to a
recoil double-counting in the production cross-register $V_{eN}$
kernel.}  Diagnostic verification
(\texttt{debug/rzg\_bug\_diagnosis\_compute.py}): the production
\texttt{geovac.cross\_register\_vne.cross\_register\_recoil\_correction}
reproduces leading-order Bethe--Salpeter recoil at $2.86\%$ for
hydrogen, $2.03\%$ for deuterium, and $8.18\%$ for muonium without
any species-specific code path; it is species-agnostic and
correct.  Six new D and muonium tests added to
\texttt{tests/test\_cross\_register\_vne.py} document and protect
this property going forward.  \textbf{Extended-fit diagnostic
finding (2026-05-09 follow-on)}: tightening $\sigma_{\mu H}$ from
1500~ppm to either 150~ppm (Krauth 2017 multi-loop QED itemization) or
to 367~ppm (5\% kernel approximation in muonic systems) drives
$\sigma(r_Z(p))$ below 50 mfm but extracts a central value
$r_Z(p)\sim 1.18$--$1.26$~fm, INCONSISTENT with both Eides
$1.045(20)$ and lattice $1.013(16)$ at $>3.85\sigma$.  The +0.22~fm
offset between the muH-alone extraction ($r_Z=1.265$~fm) and Eides
($1.045$~fm) is the structural fingerprint of recoil-mixing in the
literature Zemach formula that the framework's leading-order Eides
kernel $-2Zm_\text{red}r_Z$ does not yet capture
(Sprint MH Track B \S2.4--2.5 W1b operator-level extension point).
\textbf{Verdict: framework currently \emph{cannot} resolve the
Eides-vs-lattice 30~mfm tension at the operator-level precision of
the leading-order Eides Zemach kernel.}  Path to sub-30~mfm: extend
\texttt{magnetization\_density.py} to include recoil-mixing
(Friar--Payne 2005 \S III treatment), expected to close the +0.22~fm
muH-alone offset and bring $\sigma(r_Z(p))$ to ${\sim}30$~mfm.
\textbf{Genuine prospective contribution}: §III.18-based global fit
operator-level-localizes the ``Layer-2 itemization is the load-bearing
systematic for atomic $r_Z$ extraction'' finding --- a unified
analytical lens not currently produced by standard atomic-physics
compilations.  Note that the diagnostic exposes a SECOND load-bearing
systematic: the kernel approximation order of the framework's W1b
operator dominates over Layer-2 itemization in muonic regime
extraction.  §III.18 itself reproduces Eides analytical
$-2 Z\alpha m_e r_Z$ to floating-point precision.  Source memos:
\texttt{debug/calc\_track\_rZG\_global\_zemach\_memo.md},
\texttt{debug/rzg\_bug\_diagnosis\_memo.md} (bug fix),
\texttt{debug/calc\_track\_rZG\_extended\_memo.md} (extended-fit
diagnostic).  \textbf{Sprint W1b recoil-mixing extension
(2026-05-09 follow-on)}: \S III.18 operator extended with the
NLO recoil-mixing prefactor $m_l/(m_l + m_n)$ (per arXiv:2604.06930
eq.~95) and the Friar moment correction $(1/2)(Z\alpha m_\text{red})^2
\langle r^2\rangle_{(2)}$ (per Friar 1979 / Eides \S6.2), gated behind
opt-in \texttt{include\_recoil\_mixing} flag with backward-compatibility
preserved bit-identical (37 baseline tests, 20 new tests, 57/57 pass).
With NLO kernel enabled, the muH 1S HFS Zemach line at $r_Z=1.045$
fm reproduces Krauth~2017's full-theory line to $\sim 1\%$ (vs $\sim
3\%$ overshoot at LO).  \textbf{rZG-extended-v2 fit result}: with
NLO kernel and Layer-2 re-anchored at the Eides reference,
$\sigma(r_Z(p))$ tightens from 51~mfm to \textbf{16~mfm} (below the
30~mfm Eides-vs-lattice gap), but the central value \emph{remains}
at $r_Z(p) \approx 1.286$~fm --- only $\sim 10\%$ of the v1
$+0.22$~fm offset closes.  \textbf{Diagnostic verdict revised}: the
$+0.22$~fm offset is NOT kernel-approximation; it is a Layer-2
itemization-convention mismatch between Krauth~2017 (used for muH
Layer-2) and Eides~2024 Tab.~7.3 (used for H~21cm Layer-2), with
the convention difference at $\sim 0.01\%$ of $\nu_F$ propagating
to $\sim 25$~mfm in extracted $r_Z(p)$.  Next-step recommendation:
Layer-2 convention reconciliation sprint, NOT further kernel
extensions.  See \S~\ref{sec:convention_exposures} for catalogue
context (entry \S~\ref{sec:conv_rZG_DHFS}).  Source memo:
\texttt{debug/w1b\_recoil\_mixing\_extension\_memo.md}. \\
\bottomrule
\end{tabular}
\caption{Off-precision matches with error-source classification.
Residual error is the absolute relative deviation from the standard
physics value.  Source codes: T (truncation), B (basis quality),
A (approximation order), C (calibration mismatch), S (structural
floor).  Combinations indicate multiple independent contributions
to residual error.  This table is a complement to
Sec.~\ref{sec:catalog}; together they document the full empirical
profile of the framework.}
\label{tab:offprecision}
\end{table}

\begin{remark}[Error sources are projection-specific]
\label{rem:err_sources}
The error-source classification is not orthogonal to the projection
dictionary.  Truncation (T) errors are properties of the spectral
sums that the Sturmian or spectral-action projections invoke.
Calibration mismatch (C) is a property of how multiple projections
chain.  Structural floors (S) appear where a needed projection is
absent from Layer~2 entirely.  Approximation order (A) is the
``flow'' tier of Paper~18: it is the residual after Layer~1 and
Layer~2 do their work, and it indicates which higher-order projection
or perturbation order should be added next.
\end{remark}

% -------------------------------------------------------------------
\subsection{Roothaan autopsies: multi-component decompositions}
\label{sec:autopsies}
% -------------------------------------------------------------------

The matches catalogue rows in \S\ref{sec:catalog} and
\S\ref{sec:offprecision} each capture a single observable as a
single row.  For observables that admit a multi-component literature
decomposition --- e.g.\ Eides Tab.~7.3 for the hydrogen Lamb shift,
where the standard $1057.85$ MHz value decomposes into one-loop SE,
one-loop VP, multi-loop QED, finite nuclear size, recoil NLO, and
hyperfine averaging --- the single-row treatment hides the structure
that the projection-chain dictionary makes visible.  This subsection
hosts \emph{Roothaan autopsies}: decomposition tables that show each
measured observable as a sum of focal-length-tagged components, with
each component's projection chain made explicit.

The Roothaan autopsy is structurally the same kind of object Roothaan
(1951) made for two-electron molecular integrals: a multi-component
decomposition where each piece has its own focal length and its own
closed-form structure.  Applying the same discipline to multi-component
\emph{observables} makes the projection-chain dictionary's reach
visible at the observable level, not just at the catalogue-row level.
The autopsy partitions each component into framework-native (FN ---
computable from \S~\ref{sec:layer2} projections without external
numerical constants beyond $\alpha$) versus Layer-2 (L2 --- parameter
at operator level, projection chain framework-native), giving an
explicit operator-level reading of the structural-skeleton scope
statement (CLAUDE.md \S~1.7).  Several autopsies surface
literature convention mismatches (class~(a) findings of the
multi-observable focal-length decomposition program); these are
collected separately in \S~\ref{sec:convention_exposures}.

\subsubsection{Hydrogen 1S Lamb shift autopsy}
\label{sec:autopsy_lamb}

\textbf{Reference.} $\nu(2S_{1/2} - 2P_{1/2}) = 1057.845(9)$ MHz
(Lundeen--Pipkin 1981; NIST ASD).

\begin{table}[h]
\centering\small
\begin{tabular}{p{4cm} r p{6cm} c}
\toprule
Component & MHz & Projection chain & Status \\
\midrule
SE 2$S_{1/2}$ (Bethe + Eides 10/9) & $+1066.44$ &
\S\ref{sec:proj_fock} $\circ$ \S\ref{sec:proj_sturmian} $\circ$
\S\ref{sec:proj_spectral_action} $\circ$ \S\ref{sec:proj_drake_swainson}
& FN \\
VP 2$S_{1/2}$ (Uehling) & $-27.13$ &
\S\ref{sec:proj_fock} $\circ$ \S\ref{sec:proj_spectral_action} & FN \\
SE 2$P_{1/2}$ & $-12.88$ &
\S\ref{sec:proj_fock} $\circ$ \S\ref{sec:proj_sturmian} $\circ$
\S\ref{sec:proj_spectral_action} $\circ$ \S\ref{sec:proj_spinor}
$\circ$ \S\ref{sec:proj_drake_swainson} & FN \\
VP 2$P_{1/2}$ & $0$ &
\S\ref{sec:proj_fock} (parity zero) & FN \\
$\alpha^5$ multi-loop QED & $+1.20$ &
\S\ref{sec:proj_fock} $\circ$ \S\ref{sec:proj_sturmian} $\circ$
$(\S\ref{sec:proj_spectral_action})^2$ $\circ$ \S\ref{sec:proj_spinor}
$+$ Layer-2 $C_{2S}$ & L2 \\
FNS $r_p$ contribution & $+1.18$ &
\S\ref{sec:proj_fock} $\circ$ \S\ref{sec:proj_spinor} $\circ$
\textbf{\S\ref{sec:proj_charge_density}} $+$ Layer-2 $r_p$ & L2 \\
Recoil NLO (beyond reduced-mass) & $-2.40$ &
\S\ref{sec:proj_fock} $\circ$ \S\ref{sec:proj_restmass} $\circ$
\S\ref{sec:proj_breit_retardation} $+$ Layer-2 $R_{nl}$ & L2 \\
Hyperfine averaging (F-state mixing) & $\sim +5.00$ &
\S\ref{sec:proj_fock} $\circ$ \S\ref{sec:proj_spinor} $\circ$
\S\ref{sec:proj_3j} $\circ$ \S\ref{sec:proj_restmass} & FN \\
\midrule
\textbf{Sum} & $+1057.41$ & & \\
Experimental & $+1057.845(9)$ & & \\
\textbf{Residual} & $+0.43$ & within Eides Tab.~7.3 quoting precision & \\
\bottomrule
\end{tabular}
\caption{Hydrogen 1S Lamb shift Roothaan autopsy.  Each component is
tagged to its \S\ref{sec:layer2} projection chain.  Status: FN
$=$ framework-native; L2 $=$ Layer-2 input.  Source memo:
\texttt{debug/calc\_track\_LAR\_lamb\_autopsy\_memo.md}.}
\label{tab:autopsy_lamb}
\end{table}

\textbf{Framework-native subtotal:} $+1031.43$ MHz ($97.5\%$ of
measurement).  \textbf{Layer-2 net:} $+1.20 + 1.18 - 2.40 \approx
-0.02$ MHz --- empirical near-cancellation specific to the H 1S
Lamb shift, \emph{not} a generic property.  The framework-native
components alone reproduce the measurement at $\approx 4 \times
10^{-4}$ precision; the Layer-2 inputs cancel in this observable.

\textbf{Multi-focal-composition wall visible.}  The two non-cancelling
Layer-2 components correspond to:
\begin{itemize}\setlength\itemsep{0pt}
\item Component 5 (multi-loop QED) is the LS-8a renormalization-gap
wall (CLAUDE.md \S~1.7 LS-8a entry).  The framework reproduces the
two-loop SE structural prefactor $(\alpha/\pi)^2 (Z\alpha)^4/n^3$
but cannot autonomously generate the $Z_2/\delta m$ counterterms.
\item Component 7 (recoil NLO) is the
\S\ref{sec:proj_breit_retardation} two-body kinematic wall.  The
framework's Breit retardation projection produces the $\alpha^4$
kinematic structure but the $R_{nl}$ coefficients enter as
Layer-2 input.
\item Component 6 (FNS) is now \emph{framework-native at the
operator level} (today's \S\ref{sec:proj_charge_density} promotion);
the remaining Layer-2 input is the scalar $r_p$ value, not the
operator structure consuming it.  This is the cleanest beneficiary
of the 2026-05-09 promotion.
\end{itemize}

\textbf{Structural reading.}  The autopsy demonstrates Paper~34
\S\ref{sec:layer2}'s reach at finer grain than single-row catalogue
entries: eight components, eight different projection chains, every
Eides Tab.~7.3 row reproduced.  The structural-skeleton scope
statement is made concrete --- framework-native pieces reproduce the
observable to $\approx 4 \times 10^{-4}$, with the remaining $0.4\%$
split between renormalization (Component 5) and two-body kinematics
(Component 7), and FNS (Component 6) now in projection-chain language
rather than as a black-box external constant.

\subsubsection{Hydrogen 21~cm hyperfine autopsy}
\label{sec:autopsy_21cm}

\textbf{Reference.} $\nu_\text{HF} = 1\,420\,405\,751.768$~Hz (Hellwig
1970; CODATA / NIST), the most precisely measured atomic transition.

\begin{table}[h]
\centering\small
\begin{tabular}{p{4cm} r p{6cm} c}
\toprule
Component & MHz & Projection chain & Status \\
\midrule
Bohr--Fermi Dirac (point nucleus, $g_e=2$, no recoil) & $+1421.1595$ &
\S\ref{sec:proj_fock} $\circ$ \S\ref{sec:proj_spinor} $\circ$
\S\ref{sec:proj_3j} & FN \\
$+$ Schwinger $a_e$ (Parker--Toms-verified at $+0.5\%$) & $\times (1+\alpha/2\pi)$ &
\S\ref{sec:proj_fock} $\circ$ \S\ref{sec:proj_spinor} $\circ$
\S\ref{sec:proj_spectral_action} & FN (with calibration) \\
$+$ Reduced-mass $(1+m_e/m_p)^{-3}$ & $\times 0.998366$ &
\S\ref{sec:proj_fock} $\circ$ \S\ref{sec:proj_restmass} & FN \\
$+$ Zemach $r_Z=1.045$~fm via \S\ref{sec:proj_magnetization_density}
operator-level & $\times (1 - 39.495 \text{~ppm})$ &
\S\ref{sec:proj_fock} $\circ$ \S\ref{sec:proj_spinor} $\circ$
\textbf{\S\ref{sec:proj_magnetization_density}} & FN at op-level $+$ L2 ($r_Z$) \\
\midrule
\textbf{Final $A_\text{HF}$} & $+1420.4318$ & & \\
Experimental & $+1420.4058$ & (12-digit precision) & \\
\textbf{Residual} & $+0.026$ & $= +18.4$~ppm; within Eides Tab.~7.3 itemized $+12$--$+18$~ppm band & \\
\bottomrule
\end{tabular}
\caption{Hydrogen 21~cm hyperfine Roothaan autopsy.  Each component is
tagged to its \S\ref{sec:layer2} projection chain.  Status: FN $=$
framework-native; L2 $=$ Layer-2 input.  Source memo:
\texttt{debug/h21\_autopsy\_v1\_memo.md} (Sprint Calc-H21-Autopsy v1, May~2026).}
\label{tab:autopsy_21cm}
\end{table}

\textbf{Operator-level Zemach.}  Component 4 exercises
\S\ref{sec:proj_magnetization_density} at the operator level via the
production module \texttt{geovac.magnetization\_density}.  The bilinear
matrix element $\langle \hat{r}_Z\rangle$ on a Sturmian register at
the L=0 multipole reduces to $M_1[\rho_M] = r_Z$ at leading order in
$r_Z/a_0 \sim 2\times 10^{-5}$; the operator output $-39.495$~ppm at
$r_Z = 1.045$~fm vs Eides Tab.~7.3 reference $-39.500$~ppm gives
residual $+4.7\times 10^{-3}$~ppm $= 0.012\%$ of the LO shift --- the
operator collapses to the Eides scalar at sub-percent precision
automatically through the multipole reduction, without any external
substitution.  Profile independence verified: Gaussian and exponential
$\rho_M$ give bit-identical $A_\text{HF}$ at LO since both are
calibrated to first moment $M_1 = r_Z$.

\textbf{Framework-native subtotal:} $A_\text{HF}^\text{native}\!=\!1420.4318$~MHz
($99.998\%$ of measurement).  \textbf{Layer-2 net:} $-39.495 + 18.37
\approx -21.1$~ppm of $\nu_F$.  The Layer-2 input is the scalar $r_Z$
value; the operator structure consuming it is framework-native via
\S\ref{sec:proj_magnetization_density}.  This is the cleanest
beneficiary of the W1b operator-level extension (Sprint
Calc-rZG-extended, 2026-05-09), which promoted Zemach from
analytic-substitution to operator-level.

\textbf{Multi-focal-composition wall visible.}  The $+18.37$~ppm
residual decomposes per Eides Tab.~7.3 as:
\begin{itemize}\setlength\itemsep{0pt}
\item Multi-loop QED ($\alpha^2(Z\alpha)$, $\sim+6$~ppm) is the
LS-8a renormalization gap (CLAUDE.md \S~1.7 LS-8a entry).
\item Recoil NLO beyond reduced-mass (Bodwin--Yennie, $\sim +6$~ppm)
is the W1a-D Roothaan kernel-level recoil-mixing wall.
\item Nuclear polarizability $\Delta_\text{pol}$ ($+1.4$~ppm) and
hadronic VP ($+0.1$~ppm) are W3 inner-factor calibration data
(QCD-internal).
\item Zemach NLO recoil-mixing $m_e/(m_e+m_p) \approx 5 \times 10^{-4}$
is structural noise in the electronic regime ($+0.02$~ppm) but the
\emph{dominant} systematic in the muonic regime
($f_\text{recoil} \approx 0.092$, structural $\sim 10\%$).  The W1b
NLO opt-in flag covers this cleanly; cf. Sprint MH for the muonic
analog.
\item Convention drift on $r_Z$ ($\pm 5$~ppm; Eides 1.045 vs
Karshenboim 1.054) is the largest class~(a) literature-itemization
sensitivity in this observable.
\end{itemize}

\textbf{Structural reading.}  The autopsy demonstrates Paper~34's
multi-observable focal-length decomposition program (CLAUDE.md
\S~1.8) at the operator level on the most precisely measured atomic
transition: four components, four different projection chains, $+18$
ppm residual fully attributable to LS-8a (multi-loop QED) $+$ W1a-D
(Roothaan recoil) $+$ W3 (nuclear polarizability), exactly as
classified in CLAUDE.md \S~1.7.  The structural-skeleton scope
statement is verified: the framework reproduces the observable to
$\sim 2 \times 10^{-5}$ precision via four named projections, and
the residual budget decomposes cleanly into already-named walls
without surfacing any new class~(a) literature convention mismatch
beyond the known $r_Z$-value drift.

\subsubsection{Muonic hydrogen 2$S_{1/2}$--2$P_{1/2}$ Lamb shift autopsy}
\label{sec:autopsy_muh_lamb}

\textbf{Reference.} $\Delta E = 202.3706(23)$ meV (CREMA 2010,
Pohl~et~al.; muonic convention $E(2P_{1/2}) - E(2S_{1/2})$).

The CREMA 2010 muonic hydrogen Lamb shift decomposes under five
named projection-chain components plus three companion Layer-2
inputs.

\begin{table}[h!]
\centering\small
\caption{Five-component Roothaan decomposition of the muonic hydrogen
Lamb shift (muonic convention $E(2P_{1/2}) - E(2S_{1/2})$).  The
framework-native subtotal $+200.50$ meV comes from full Uehling
integration via \S\ref{sec:proj_spectral_action}, the self-energy
bracket under rest-mass projection
(\S\ref{sec:proj_restmass}~+~\S\ref{sec:proj_sturmian}), and the
operator-level Friar moment via the magnetization-density module
(\S\ref{sec:proj_magnetization_density}; cf.\
\S\ref{sec:proj_charge_density} for the parent charge-density
projection).  The $+1.67$ meV gap to the $+202.17$ meV
framework-plus-literature total is closed by Antognini~2013
literature inputs covering multi-loop QED (K\"all\'en--Sabry,
mixed VP-VP/VP-SE, $\alpha^7$) and QCD-internal proton
polarizability (W3 inner-factor).  Source memo:
\texttt{debug/muH\_Lamb\_autopsy\_v1\_memo.md}.}
\label{tab:autopsy_muh_lamb}
\begin{tabular}{p{4.5cm} r p{4.8cm} l}
\toprule
Component & Value (meV) & Projection chain & Status \\
\midrule
1.~Full Uehling VP & $+205.0074$ &
\S\ref{sec:proj_fock} $\circ$ \S\ref{sec:proj_spectral_action}
& FN ($-0.10$ ppm vs Antognini) \\
2.~Self-energy (Eides bracket) & $-0.8295$ &
\S\ref{sec:proj_fock} $\circ$ \S\ref{sec:proj_sturmian}
$\circ$ \S\ref{sec:proj_restmass}
& FN ($+24\%$ vs Antognini, recoil-mixing gap) \\
3.~Friar moment (closed-form) & $-3.6752$ &
\S\ref{sec:proj_fock} $\circ$ \S\ref{sec:proj_charge_density}
$/$ \S\ref{sec:proj_magnetization_density}
& FN ($-4.3\%$ vs Antognini) \\
4.~K\"all\'en--Sabry 2-loop VP & $+1.5081$ & Layer-2 input &
LS-8a renormalization wall \\
5.~Proton polarizability & $+0.0129$ & Layer-2 input &
W3 inner-factor (QCD) \\
Companions (recoil $+$ multi-loop) & $+0.144$ & Layer-2 input &
Antognini 2013 catalogue \\
\midrule
\textbf{Total} & $+202.17$ & \multicolumn{2}{l}{Residual $-0.10\%$
vs CREMA $+202.3706(23)$ meV} \\
\bottomrule
\end{tabular}
\end{table}

\textbf{Headline framework computation.}  The full Uehling integration
matches the Antognini~2013 / Pachucki~1996 reference at $-0.10$ ppm
--- the cleanest single-component match in the muonic catalogue.
The $\beta = 2/(m_{\text{red},\mu}\,\alpha) = 1.475$ muonic regime,
where the muonic Bohr scale and the $e^+ e^-$ Compton wavelength are
comparable, is handled correctly by the full kernel where the
contact-form Uehling overshoots by $\sim 3.6\times$.  This is the
muonic-regime analog of the electronic Uehling row in
\S\ref{sec:autopsy_lamb} (Component 2 of \S\ref{tab:autopsy_lamb}),
where contact form is sufficient because $\beta_\text{ep} \approx 274$.

\textbf{Convention-mismatch surface.}  Antognini~2013 and Krauth~2017
itemizations agree on the bottom-line total but differ at $\sim 100$
ppm in per-component split (VP one-loop and SE-recoil aggregation).
Specifically, Krauth absorbs $\sim 0.021$~meV of higher-order VP into the
one-loop line ($+205.0282$~meV), while Antognini itemizes it separately
as ``VP-VP mixed'' ($+0.150$~meV), with one-loop reported at
$+205.0074$~meV.  The framework's $-0.10$~ppm match against Antognini
is $1000\times$ tighter than the $102$~ppm Antognini-vs-Krauth gap.
At this precision level the choice of Layer-2 reference compilation
matters; this autopsy anchors throughout to Antognini~2013.  See
\S~\ref{sec:convention_exposures} for catalogue context (entry
\S~\ref{sec:conv_muH_VP}).

\textbf{Operator-level Friar moment.}  Computed via the
\S\ref{sec:proj_magnetization_density} module
\texttt{geovac.magnetization\_density} with
\texttt{lepton\_mass}~=~$m_{\text{red},\mu}$.  The operator-level
result $-4.33$ meV vs closed-form $-3.68$ meV reflects the Gaussian
profile's $\langle r^2\rangle / \langle r\rangle^2 = 4/\pi \approx
1.27$ factor --- a profile-convention issue between calibrating to
the first moment $\langle r\rangle$ versus the RMS radius
$\sqrt{\langle r^2\rangle}$.  Eides Eq.~2.35 closed-form derivation
takes $\langle r^2\rangle^{1/2} = r_p$ (RMS radius); the operator-level
realization via \S\ref{sec:proj_magnetization_density} takes
$\langle r\rangle = r_p$ (first moment).  The discrepancy is exactly
the Gaussian profile factor.  The flagged extension is a sibling
\texttt{MagnetizationDensitySpec} calibrated to RMS radius for the
Lamb-shift FNS channel, distinct from the existing first-moment
calibration for the HFS Zemach channel.  See
\S~\ref{sec:convention_exposures} for catalogue context (entry
\S~\ref{sec:conv_friar_profile}).

\textbf{Structural reading.}  The autopsy realizes the
structural-skeleton-scope statement at the operator level for a
cross-register multi-focal observable in the muonic regime.  Three
framework-native components reproduce $+200.50$ meV (the
five-projection-chain skeleton); three Layer-2 inputs close the
remaining $+1.67$ meV gap, attributable cleanly to LS-8a multi-loop
renormalization (Component 4) and W3 inner-factor (Component 5).
The same architecture that closes the H 1S Lamb shift at
\S\ref{sec:autopsy_lamb} closes $\mu$H 2$S$--2$P$ here, with
Uehling-kernel kinematics (contact-form $\to$ full kernel) and
Friar-profile calibration (first moment $\to$ RMS) being the only
two ingredients that change between regimes.

\subsubsection{Helium $2{}^3P$ fine-structure autopsy (first internal multi-focal entry)}
\label{sec:autopsy_he_2_3p}

\textbf{Reference.}  $\nu(2{}^3P_0 - 2{}^3P_1) = +29\,616.951(6)$~MHz,
$\nu(2{}^3P_1 - 2{}^3P_2) = +2\,291.176(15)$~MHz,
$\nu(2{}^3P_0 - 2{}^3P_2) = +31\,908.131(6)$~MHz (NIST compilation;
Pachucki--Yerokhin 2010 theoretical at sub-kHz vs experiment).

This is the first \emph{internal multi-focal} Roothaan autopsy: two electrons
on the same nucleus at structurally distinct effective nuclear charges
($Z_\text{eff}(1s) = 2$ full-nuclear, $Z_\text{eff}(2p) = 1$ full-shield).
Distinct from the cross-register multi-focal autopsies of \S\ref{sec:autopsy_lamb}
(electron $+$ I$=$1/2 nucleus) and \S\ref{sec:autopsy_muh_lamb} (electron $+$
muonic nucleus), which couple particles on different registers.

The operator-level Drake (1971) decomposition gives three components:
\begin{equation}
E({}^3P_J) = \tfrac{\zeta_{2p}}{2}\, X(J) + A_{SS}\, f_{SS}(J)
              + A_{SOO}\, f_{SOO}(J)
\end{equation}
with
$X(J) = J(J+1) - L(L+1) - S(S+1)$,
$f_{SS}(J) = (-2, +1, -\tfrac{1}{5})$ from the rank-2 6j symbol
$(-1)^{L+S+J}\cdot 6j\{1,1,J;1,1,2\}$ at $L = S = 1$ (Sprint 4 DD,
sympy-exact),
$f_{SOO}(J) = (+2, +1, -1)$ from the rank-1 6j,
and Drake combining coefficients
$A_{SS} = \alpha^2 (\tfrac{3}{50}\, M^2_\text{dir} - \tfrac{2}{5}\, M^2_\text{exch})$,
$A_{SOO} = \alpha^2 (\tfrac{3}{2}\, M^1_\text{dir} - M^1_\text{exch})$
(Sprint 3 BF-D, intrinsic-tier rationals).

\begin{table}[h]
\centering\small
\begin{tabular}{p{2.6cm} r r r r p{4.5cm} c}
\toprule
Interval & $E_{SO}$ (MHz) & $E_{SS}$ (MHz) & $E_{SOO}$ (MHz) & Total (MHz) & Projection chain & Status \\
\midrule
$P_0 - P_1$ & $-29\,198.090$ & $+23\,747.770$ & $+35\,063.226$ & $+29\,612.906$ &
\S\ref{sec:proj_fock} $\circ$ \S\ref{sec:proj_spinor} $\circ$
\S\ref{sec:proj_sturmian} (SO); $+$ (SS, SOO chains via
\S\ref{sec:proj_tensor_multipole} $\circ$ \S\ref{sec:proj_3j}) & FN \\
$P_1 - P_2$ & $-58\,396.180$ & $-9\,499.108$ & $+70\,126.452$ & $+2\,231.165$ &
(same as $P_0$-$P_1$) & FN \\
$P_0 - P_2$ & $-87\,594.270$ & $+14\,248.662$ & $+105\,189.679$ & $+31\,844.070$ &
(same as $P_0$-$P_1$) & FN \\
\midrule
\textbf{NIST} & & & & & & \\
$P_0 - P_1$ & & & & $+29\,616.951(6)$ & & \\
$P_1 - P_2$ & & & & $+2\,291.176(15)$ & & \\
$P_0 - P_2$ & & & & $+31\,908.131(6)$ & & \\
\midrule
\textbf{Residual} & & & & & & \\
$P_0 - P_1$ & & & & $-4.05$ MHz / $-0.014\%$ & & \\
$P_1 - P_2$ & & & & $-60.01$ MHz / $-2.62\%$ & partial cancellation $61.9\times$ & \\
$P_0 - P_2$ & & & & $-64.06$ MHz / $-0.20\%$ & & \\
\bottomrule
\end{tabular}
\caption{He $2{}^3P$ fine-structure operator-level Roothaan autopsy.
The angular content (Drake $J$-pattern, spin reduced m.e., bipolar
Gaunt selection) is sympy-exact at machine precision; residuals live
entirely in the radial sector (multi-electron correlation, higher-order
Breit-Pauli, $\alpha^3/\pi$ QED, recoil).  Source memo:
\texttt{debug/he\_2\_3P\_autopsy\_v1\_memo.md}.}
\label{tab:autopsy_he_2_3p}
\end{table}

\textbf{Bipolar (k$_1$, k$_2$) decomposition.}
For the (1s)(2p) configuration with $L_\text{max} = 2 l_\text{max} = 2$:
SS (K$=$2) admits direct (k$_1=0$, k$_2=2$) and exchange (k$_1=1$, k$_2=1$)
channels only; SOO (K$=$1) admits exchange (k$_1=1$, k$_2=1$) only (the
direct path is empty by parity selection).  All higher (k$_1$, k$_2$)
channels are Gaunt-forbidden.  This is the INTERNAL analog of the
cross-register Roothaan termination (B-W1a-diag finding): the
multipole termination at $L_\text{max} = 2 l_\text{max}$ holds
\emph{independent of the focal-length ratio} $Z_\text{eff}(1s) /
Z_\text{eff}(2p) = 2$ --- the channel structure is angular content,
not radial content.

\textbf{Partial-cancellation mechanism ($P_1$-$P_2$ amplification).}
The small interval is the difference SO ($-58.4$ GHz) $+$ SS
($-9.5$ GHz) $+$ SOO ($+70.1$ GHz) $= +2.2$ GHz, a structural
cancellation by factor $61.9\times$.  The framework's $60$ MHz
absolute residual on $P_1 - P_2$ is the same order as the $64$ MHz
residual on $P_0 - P_2$; the $-2.62\%$ fractional residual is the
$60$ MHz absolute divided by the $14\times$-smaller small interval.
Generic feature of multi-component fine-structure splittings (Li
$2{}^2P$ and Be $2s2p\,{}^3P$ show analogous partial-cancellation
amplification, Sprint 5 CP).

\textbf{Three-class diagnosis (per CLAUDE.md \S~1.8 directive).}
Class A (literature convention mismatch): NEGATIVE --- Drake (1971)
and Pachucki--Yerokhin (2010) agree on leading-order Breit-Pauli
itemization.  Class B (framework gap, RADIAL): POSITIVE --- residual
attributes to (i) multi-electron correlation in the (1s)(2p)
configuration (largest source, $\sim 45$ MHz on $P_0$-$P_2$ scale,
not in framework's single-configuration hydrogenic basis), (ii)
$\alpha^3(Z\alpha)^2$ QED on Breit-Pauli (LS-8a vertex sector,
$\sim 15$ MHz), (iii) $\alpha^3$ recoil (W1a cross-register
kinematic wall, $\sim 2$ MHz), (iv) higher-order $\alpha^4 + \alpha^5
+ \alpha^6$ NRQED (LS-8a multi-loop, sub-MHz).  Class C (focal-length
cataloguing): POSITIVE --- partial-cancellation amplification is
generic across atomic fine-structure splittings.

\textbf{Structural reading.}  The autopsy realizes the
structural-skeleton-scope statement at the operator level for an
internal multi-focal observable.  Five projection chains, three
operators, sympy-exact angular content, residuals in the radial
sector --- the same architecture that closes hydrogen Lamb shift at
\S\ref{sec:autopsy_lamb} closes He $2{}^3P$ fine structure here.
The internal multi-focal angular-only prediction (Observation 3)
holds at machine precision: angular Roothaan termination is
independent of focal-length ratio.

\subsubsection{Helium $2{}^1\!P \to 1{}^1\!S$ oscillator strength}
\label{sec:autopsy_he_oscillator}

\textbf{Reference.}  $f_\text{length} = 0.27616$ (Drake handbook;
Theodosiou 1984), oscillator strength for the He $2{}^1P \to 1{}^1S$
electric-dipole transition; $\omega_\text{exact} = 0.7799$ Ha.

The first multi-electron extension of Sprint Calc-L (hydrogen Lyman~$\alpha$,
$+0.055\%$ at depth-3) and Sprint Calc-P (hydrogen polarizability,
machine-exact at $N_{\rm basis}=2$ in the Sturmian basis).  The
structural question: does the Sturmian continuum-closing property
extend to multi-electron transition matrix elements?

The framework-native length-form oscillator strength is computed via
\[
f = 2\omega\, |\langle 1{}^1\!S | z | 2{}^1\!P_{m=0}\rangle|^2
\]
in the graph-native CI sector at fixed $M_L$, with the $1{}^1\!S$ state in
the $(l_1=l_2=0)$ ss singlet sub-block, the $2{}^1\!P$ state in the
$(l_1=0, l_2=1, M_L=0)$ sp singlet sub-block, and the transition dipole
assembled by Slater--Condon expansion of the multi-electron CI vectors.

\begin{table}[h]
\centering\small
\begin{tabular}{c c c c c c}
\toprule
$n_{\max}$ & configs & $\omega$ (Ha) & $|\langle z\rangle|^2$ & $f_\text{length}$ & err vs Drake \\
\midrule
2 & 5 & 0.9044 & 1.107 & 2.002 & $+625\%$ \\
3 & 12 & 0.7045 & 0.420 & 0.592 & $+114\%$ \\
4 & 22 & 0.6930 & 0.397 & 0.550 & $+99\%$ \\
5 & 35 & 0.6891 & 0.323 & 0.446 & $+61\%$ \\
6 & 51 & 0.6893 & 0.334 & 0.460 & $+67\%$ \\
7 & 70 & 0.6893 & 0.322 & 0.444 & $+61\%$ \\
\midrule
exact & --- & 0.7799 & 0.177 & 0.27616 & 0 \\
\bottomrule
\end{tabular}
\caption{Coulomb Sturmian convergence for He $2{}^1P \to 1{}^1S$
oscillator strength at single exponent $k_\text{orb} = Z = 2$.  Both
$\omega$ and $|\langle z\rangle|^2$ plateau by $n_{\max} = 6$ at
\emph{wrong limits} ($\omega$ low by $-12\%$, $|\langle z\rangle|^2$
high by $+82\%$).  Source memo:
\texttt{debug/he\_oscillator\_v1\_memo.md}.}
\label{tab:autopsy_he_oscillator}
\end{table}

Both $\omega$ and $|\langle z\rangle|^2$ \emph{plateau} by $n_{\max} = 6$ at
\emph{wrong limits}: $\omega$ converges $-12\%$ low and $|\langle z\rangle|^2$
converges $+82\%$ high.  The $2{}^1\!P$ state is variationally violated by
$\sim 0.08$ Ha ($E(2{}^1\!P)_\text{GeoVac} = -2.204$ vs Drake $-2.124$ Ha) ---
the same $\kappa$-induced over-binding mechanism documented for
the He $2{}^1\!S - 2{}^3\!S$ splitting (small-$Z$
graph-validity-boundary artifact, $Z_c \approx 1.84$; He at $Z = 2$
is just above; cf.\ \S\ref{sec:offprecision} He singlet-triplet row).
The angular (Wigner $3j$) and vector-photon
(\S\ref{sec:proj_vector_photon}) projections are exact; the failure
is in the CI sector.

\textbf{Structural reading.}  Sturmian's continuum-closing is a
\emph{single-focal-length} property.  For multi-electron transitions where
initial and final states have structurally different effective focal lengths
(here, doubly-screened $1s$ and singly-screened $2p$), the single-exponent
Sturmian basis cannot represent both, and the framework's CI eigenvectors
overlap incorrectly even at large $n_{\max}$.  The natural extension is the
\emph{multi-focal} architecture (\verb|geovac/cross_register_vne.py| for
chemistry; internal split-$Z$ for atomic excited states is a named
follow-on track per Sprint Calc-He recommendation).

The $+61\%$ residual at $n_{\max} = 7$ does not improve at higher $n_{\max}$
within the single-exponent architecture and is therefore tagged error code
\texttt{B} (basis architecture) in
Table~\ref{tab:catalog_off}.  This entry is a \emph{structural} negative
result that complements the He $2{}^1\!S - 2{}^3\!S$ off-precision row:
both expose the small-$Z$ graph-validity-boundary at He as the dominant
limitation for excited-state precision, separable from the framework's
clean closures on hydrogenic and ground-state observables.

\subsubsection{Cesium 6S$_{1/2}$ hyperfine (atomic clock; SCOPING)}
\label{sec:autopsy_cs_hfs}

\textbf{Reference.}  $\Delta\nu = 9\,192\,631\,770$ Hz (BIPM 1967,
SI-second-defining transition), equivalently
$A(6S_{1/2}) = \Delta\nu/4 = 2298.158$ MHz for the
Cs-133 $6S_{1/2}$, $F=4 \leftrightarrow F=3$ hyperfine transition.

A scoping pass (Track Cs-HFS-v1, May 2026) confirms the framework's
projection-chain machinery (\texttt{atomic\_classifier}, [Xe]
\texttt{FrozenCore}, Bohr--Fermi formula) extends to $Z=55$ without
structural failure: the prediction lives in $A \in [793, 2931]$ MHz
across the natural range of effective screening charges
$Z_\text{eff} \in [9.7, 15]$, bracketing the experimental value at
the inverse-solve $Z_\text{eff} \approx 10$--$11$ (consistent with
Roberts--Ginges 2019).

A clean framework-native compute is blocked by two engineering items:
\begin{itemize}\setlength\itemsep{0pt}
\item The screened radial solver \texttt{\_solve\_screened\_radial}
rejects $\ell=0$ by design, requiring extension to log-grid $+$
Numerov $+$ analytical small-$r$ fit for the $s$-wave $|\psi(0)|^2$.
\item The Track NI hyperfine Pauli wrapper
\texttt{hyperfine\_coupling\_pauli} is specialized to the 1p$+$1e
architecture, requiring a generic $A \cdot \vec{I}\cdot \vec{J}$
wrapper for atomic HFS observables.
\end{itemize}
Both are mechanical (~2--3 week sprint).

\textbf{Z-scaling assessment.}  The heavy-atom regime does not expose
structural limitations beyond those already identified at $Z=20$--$56$
in Sprint 7b screened SO splittings (Clementi--Raimondi single-zeta
screening qualitative; W1c-residual core-polarization correction
needed for sub-percent precision).

\textbf{Diagnostic-instrument angle.}  The framework's projection-chain
decomposition would isolate three known atomic-structure
dependencies of the percent-level Cs PNC uncertainty:
\begin{itemize}\setlength\itemsep{0pt}
\item radial wavefunction at the origin (Bohr--Weisskopf via
\S\ref{sec:proj_magnetization_density}),
\item core polarization (relativistic enhancement via
\S\ref{sec:proj_spinor}),
\item multi-loop QED in the heavy-atom regime (two-body Dirac/Breit
Layer-2 input via \S\ref{sec:proj_breit_retardation}),
\end{itemize}
each tied to a different Layer-2 input.  This is the prospective
catalogue's first heavy-atom test of the framework's projection-chain
machinery in a regime where simple hydrogenic formulas break down,
and connects to atomic-clock-grade precision ($10^{-15}$ relative,
the SI-second definition).  Source memo:
\texttt{debug/cs\_hfs\_v1\_memo.md}.

% -------------------------------------------------------------------
\subsection{Literature convention exposures (running catalogue)}
\label{sec:convention_exposures}
% -------------------------------------------------------------------

The framework's projection-chain decomposition acts as a unique
diagnostic instrument for precision atomic and molecular physics:
multi-observable global fits using a single structural dictionary
(the \S~\ref{sec:layer2} projection family) expose itemization
choices in standard-physics compilations that single-observable
analyses cannot see.  This is the directive's class~(a) finding
type --- ``literature convention mismatches'' between independent
compilations (Eides~2024, Karshenboim~2005, Krauth~2017,
Pachucki--Yerokhin~2010, Antognini~2013, Friar--Payne~2005)
on Layer-2 itemization at the sub-percent precision the framework
operates at.  Frameworks that consume single-observable Layer-2
inputs in their native conventions cannot expose these mismatches
structurally.  This subsection consolidates the four exposures
catalogued to date (2026-05-09); each future sprint that surfaces a
class~(a) finding adds a new \S~V.D.$N$ subsection with the
five-element template \emph{Observable~$|$~Source compilations~$|$~
Convention difference~$|$~Magnitude~$|$~Sprint reference $+$ Closure verdict.}

\begin{table}[h]
\centering\small
\caption{Literature convention exposures catalogued through 2026-05-09.
Each row is documented in detail in the indicated subsection
(originating sprint memo) and cross-references the
\S~\ref{sec:layer2} projection chain whose operator-level
decomposition exposed it.}
\label{tab:convention_exposures}
\begin{tabular}{p{0.6cm} p{2.6cm} p{3.0cm} p{2.6cm} p{1.6cm} p{1.6cm}}
\toprule
ID & Observable & Source compilations & Convention difference &
Magnitude & Sprint ref. \\
\midrule
\S\ref{sec:conv_rZG_DHFS} & D 1S HFS Layer-2 budget &
Eides~2024 Tab.~7.3 vs Pachucki--Yerokhin~2010 &
itemization of recoil NLO + multi-loop QED + finite-size correction &
$\sim 0.01\%$ of $\nu_F$ $\to$ $\sim 25$ mfm in $r_Z(p)$ &
Calc-rZG (\S\ref{sec:offprecision}) \\
\S\ref{sec:conv_muH_VP} & $\mu$H 2$S$-2$P$ VP one-loop &
Antognini~2013 vs Krauth~2017 &
absorption of higher-order VP into one-loop line vs separate
``VP-VP mixed'' itemization &
$\sim 100$ ppm in per-component split (totals agree) &
Sprint MH Track A / Calc-muH Lamb autopsy v1 \\
\S\ref{sec:conv_karshenboim_recoil} & $\mu$ Lamb shift SE bracket &
Framework Eides~\S 3.2 vs Karshenboim~2005 &
$m_\text{red}$ rescaling absorbs recoil NLO at bracket level vs
separately-itemized ``Recoil NLO'' line &
$\sim 11$ MHz at $\mu$ Lamb (exact match between framework
$+1074.13$ and Karshenboim $+1085.84 - 11.30$) &
Sprint precision-catalogue Mu Lamb 2026-05-08 \\
\S\ref{sec:conv_friar_profile} & $\mu$H Friar moment &
Eides Eq.~2.35 closed-form vs operator-level
\S\ref{sec:proj_magnetization_density} &
$\langle r^2\rangle^{1/2} = r_p$ (RMS) vs $\langle r\rangle = r_p$
(first moment) for the Gaussian profile &
$\langle r^2\rangle/\langle r\rangle^2 = 4/\pi \approx 1.27$
$\to$ $+12.7\%$ on Friar moment &
Sprint MH / Calc-muH Lamb autopsy v1 \\
\bottomrule
\end{tabular}
\end{table}

\subsubsection{D 1S HFS Layer-2 itemization (Eides vs Pachucki--Yerokhin)}
\label{sec:conv_rZG_DHFS}

\textbf{Observable.}  Deuterium 1S hyperfine splitting
$\nu_\text{HFS}(D) = 327\,384\,352.522$ kHz (Wineland--Ramsey 1972),
used in the Sprint Calc-rZG global fit (\S\ref{sec:offprecision}, row
``Self-consistent $r_Z$ global extraction'') to extract
$r_Z(D)$ via the framework's \S\ref{sec:proj_magnetization_density}
operator.

\textbf{Source compilations.}  Eides~2024 Tab.~7.3 (used as the
canonical Layer-2 budget for H 21cm in this paper) and
Pachucki--Yerokhin~2010 (the standard atomic-physics itemization for
deuterium HFS).

\textbf{Convention difference.}  The two compilations use different
groupings for the recoil-NLO + multi-loop QED + finite-size
contributions to $\nu_F^D$.  Eides Tab.~7.3 itemizes the leading
Zemach line at $-150$~ppm; Pachucki--Yerokhin~2010, when fully
expanded with the Friar--Payne~2005 closure analysis, places the
itemization at $-286$~ppm.  Both compilations close to the same
total $\nu_\text{HFS}(D)$ at the precision the deuterium HFS
measurement reports; the difference is in how sub-leading recoil and
multi-loop contributions are aggregated.

\textbf{Magnitude.}  $\sim 0.01\%$ of $\nu_F$, propagating to
$\sim 25$~mfm in extracted $r_Z(p)$ via the global fit.

\textbf{Sprint reference and closure.}  The original Sprint Calc-rZG
fit (2026-05-09 morning) used the Eides $-150$~ppm itemization for
deuterium; this gave $r_Z(D) = 6.18(132)$~fm, the well-known
$\sim 5$~fm artifact.  The diagnostic memo
\texttt{debug/rzg\_bug\_diagnosis\_memo.md} verified that the
production cross-register $V_{eN}$ kernel
(\texttt{geovac.cross\_register\_vne}) reproduces leading-order
Bethe--Salpeter recoil at $2.86\%$ for hydrogen, $2.03\%$ for
deuterium, and $8.18\%$ for muonium, species-agnostically and
correctly.  The artifact came from the Layer-2 budget specification
for D HFS, not from the framework kernel.  Re-anchoring to the
Pachucki--Yerokhin~2010 budget at $-286$~ppm closed the extraction
to $r_Z(D) = 2.583(1588)$~fm (within $0.01\sigma$ of Friar--Payne
$2.593(16)$~fm).  \textbf{Closure verdict:} resolved by Layer-2
budget re-anchoring.  The original ``recoil double-counting''
diagnosis was a misdiagnosis; the issue was Layer-2 itemization
specification, not a kernel bug.

\textbf{Cross-references.}  Originating documentation:
\S\ref{sec:offprecision} Sprint Calc-rZG row.  Source memos:
\texttt{debug/calc\_track\_rZG\_global\_zemach\_memo.md},
\texttt{debug/rzg\_bug\_diagnosis\_memo.md}.

\subsubsection{$\mu$H 2$S_{1/2}$--2$P_{1/2}$ VP one-loop split (Antognini vs Krauth)}
\label{sec:conv_muH_VP}

\textbf{Observable.}  Muonic hydrogen Lamb shift VP one-loop
component, the dominant ($\sim +205$~meV out of $+202.37$~meV total)
contribution to the CREMA 2010 measurement.  Documented in the
\S\ref{sec:autopsy_muh_lamb} Roothaan autopsy.

\textbf{Source compilations.}  Antognini~2013 (the canonical
muonic-hydrogen reference for CREMA Lamb-shift theory; used as the
Layer-2 anchor in \S\ref{sec:autopsy_muh_lamb}) and Krauth~2017 (the
canonical muonic-hydrogen reference for the parallel HFS measurement
and the most thorough recent itemization).

\textbf{Convention difference.}  Antognini reports VP one-loop at
$+205.0074$~meV with ``VP-VP mixed'' itemized separately at
$+0.150$~meV.  Krauth absorbs $\sim 0.021$~meV of higher-order VP
into the one-loop line and reports $+205.0282$~meV.  Both compilations
agree on the bottom-line $\mu$H Lamb total to printed precision; the
difference is in per-component itemization.

\textbf{Magnitude.}  $\sim 100$~ppm in per-component split (totals
agree at much tighter precision).

\textbf{Sprint reference and closure.}  Sprint MH Track A
(\S\ref{sec:offprecision}, row ``$\mu$H 2$S$--2$P$ Lamb shift'') and
Sprint Calc-muH Lamb autopsy v1 (\S\ref{sec:autopsy_muh_lamb}) both
operationally surfaced the convention difference: the framework's
full-Uehling kernel matches Antognini~2013 at $-0.10$~ppm, which is
$1000\times$ tighter than the Antognini-vs-Krauth $\sim 100$~ppm
itemization gap.  This means the framework's match precision is
sensitive to the choice of reference compilation for Layer-2 inputs.
\textbf{Closure verdict:} anchored to Antognini~2013 throughout the
$\mu$H Lamb autopsy.  No reconciliation against Krauth attempted at
this time; flagged for future sprint if Krauth-anchored cross-checks
become relevant.

\textbf{Cross-references.}  Originating documentation:
\S\ref{sec:autopsy_muh_lamb} Convention-mismatch surface paragraph.
Source memo: \texttt{debug/muH\_Lamb\_autopsy\_v1\_memo.md} (LCM tag).

\subsubsection{$\mu$ Lamb shift SE-recoil aggregation (framework vs Karshenboim)}
\label{sec:conv_karshenboim_recoil}

\textbf{Observable.}  Muonium 2$S_{1/2}$-2$P_{1/2}$ Lamb shift
self-energy bracket.  Documented in the \S\ref{sec:offprecision} row
``Mu 2$S_{1/2}$-2$P_{1/2}$ Lamb $1047.63$~MHz vs Mariam~1982
experiment $1054(22)$~MHz.''

\textbf{Source compilations.}  Framework Eides~\S 3.2 SE bracket
(used in this paper's $\mu$ Lamb-shift framework calculation) and
Karshenboim~2005 (the canonical muonium-Lamb reference).

\textbf{Convention difference.}  The framework's $m_\text{red}$
rescaling at the Eides bracket level absorbs the leading
recoil-correction beyond reduced-mass automatically, giving
SE bracket $= +1074.13$~MHz.  Karshenboim itemizes ``SE one-loop''
at $+1085.84$~MHz with a separately-itemized ``Recoil NLO''
line at $-11.30$~MHz.  The two reach exact agreement at the bracket
level: $+1085.84 - 11.30 = +1074.54$~MHz vs framework $+1074.13$~MHz
(within bracket-evaluation precision).

\textbf{Magnitude.}  $\sim 11$~MHz on the $\mu$ Lamb-shift SE bracket
($\sim 1\%$ of the SE component, $\sim 0.01$ of the total Lamb shift);
the difference is exactly the recoil-NLO itemization line.

\textbf{Sprint reference and closure.}  Sprint precision-catalogue
2026-05-08 (Mu Lamb shift row, \S\ref{sec:offprecision}) operationally
surfaced this: the framework total at $1047.63$~MHz matches
Karshenboim~2005 theory at $1047.49(5)$~MHz to $+0.013\%$.  Naively
adding Karshenboim's separately-itemized ``Recoil NLO $-11.30$~MHz''
to the framework total double-counts the same recoil correction that
the framework's $m_\text{red}$ rescaling already absorbs.
\textbf{Closure verdict:} resolved by recognizing the framework's
$m_\text{red}$-bracket convention is structurally equivalent to
Karshenboim's SE one-loop $+$ Recoil NLO sum, and not adding the
recoil NLO twice.

\textbf{Cross-references.}  Originating documentation:
\S\ref{sec:offprecision} Mu 2$S_{1/2}$-2$P_{1/2}$ Lamb row.

\subsubsection{Friar moment profile convention (Eides closed-form vs operator-level)}
\label{sec:conv_friar_profile}

\textbf{Observable.}  Muonic hydrogen Friar moment, the
$O((Z\alpha)^2)$ finite-nuclear-size correction in the
\S\ref{sec:autopsy_muh_lamb} Lamb-shift autopsy (Component~3).

\textbf{Source compilations.}  Eides Eq.~2.35 closed-form derivation
(used as the Antognini~2013 reference value $-3.84$~meV at $r_p =
0.84087$~fm) and the framework's operator-level realization via the
\S\ref{sec:proj_magnetization_density} module
\texttt{geovac.magnetization\_density}.

\textbf{Convention difference.}  Eides Eq.~2.35 is derived assuming
the input radius is $\langle r^2\rangle^{1/2}$ (RMS charge radius);
the operator-level
\S\ref{sec:proj_magnetization_density} module is calibrated to the
first moment $\langle r\rangle = r_p$ (consistent with its
existing HFS Zemach use case).  For a Gaussian profile, the two
calibrations differ by exactly the factor
$\langle r^2\rangle / \langle r\rangle^2 = 4/\pi \approx 1.27$.

\textbf{Magnitude.}  $+12.7\%$ overshoot on the Friar moment
(operator-level $-4.33$~meV vs Eides closed-form $-3.68$~meV at
fixed $r_p$).  This is exactly $1.27\times$ the closed-form value to
$\sim 0.5\%$ relative precision.

\textbf{Sprint reference and closure.}  Sprint Calc-muH Lamb autopsy v1
(2026-05-08, \S\ref{sec:autopsy_muh_lamb} Operator-level Friar moment
paragraph) surfaced the discrepancy by computing the Friar moment
both ways: closed-form Eides reproduction matches Antognini~2013 at
$-4.3\%$, while operator-level realization overshoots Antognini by
$+12.7\%$, with the ratio between the two paths equal to the Gaussian
profile factor $4/\pi$ to sub-percent precision.
\textbf{Closure verdict:} flagged for a sibling
\texttt{MagnetizationDensitySpec} calibrated to RMS radius for the
Lamb-shift FNS channel, distinct from the existing first-moment
calibration for the HFS Zemach channel.  The HFS Zemach channel is
unaffected (cf. \S\ref{sec:autopsy_21cm}, where the operator-level
Zemach result $-39.495$~ppm matches Eides $-39.500$~ppm at
$0.012\%$).  Resolution is mechanical (~1 day sprint).

\textbf{Cross-references.}  Originating documentation:
\S\ref{sec:autopsy_muh_lamb} Operator-level Friar moment paragraph.
Source memo: \texttt{debug/muH\_Lamb\_autopsy\_v1\_memo.md}
(profile-convention finding).

\paragraph{Cross-pattern observation.}
The four exposures partition cleanly into two structural classes.
\emph{Class~(i) inter-compilation itemization} (entries
\S\ref{sec:conv_rZG_DHFS}, \S\ref{sec:conv_muH_VP},
\S\ref{sec:conv_karshenboim_recoil}): three of the four exposures
are between Eides~2024, Krauth~2017, Karshenboim~2005,
Pachucki--Yerokhin~2010, and Antognini~2013 compilations on items at
sub-percent precision.  The pattern is consistent: most precision-physics
compilations agree on bottom-line totals but differ in itemization at
the $10^{-4}$ to $10^{-3}$ level.  This is exactly the precision range
the framework's projection-chain decomposition operates at, hence the
surface contact between the framework's diagnostic instrument and
the literature compilations.  \emph{Class~(ii) closed-form vs
operator-level convention} (entry \S\ref{sec:conv_friar_profile}):
one exposure is between Eides closed-form derivation (which assumes
RMS-radius input) and the framework's operator-level realization
(which calibrates to first-moment input), with the two paths
differing by a profile-dependent factor that the closed-form
derivation hides.  Class~(ii) is internal to the framework's
operator-level realization choices, not a literature convention
mismatch in the strict sense, but it is the same kind of finding:
sub-percent agreement on totals masks per-channel calibration choices
that surface only when the two paths are compared at operator level.
Both classes vindicate the directive's class~(a) finding type:
multi-observable global fits using a single dictionary expose
itemization choices that single-observable analyses cannot see.

\paragraph{Protocol for adding new exposures.}
Each future sprint that surfaces a class~(a) finding adds a new
\S~V.D.$N$ subsection following the five-element template:
\emph{Observable~$|$~Source compilations~$|$~Convention difference~$|$
~Magnitude~$|$~Sprint reference $+$ Closure verdict.}  The exposure
is also added as a row to Table~\ref{tab:convention_exposures}.
Each entry should cross-reference the \S~\ref{sec:layer2} projection
chain whose operator-level decomposition surfaced the convention
difference, and the originating documentation in the catalogue
(\S\ref{sec:catalog}, \S\ref{sec:offprecision}, or \S\ref{sec:autopsies}).
The catalogue is a running register; expect it to grow as the
multi-observable focal-length decomposition program (CLAUDE.md
\S~1.8) advances.

% -------------------------------------------------------------------
\subsection{Falsified candidate: W3 spectral-zeta identification of
SM calibration data (sprint 2026-05-08, falsified same evening)}
\label{sec:calibration_candidates}
% -------------------------------------------------------------------

This subsection preserves the record of a single-session candidate
that was \emph{falsified} the same evening by three independent
follow-up tracks.  It is retained in the catalogue because the
methodological lesson -- how a $6$--$10\sigma$ signal under a
hand-curated prespecified basis can vanish under a mechanically
generated basis -- is a sharper instance of the curve-fit-audit
discipline (see \texttt{docs/curve\_fit\_audit\_memo.md}) than any
abstract description.

\paragraph{The original candidate.}
On 2026-05-08 evening, a conversational sprint produced candidate
spectral-zeta identifications of all four CKM Wolfenstein parameters
with master-Mellin-engine M1 and M2 ring elements:
$\lambda^2 \approx 1/\mathrm{Vol}(S^3) = 1/(2\pi^2)$ and
$\bar\rho \approx 1/(2\pi)$ in the M1 Hopf-base family;
$\bar\eta \approx (\ln 2)/2 = -\zeta'(0,1/2)$ and
$A^2 \approx \ln 2$ in the M2 chirality family.  The four-form fit
gave 7 hits at $<1\%$ in a hand-curated 30-form prespecified basis
(random expectation $0.7\pm 1.0$); a derived structural relation
$A^2 = 2\bar\eta$ survived PDG at $0.52\sigma$; a closed form
$\tan(\delta_{CP}) = \pi\cdot\ln(2)$ at $0.02\sigma$ provided a
striking single-number target.  The aggregate signal was reported as
approximately $6$--$10\sigma$ above chance for that basis size.  The
candidate motivated a draft addition to CLAUDE.md~\S 1.7 (WH7) and
prompted a forward plan in
\texttt{debug/w3\_forward\_plan\_memo.md} listing six tests.

\paragraph{Three-track falsification (same evening).}
The first three tests of the forward plan ran in parallel as
independent sub-agent dispatches.  All three returned negatives.

\begin{itemize}
\setlength\itemsep{2pt}
\item \textbf{PMNS recheck.}  The same 30-form prespecified basis
applied to the PMNS angles
$(\sin^2\theta_{12}, \sin^2\theta_{23}, \sin^2\theta_{13},
\delta_{CP}^{\rm PMNS})$ at PDG~2024 NH best-fit values returned
\emph{zero} matches at $<2\%$ across all four parameters.  Aggregate
signal: $-0.68\sigma$ below chance.  PMNS sits closer to tri-bimaximal
pure rationals $(1/3, 1/2, 0)$ than to any M1/M2 form.  The cleanest
single PMNS hit is $\theta_{12}\approx \arcsin(1/\sqrt{3})$ at
$0.11\,\sigma$ -- the TBM solar-angle prediction, not in the master
Mellin rings.  W3 universality across mixing matrices is falsified.
Files: \texttt{debug/w3\_pmns\_recheck.\{py,\_memo.md\}}.

\item \textbf{Lepton mass spectrum.}  The same methodology applied to
$\sim$30 charged-lepton mass observables (ratios, log spacings,
sqrt-mass ratios, Koide $K$, cone position) crossed with the 30-form
basis and 10 small-rational factor multipliers (9000 tests total)
returned only $2$ within-$1\sigma$ matches, both encoding the same
Koide $45^\circ$ cone fact via algebraically equivalent identities
($2/3 \times (\varphi^2 - \varphi) = 2/3$ and
$4 \times \pi/16 = \pi/4$).  Statistical signal: $+0.85\sigma$ above
chance, indistinguishable from random.  Hit-density gap CKM~vs~lepton:
$12\%$ vs $0.22\%$ ($50\times$).  W3 identification of mass spectra
is falsified.  Files:
\texttt{debug/w3\_lepton\_mass\_recheck.\{py,\_memo.md\}}.

\item \textbf{Mechanically-generated basis sweep.}  The hand-curated
30-form basis was replaced by a mechanically generated 21,448-form
basis: combinations of M1/M2/M3 seeds (Hopf-base measures, Hurwitz
shifts, Catalan~$G$, Dirichlet~$\beta$, $\zeta_R$ at integer points)
with small-integer rational prefactors and small algebraic factors
($\sqrt{2},\sqrt{3},\sqrt{5},\varphi$), plus a pure-rational control
class.  The full basis was frozen and saved to JSON before any test
run.  Re-running the four-Wolfenstein search returned aggregate
$z = -0.52$ (slightly \emph{below} random chance); per-parameter
z-scores all within $\pm 1\sigma$ of chance ($\lambda$ $+0.74$,
$A$ $-0.02$, $\bar\rho$ $-1.01$, $\bar\eta$ $+0.82$).  The original
M1$\leftrightarrow$real / M2$\leftrightarrow$CP-imaginary structural
reading does not survive: $\bar\eta$'s top mechanical match is in the
M1+algebraic class ($\pi^{3/2}/16$), not M2; $A$'s original
$\sqrt{\ln 2}$ does not appear in the top-20 sigma-ordered slice; and
$\lambda$'s top match is the pure rational $9/40 = 0.225$ exactly,
beating every transcendental form.  Selection-bias verdict: fully
attributable.  Files:
\texttt{debug/w3\_mechanical\_basis.\{py,\_memo.md\}}.
\end{itemize}

\paragraph{Three independent tests, three negatives.}
The W3 spectral-zeta candidate is falsified.  The original signal was
basis-selection bias.  WH7 is not promoted to CLAUDE.md~\S 1.7.

\paragraph{Methodological lesson.}
The original ``$5$--$10\sigma$ above chance'' calculation was
statistically correct \emph{for that basis size}.  The error was the
implicit assumption that the 30 hand-curated forms were drawn without
knowledge of the data; in fact, the basis was disproportionately near
the empirical parameter values.  Mechanical generation removes the
implicit selection, and when removed, the signal vanishes.  This is
the cleanest illustration the project has produced of why the
curve-fit-audit discipline (independent re-run with strictly
mechanically-generated basis) is necessary, not optional, for
empirical pattern claims that touch transcendental rings of any
size.

\paragraph{Surviving observations independent of W3.}
The charged-lepton Koide $45^\circ$ cone (1 arcsec) is independent of
the W3 sprint and stands.  One isolated match is preserved as a
record (no structural reading attached): the CKM Wolfenstein
parameter $\lambda = 0.225$ is exactly $9/40$ to $5$-digit PDG
precision, which in GeoVac vocabulary reads as $\lambda = 9\Delta$
with $\Delta^{-1} = g_3^{\rm Dirac} = 40$ (Paper~2 Dirac-mode
boundary count); whether this is structural or Diophantine
coincidence at PDG precision is not addressable from a single data
point and the framework does not commit to either reading.

\paragraph{Standing position on calibration data.}
GeoVac's structural-skeleton scope (Bertrand truncation, inner-factor
Mellin trivialization, H1/G3 verdicts, LS-8a renormalization gap)
is restored as the standing position on calibration data: parameter
values that are not autonomously generated by the framework are
genuinely external input.  The second-packing-axiom question for
calibration data remains open at the same level it was on
2026-05-07.

% ===================================================================
\section{Dimensional Axiomatization}
\label{sec:axiomatization}
% ===================================================================

The projection dictionary of \S~\ref{sec:dictionary} treats each
projection as a graded map along three axes: variables added,
dimensions added, and transcendental class introduced.  Sprint TX-A
(May~2026) treated this graded structure as an axiomatic system and
tested three theorems analytically.  This section consolidates the
verdicts; full data are in
\texttt{debug/data/tx\_a\_projection\_table.json} and
\texttt{debug/data/tx\_a\_composition\_table.json}, with verification
driver \texttt{debug/tx\_a\_signature\_arithmetic.py}.

\begin{theorem}[Dimensional Consistency]
\label{thm:dimensional_consistency}
For any projection chain $S = P_n \circ \cdots \circ P_1$ applied to
the bare graph (Layer~1, dimensionless), the variable, dimension, and
transcendental signatures of $S$ equal the union of the signatures of
the constituent $P_i$:
\[
  \mathrm{sig}(P_n \circ \cdots \circ P_1)
   = \bigcup_{i=1}^{n} \mathrm{sig}(P_i).
\]
The catalogue rows of \S~\ref{sec:catalog} satisfy this consistency
check across all eight fully-tagged entries with non-trivial chains
(Lamb shift, $K = \pi(B + F - \Delta)$, Stefan-Boltzmann, He CI,
hydrogen $E_n$, scalar/Dirac spatial Casimir, Uehling VP).
\end{theorem}

The theorem is true by construction at the level of the dictionary
metadata: each projection's signature is fixed in
\S~\ref{sec:layer2}, so set/list union is a tautology.  Its
non-trivial content is the empirical check that the catalogue's
expected target signature matches the chain signature; this passes
for all eight tested rows.

\begin{theorem}[Weak Minimality]
\label{thm:minimality}
Two chains containing the same multiset of projections produce the
same signature, hence are observationally indistinguishable at the
level of dimensional types and transcendental classes.  They may
differ at the operator-system level when bundle-non-trivial
projections (Hopf, spectral action, Wilson) are composed with the
spinor lift; these inequivalences are not visible to signature
arithmetic and require explicit operator-system analysis (cf.\ the
WH1 spectral truncation rounds 1--3 documented in
\texttt{CLAUDE.md} \S~2).
\end{theorem}

The strong form of minimality (uniqueness up to commutation as
operators on the Hilbert space) fails for the Paper~2 $K$-chain:
\texttt{hopf} composed with \texttt{spinor\_lift} acts differently in
the two orderings (Hopf-then-spinor produces Hopf-$S^2$ with Dirac
fiber; spinor-then-Hopf gives a $\mathbb{Z}_2$-quotient of Dirac
$S^3$ with half-integer charges).  After the addition of
\texttt{breit\_retardation} (\S~\ref{sec:proj_breit_retardation}),
a second non-commuting pair appears:
\texttt{rest\_mass} composed with \texttt{breit\_retardation} is
ordered, with rest-mass first.  The non-commutation is semantic:
Breit retardation's variable $m_l/m_n$ is a refinement of the
rest-mass variable~$m$ (it is the ratio between two species' masses,
both of which are introduced by the rest-mass projection), so it is
not well-defined until rest-mass has acted.  The metadata-level
test correctly flags both
$(\texttt{hopf}, \texttt{spinor\_lift})$ and
$(\texttt{rest\_mass}, \texttt{breit\_retardation})$ as the
non-commuting pairs in the worked examples; the underlying
mechanisms differ
(bundle-non-trivial composition for the first, variable-refinement
dependence for the second).

\begin{theorem}[Composition Algebra]
\label{thm:composition_algebra}
The 16$\times$16 composition table classifies each ordered pair
$(P_1, P_2)$ as one of: \emph{commute} (admissible in both orders,
disjoint variables); \emph{one-way} (admissible in only one order,
constrained by Layer transitions or by semantic dependence between
variables); \emph{not composable}; \emph{idempotent} (Layer-1
self-composition); \emph{self-compose} (non-trivial diagonal);
\emph{variable conflict} (same variable name added by both);
\emph{inverse pair} (one is the inverse of the other).  Of
$16^2 = 256$ cells:
\begin{itemize}
\setlength\itemsep{0pt}
\item $\sim 196$ commute ($\sim 76.6\%$; updated from 15$\times$15
count of 182 by adding the 14 new commuting cells in P$_{16}$'s row
and column with projections that compose layer-cleanly under disjoint
variable axes);
\item $\sim 26$ one-way pairs (one-way cells flagged either by Layer
transitions or by semantic dependence between variables; updated by
$+4$ from the 15$\times$15 count of 22 to account for the four new
one-way cells flagged below in P$_{16}$'s row and column);
\item $6$ variable conflicts (energy-scale-conflating compositions
of Fock, Bargmann--Segal, and spectral action; the new P$_{16}$
introduces no further variable conflicts since its variable
$m_l/m_n$ is distinct from the existing variable set);
\item $4$ idempotent projections on the diagonal (Sturmian,
Wigner~$3j$, Drake--Swainson, and now two-body Dirac / Breit
retardation: applying the Breit Hamiltonian twice gives the same
two-body output);
\item $3$ inverse pairs (stereographic, Wigner~$3j$, Wigner~$D$;
unchanged by the new projection).
\end{itemize}
The previous version of this theorem flagged \texttt{drake\_swainson}
as the unique projection with input\_layer = output\_layer = 2.
After the addition of \texttt{breit\_retardation}
(\S~\ref{sec:proj_breit_retardation}), \emph{two} projections now
have input\_layer = output\_layer = 2: Drake--Swainson and Breit
retardation.  The structurally distinguishing fact is that they
operate on disjoint Layer-2 content (Drake--Swainson regularizes
divergent Bethe-log spectral sums; Breit retardation adds two-body
kinematic content to the framework's single-particle Lamb-shift
bracket), so they commute as a $(P_1, P_2)$ pair; their
non-uniqueness on the layer index is harmless for signature
arithmetic.

\textbf{New structural finding (the load-bearing T3 update from
this addition).}  The pair
$(P_{14}, P_{16}) = (\text{rest-mass projection},
\text{Breit retardation})$ does \emph{not} commute, but the
non-commutation is \emph{not} a Layer transition --- both projections
output to Layer~2 from input layers 1 and 2 respectively, and the
Layer indices alone would mark this as a one-way pair under the
existing classification.  The \emph{semantic} reason for one-way
ordering is sharper: the Breit projection's variable $m_l/m_n$ is
the rest-mass ratio between two species of particle in the bound
state, which is well-defined only after rest-mass projection has
assigned masses to each species.  In the reverse order
$P_{14} \circ P_{16}$ the variable $m_l/m_n$ would have to be
introduced before the masses themselves exist, which is semantically
inadmissible.  This is the first instance in the Paper~34 dictionary
of a one-way pair driven by a \emph{variable refinement} dependence
rather than by a Layer transition.  The signature-arithmetic test
correctly flags both the
$(P_{14}, P_{16})$ and the
$(P_{16}, P_{14})$ cells, and the test agrees with the structural
reading: Breit retardation must follow rest-mass projection in any
admissible chain.

\textit{Erratum (Track TS-C, 2026-05-04):} an earlier draft of this
section claimed that \texttt{spectral\_action} and
\texttt{observation\_window} were mutually one-way, encoding the
field-theoretic UV/IR ordering anomaly.  That claim was retracted
on direct inspection of the composition data: both projections
target output\_layer=2 from input\_layer=1 with disjoint variable
sets, and signature arithmetic correctly returns \texttt{commute}
in both directions.  The UV/IR anomaly ($\beta \to \infty$ does not
commute with $\Lambda \to \infty$) is real physics that the current
coarse layer-grading cannot see; refining the output\_layer index
to distinguish UV-cutoff, thermal, and flow Layer-2 outputs is the
natural next step.

\textit{Note on cell counts.}  The exact 16$\times$16 cell counts
are stated above as approximate ($\sim$) totals; the load-bearing
update is the 15$\times$15 $\to$ 16$\times$16 expansion plus the
qualitative shifts (new idempotent, new one-way pair driven by
variable refinement, no new variable conflicts).  Definitive
recomputation against the metadata in
\texttt{debug/data/tx\_a\_composition\_table.json} is deferred until
the next TX-A-style audit.

\textit{Pending T3 expansion (2026-05-09).}  With the addition of
the three intra-nuclear projections
(\S~\ref{sec:proj_charge_density},
\S~\ref{sec:proj_magnetization_density},
\S~\ref{sec:proj_tensor_multipole}), the T3 composition table needs
expansion from 16$\times$16 to 19$\times$19.  The qualitative
expectation: the three new projections all introduce intra-nuclear
$[L]$ (or $[L]^2$) variables disjoint from the existing length
axes, so they should commute layer-cleanly with most existing
projections; expected new entries are $\sim 50$ commute cells in
each new row and column, with possible variable-refinement one-way
pairs against rest-mass (the intra-nuclear distribution radii are
defined per nuclear species, so they require rest-mass projection
to assign mass to each species first --- structurally analogous to
the $(P_{14}, P_{16})$ one-way pair documented above).  No new
variable conflicts are expected since intra-nuclear $[L]$ scales
are categorically disjoint from atomic / internuclear $[L]$ scales.
Definitive cell-by-cell verification deferred to the next TX-A
audit.
\end{theorem}

\begin{observation}[axes are not independent]
\label{obs:axes_correlated}
Of the three axes (variable, dimension, transcendental), only the
transcendental axis is genuinely independent.  Cross-tabulating the
nineteen projections by axis-flag pattern, no projection adds a
dimension without also adding the variable that carries it (perfect
V/D correlation, 12/12); zero projections fall into the
$(\text{D-only})$ bucket.  The transcendental axis admits an
independence witness: vector-photon promotion (\S~\ref{sec:layer2},
projection~11) introduces $1/(4\pi)$ without adding any variable
or dimension.  Two projections (\texttt{spectral\_action},
\texttt{observation\_window}) move all three axes; these are the
``heavy'' projections, carrying the most physical content per step.
The newly added \texttt{breit\_retardation}
(\S~\ref{sec:proj_breit_retardation}) preserves the V/D correlation
(adds the variable $m_l/m_n$ together with the dimension energy)
and lives in the ring-preserving $\alpha^4 \cdot \mathbb{Q}$ tier on
the transcendental axis; it is therefore axis-pattern-equivalent to
the rest-mass projection at the V/D level, with the transcendental
class differing only by the leading $\alpha^4$ factor inherited from
its calibration sibling, the spectral action.
\end{observation}

This observation has structural force.  The variable axis and
dimension axis collapse to a single \emph{physical-content} axis
without information loss; the transcendental axis remains
independent.  The three-axis dictionary is therefore mildly
overcomplete on the V/D side but irreducible on the
transcendental side.  See Remark~\ref{rem:vd_correlation} in
\S~\ref{sec:dictionary} for the parallel statement at the
projection-row level.

% ===================================================================
\section{Falsifiable Prediction: Error Compounds with Projection Depth}
\label{sec:prediction}
% ===================================================================

The catalogue groups matches by projection depth: zero, one, two,
three.  Reading down the column of match precision gives:
\begin{itemize}\setlength\itemsep{0pt}
\item Zero-projection (Layer~1 only): all exact.
\item Single-projection: all exact at the symbolic level (modulo
truncation at finite $n_\text{max}$).
\item Two-projection: residual error $1$--$3\%$.
\item Three-projection: residual error $\approx 2$--$3\%$.
\end{itemize}

\begin{prediction}[Error compounds with projection depth]
\label{pred:error_depth}
For matches involving $k$ chained projections, the residual error
$\epsilon_k$ scales roughly as $\epsilon_k \sim k \cdot \epsilon_1$,
where $\epsilon_1 \sim 1\%$ is the per-projection truncation error
at finite~$n_\text{max}$.  Equivalently: every additional projection
introduces an independent $O(n_\text{max}^{-p})$ residual that does
not cancel against earlier projections.
\end{prediction}

This is a falsifiable structural claim about the framework, not a
phenomenological fit.  Existing data are consistent (LS-2 Lamb shift
at three projections, $-2.0\%$).  Tests for the prediction:
\begin{itemize}\setlength\itemsep{0pt}
\item Compute Schwinger anomalous moment $\alpha/(2\pi)$ via
two chained projections (vector photon $\circ$ spectral action).
The prediction sets a $1$--$2\%$ ceiling on the achievable accuracy
at fixed $n_\text{max}$.
\item Compute the H$^+_2$ ground-state energy (single projection,
Fock + spectral); the prediction sets the accuracy at the
single-projection symbolic limit.
\item Compute Casimir energy on $S^3 \times \mathbb{R}$ (zero
chained projections beyond Fock); should be exact symbolic.
\end{itemize}
A failure of the prediction (a multi-projection match found exact,
or a single-projection match found with residual error) would
indicate either an incorrect projection count for that observable
or a structural cancellation we have missed.

\begin{remark}[Refinement from Sprint Calc-P, 2026-05-09]
\label{rem:depth_refinement}
The H~1S static dipole polarizability calculation (Sprint Calc-P,
\S~\ref{sec:catalog} three-projection chain row) returned
$\alpha_p(1S) = \tfrac{9}{2}\,a_0^3$ \emph{exactly} via a depth-3
chain (Fock $\circ$ Wigner $3j$ $\circ$ Sturmian), residual~$=~0$.
This is structurally a counter-example to a strong reading of
Prediction~\ref{pred:error_depth}: depth alone does not bound
residual from below.

The mechanism is sharp.  The Dalgarno--Lewis solution
$F|1s\rangle = -\tfrac{1}{2}\, z(r{+}2)|1s\rangle$ has $p$-wave
radial part $(r^2 + 2r) e^{-r}$ that lies \emph{exactly} in the
span of the first two Sturmian basis functions $\{S_2, S_3\}$ at
$\lambda = Z$.  Combined with the framework's symbolic-exact
matrix-element machinery (rational arithmetic via
\texttt{geovac/hypergeometric\_slater.py} and
\texttt{geovac/dirac\_matrix\_elements.py}), the projection chain
produces the analytic answer with no truncation residual.

The conditional refinement: the depth-prediction holds for
observables whose analytic answer requires \emph{irrational
content} (continuous integration over a temporal/spectral parameter
per Paper~35 Prediction~1, transcendental ring per Paper~18, or
limits inaccessible at finite $n_\text{max}$).  For observables
whose analytic answer is a pure rational with a closed-form
representation that lies exactly within the framework's truncated
basis, the residual is zero at any depth.  Polarizability is the
first identified example of the latter class; static linear-response
observables that admit Dalgarno--Lewis-type closed forms are
candidates for the same exact-at-finite-$n_\text{max}$ behaviour.

The refined prediction: $\epsilon_k \sim k \cdot \epsilon_1$ for
observables in the irrational-analytic-answer class; $\epsilon_k =
0$ at sufficient $n_\text{max}$ for observables in the
rational-analytic-answer class with finite-basis representation.
The dividing line is empirically diagnosable from the
transcendental signature: rational-class observables have pure
$\mathbb{Q}$ transcendental signature, irrational-class observables
have at least one transcendental factor.
\end{remark}

% ===================================================================
\section{Connection to Paper 18 (Transcendentals) and Paper 32 (Spectral Triple)}
\label{sec:connections}
% ===================================================================

This paper does not supersede Paper~18 or Paper~32; it extends
them with a complementary axis.

\paragraph{Versus Paper~18.}  Paper~18 classifies transcendental
content of GeoVac observables (intrinsic, calibration, embedding,
flow, composition).  Paper~34 classifies the projections that
\emph{produce} that content.  The two classifications are
duals: Observation~\ref{obs:duality} maps each Paper~18 tier to
specific Paper~34 projections.  Reading direction~A
(transcendental $\to$ projection) lets you predict which projection
must be active given an observed transcendental.  Reading
direction~B (projection $\to$ transcendental) lets you predict the
transcendental signature of a calculation before running it.  The
audit memo of 2026-05-02 documents what happens when this dictionary
is unavailable: PSLQ searches against an arbitrary basis surface
near-misses that look like identifications, because the right basis
(determined by the projection) was never specified.

\paragraph{Versus Paper~32.}  Paper~32 constructs the GeoVac
spectral triple $(\mathcal{A}_{\text{GV}}, \mathcal{H}_{\text{GV}},
D_{\text{GV}})$ and audits it against the Connes axioms.  In that
language, Layer~1 of this paper is the algebra
$\mathcal{A}_{\text{GV}}$ and the Hilbert space
$\mathcal{H}_{\text{GV}}$, while Layer~2 is the family of
spectral functions, gauge sectors, and embedding maps that take
$(\mathcal{A}, \mathcal{H}, D)$ to physical observables.
Paper~34 is therefore the projection-side reading of the spectral
triple synthesis.

\paragraph{Versus Paper~31.}  Paper~31's universal/Coulomb
partition is the same observation phrased on the algebra/Dirac side:
features that depend only on $\mathcal{A}$ (e.g.\ angular sparsity)
are universal, features that depend on $D_{\text{GV}}$ (e.g.\
calibration $\pi$, the Hopf bundle's projection-$\pi$) are
potential-specific.  Paper~34's claim is the same in projection
language: zero-projection matches and Wigner-$3j$-only matches
transfer between potentials, while spectral-action and Fock-conformal
matches do not.

% ===================================================================
\section{Open Questions}
\label{sec:open}
% ===================================================================

\begin{enumerate}
\setlength\itemsep{0pt}

\item \textbf{Closure of the projection list.}  Sixteen projections
are now documented.  The thirteenth (Drake--Swainson asymptotic
subtraction) was added in sprint LS-4 to close the 2P Bethe-log
structural floor of LS-3.  The fourteenth (rest-mass projection) and
fifteenth (observation / temporal-window projection) were added in
sprint KG / Paper~35 (May~2026) after verifying that the bare
relativistic Klein-Gordon spectrum on $S^3 \times \mathbb{R}$ is
$\pi$-free in $\mathbb{Q}[\sqrt{d_1}, \sqrt{d_2}, \ldots]$ and that
$\pi$ enters the spectrum at exactly one step --- temporal
compactification.  These two projections were previously bundled
implicitly with other parameter-projection rows; the split was
forced by the Casimir computation $E_\text{Cas}^{S^3} = 1/240$
(spatial only, exact rational, no $\pi$) versus the Stefan-Boltzmann
$\pi^2/90$ (high-$T$ limit on $S^3 \times S^1_\beta$, $\pi$ from
Matsubara).  The sixteenth (two-body Dirac / Breit retardation,
\S~\ref{sec:proj_breit_retardation}) was added in the
precision-catalogue sprint of 2026-05-08, motivated by the
positronium 1S--2S framework-native residual at $+64.75$ ppm (the
sixth and structurally cleanest instance of the
multi-focal-composition wall).  Are there projections in the
framework that are still not yet in this list?  Two candidates that
may need to be added on review: (a)~the Mellin / heat-kernel
evaluation at fractional~$s$ (currently treated as part of
spectral-action), and (b)~the direction-resolved Hodge decomposition
of vector-photon edges (Paper~28 \S~vq\_sprint).  Decision pending.

\item \textbf{Quantitative form of Prediction~\ref{pred:error_depth}.}
The prediction states that error compounds with projection depth.
The current data are consistent with linear compounding
($\epsilon_k \sim k \cdot \epsilon_1$), but quadratic
($\sqrt{k} \cdot \epsilon_1$) or geometric
($\epsilon_1^k$) compounding are not yet ruled out.  A six-row
benchmark across $k = 0, 1, 2, 3, 4, 5$ would distinguish.

The LS-4 Lamb shift datapoint adds new information: the LS-4
chained projection (Fock + spectral action + Sturmian +
Drake--Swainson, $k=4$ projections) achieves $-0.39\%$ residual at
the sweet-spot $N=40$.  This is BELOW the linear prediction
$\epsilon_4 \sim 4\%$, but it arises from accidental
truncation/approximation cancellation rather than genuine reduction
of per-projection error.  At converged $N$, the LS-4 result
reverts to the LS-1 ceiling $-3.10\%$, consistent with the linear
compounding prediction taken over distinct uncorrelated
projections.  The structural reading: prediction~\ref{pred:error_depth}
applies to converged-basis residuals, not to transient
finite-basis cancellations.

\item \textbf{$K = \pi(B + F - \Delta) \approx 1/\alpha$.}
The Paper~2 conjecture is the most extreme multi-projection
coincidence in the catalogue: three structurally independent
projections (finite Casimir trace at $m=3$, infinite scalar
Dirichlet at $s = d_\text{max}$, Dirac mode degeneracy at the
spinor cutoff), all evaluated on the same $S^3$, summing to
$1/\alpha$ at $8.8 \times 10^{-8}$ relative residual.  Under
Prediction~\ref{pred:error_depth} a three-projection match
should carry $\sim 3\%$ error, not $10^{-7}$.  Either the
projection-depth model under-counts what's chained here, or
something structural we have not identified is forcing the
combination, or the match is a coincidence below the framework's
intrinsic resolution.  Working hypothesis WH5 (Paper~2 in
Observations status as of 2026-05-02) reads it as a coincidence;
this paper records the structural anomaly without resolving it.

\textit{Track TS-C update (2026-05-04).} The ``depth model
under-counts'' reading was tested by examining whether the three
chains share a fourth implicit projection.  They do not: $F$ is
derived from a Dirichlet zeta on the $(n,l)$ lattice with no Hopf
invocation (Phase~4F $\alpha$-J); $\Delta^{-1}$ is a
Camporesi--Higuchi Dirac mode degeneracy with no Hopf invocation
(Phase~4H SM-D); and the Hopf-quotient graph at $n_{\max} = 3$ does
not encode $K$ through any independent spectral invariant beyond
a Casimir relabeling (Phase~4B $\alpha$-D).  The outer $\pi$
prefactor of $K$ \emph{is} identified with the Hopf-base measure
(Sprint~A $\alpha$-PI), but this is already counted as the Hopf
projection in the depth-3 chain.  Under-counting is therefore not
the mechanism.  Combined with the K-CC clean negative on a unified
Connes--Chamseddine spectral-action source, twelve mechanisms total
have been eliminated.  The remaining live readings are (ii) and
(iii) above; the standing interpretation is (iii), per WH5.

\item \textbf{Unit-introduction asymmetry between Fock and Bargmann.}
Both Fock and Bargmann--Segal projections introduce energy as the
dimension; the former adds $Z$ and produces $\pi$ via
$\mathrm{Vol}(S^3)$, the latter adds $\hbar\omega$ and produces
no $\pi$ at all.  The asymmetry is captured by the HO/Coulomb
rigidity theorems (Paper~24, Paper~31) at the operator-order level
(first-order complex versus second-order real).  Whether the
projection-dictionary axis can independently predict the asymmetry
without invoking operator order is open.

\item \textbf{Pure-axis projections.}  The current dictionary contains
no projection that moves only the dimension axis (e.g., a
unit-rescaling that introduces a length scale without a coordinate
variable).  Whether such projections are structurally absent from the
GeoVac framework, or merely absent from the current implementation,
is open.  Two candidate additions surfaced by the TX-A axiomatization
sprint: (a)~a gauge-choice projection (Coulomb vs Lorenz) that would
sit in the pure-transcendental bucket alongside vector-photon, (b)~a
spin-statistics projection that splits boson/fermion Matsubara
content out of observation/temporal-window.  Neither is forced by
current data.

\end{enumerate}

% ===================================================================
\section{Conclusion}
\label{sec:conclusion}
% ===================================================================

GeoVac is a two-layer framework.  Layer~1 is a bare combinatorial
graph: dimensionless integers, rational matrix elements, no
physics, no $\pi$.  Layer~2 is a finite family of named projections;
each adds specific physical variables, specific physical
dimensions, and specific transcendental content.  Every physical
quantity in any GeoVac calculation enters through exactly one
projection in Layer~2 and is fully traceable to it.

The dictionary in this paper makes that statement operational.
Nineteen projections are listed, each tagged by variables, dimension,
and transcendental signature.  An empirical matches catalogue uses
these tags to organise the framework's known machine-precision
agreements with standard physics values.  Sprint TX-A's
dimensional axiomatization (\S~\ref{sec:axiomatization}) verified
three structural theorems on this dictionary: dimensional consistency
holds across all worked examples, weak minimality holds at the
signature level (with the strong form failing on two non-commuting
pairs --- the Paper~2 \texttt{hopf}--\texttt{spinor\_lift} pair under
bundle-non-trivial composition, and the
\texttt{rest\_mass}--\texttt{breit\_retardation} pair under
variable-refinement dependence).  The 16$\times$16 composition table
expanded from the previous 15$\times$15 reflects the new sixteenth
projection (two-body Dirac / Breit retardation, added 2026-05-08);
under this expansion Drake--Swainson is no longer the unique
projection with input\_layer = output\_layer = 2, with Breit
retardation joining as the second.  Both projections continue to
account for all asymmetric (one-way) cells in the 16$\times$16
table.  The 2026-05-09 addition of three intra-nuclear projections
(charge-density / Foldy--Friar, magnetization-density / Zemach,
tensor multipole) extends the dictionary further to nineteen, and
sharpens the $L$--$E$ asymmetry of \S~\ref{sec:base_units}: $[L]$
admits six scalar instances plus one rank-2 tensor instance, while
$[E]$ admits four.  A boundary subsection
(\S~\ref{sec:layer2_boundary}) names the structural distinction
between intrinsic graph-derived variables (which the framework
projects natively) and Layer-2 inputs (external classical fields,
multi-loop QED coefficients, lattice-QCD-derived radii); the same
day's external-field hunting verified that uniform Zeeman / Stark
fields are Layer-2, not new projections.  T3 expansion to
19$\times$19 is flagged for the next TX-A audit.  The
field-theoretic UV/IR ordering anomaly between \texttt{spectral\_action}
and \texttt{observation\_window} is real physics but is \emph{not}
encoded by the current coarse layer-grading (Track TS-C, 2026-05-04);
refining the output\_layer index to distinguish UV-cutoff, thermal,
and flow Layer-2 outputs is the natural next step.  Of the three axes, only the
transcendental axis is genuinely independent --- variable and
dimension are perfectly correlated.  A falsifiable structural
prediction (\S~\ref{sec:prediction}) emerges from the depth-grouped
catalogue: error compounds with projection depth.  The prediction is
consistent with current data, including the
208/208-confirmation panel of Paper~35 Prediction~1
(Sprint TX-B, May~2026), and refinable by future calculations.

This paper introduces no new computational result.  Its contribution
is a vocabulary that the framework has been needing.  With the
two-layer view in place, the question ``where does $\pi$ come from
in this observable?'' has a definite answer for every observable
GeoVac produces.  The same is true for $\alpha$, $Z$, $\hbar\omega$,
$\Lambda$, $\zeta(3)$, Catalan's $G$, and any other quantity that
enters from outside the bare graph.

% ===================================================================
\begin{thebibliography}{99}
% ===================================================================

\bibitem{paper0}
GeoVac Project, ``Geometric Packing Construction,'' Paper~0 (2025).

\bibitem{paper2}
GeoVac Project, ``The fine structure constant from Hopf bundle
spectral invariants,'' Paper~2 (2025; observation status, May~2026).

\bibitem{paper7}
GeoVac Project, ``The Dimensionless Vacuum: $S^3$ Conformal
Equivalence,'' Paper~7 (2025).

\bibitem{paper14}
GeoVac Project, ``Structurally Sparse Qubit Hamiltonians from Natural
Geometry,'' Paper~14 (2025).

\bibitem{paper15}
GeoVac Project, ``Level~4 Geometry: Molecule-Frame Hyperspherical
Coordinates,'' Paper~15 (2025).

\bibitem{paper17}
GeoVac Project, ``Composed Geometries,'' Paper~17 (2025).

\bibitem{paper18}
GeoVac Project, ``Spectral-Geometric Exchange Constants,'' Paper~18
(2026).

\bibitem{paper19}
GeoVac Project, ``Coupled Composition: Two-Center Integrals and
Sturmian--Wulfman Connection,'' Paper~19 (2026).

\bibitem{paper20}
GeoVac Project, ``Resource Benchmarks for Quantum Computing,''
Paper~20 (2026).

\bibitem{paper22}
GeoVac Project, ``Potential-Independent Angular Sparsity Theorem,''
Paper~22 (2026).

\bibitem{paper23}
GeoVac Project, ``Nuclear Shell Model on the GeoVac Framework,''
Paper~23 (2026).

\bibitem{paper24}
GeoVac Project, ``The Bargmann--Segal Lattice: $\pi$-Free
Discretization of the 3D Harmonic Oscillator on $S^5$,'' Paper~24
(2026).

\bibitem{paper25}
GeoVac Project, ``The Hopf Graph as a Lattice Gauge Structure,''
Paper~25 (2026).

\bibitem{paper27}
GeoVac Project, ``Entropy as Projection Artifact,'' Paper~27 (2026).

\bibitem{paper28}
GeoVac Project, ``QED on $S^3$: Spectral Geometry of Perturbative
Transcendentals,'' Paper~28 (2026).

\bibitem{paper30}
GeoVac Project, ``SU(2) Wilson Lattice Gauge Structure on the GeoVac
Hopf Graph,'' Paper~30 (2026).

\bibitem{paper31}
GeoVac Project, ``Universal/Coulomb Partition: The $\mathcal{A}/D$
Split of the GeoVac Spectral Triple,'' Paper~31 (2026).

\bibitem{paper32}
GeoVac Project, ``The GeoVac Spectral Triple: A Synthesis Paper,''
Paper~32 (2026).

\bibitem{paper33}
GeoVac Project, ``The 1+6+1 Partition of QED Selection Rules
on $S^3$,'' Paper~33 (2026).

\bibitem{paper35}
GeoVac Project, ``Time as Projection: The Klein-Gordon Spectrum on
$S^3 \times \mathbb{R}$ and the Algebraic / Observation Split,''
Paper~35 (2026).

\bibitem{ls1_memo}
GeoVac Project, ``LS-1 Sprint: Hydrogen Lamb Shift on the Dirac-S$^3$
Framework,'' \texttt{debug/ls1\_lamb\_shift\_memo.md} (May 2026).

\bibitem{ls2_memo}
GeoVac Project, ``LS-2 Sprint: Bethe Logarithm from GeoVac Spectral
Machinery,'' \texttt{debug/ls2\_bethe\_log\_memo.md} (May 2026).

\bibitem{ls3_memo}
GeoVac Project, ``LS-3 Sprint: Regularized Bethe Logarithm via
Acceleration Form,'' \texttt{debug/ls3\_bethe\_log\_regularized\_memo.md}
(May 2026).

\bibitem{ls4_memo}
GeoVac Project, ``LS-4 Sprint: Bethe Logarithm via Schwartz/Drake--Swainson
Regularization,'' \texttt{debug/ls4\_bethe\_log\_drake\_memo.md} (May 2026).

\bibitem{vp1_memo}
GeoVac Project, ``VP-1 Sprint: Vector-Photon Promotion on the
Graph-Native QED Pipeline,'' \texttt{debug/vp1\_vector\_graph\_native\_memo.md}
(May 2026).

\bibitem{vp2_memo}
GeoVac Project, ``VP-2 Sprint: Topology-Specific Projection Theory,''
\texttt{debug/vp2\_topology\_projections\_memo.md} (May 2026).

\bibitem{audit_memo}
GeoVac Project, ``Curve-Fit Audit (2026-05-02),''
\texttt{docs/curve\_fit\_audit\_memo.md} (May 2026).

\bibitem{fock1935}
V.~Fock, ``Zur Theorie des Wasserstoffatoms,''
\textit{Z.\ Phys.}\ \textbf{98}, 145 (1935).

\bibitem{bargmann1961}
V.~Bargmann, ``On a Hilbert space of analytic functions and an
associated integral transform, Part~I,'' \textit{Comm.\ Pure Appl.\
Math.}\ \textbf{14}, 187 (1961).

\bibitem{chamseddine_connes2010}
A.~H.~Chamseddine and A.~Connes, ``The spectral action principle in
noncommutative geometry,'' \textit{Phys.\ Rev.\ D}\ \textbf{83},
045001 (2011).

\bibitem{marcolli_vansuijlekom2014}
M.~Marcolli and W.~D.~van~Suijlekom, ``Gauge networks in
noncommutative geometry,'' \textit{J.\ Geom.\ Phys.}\ \textbf{75},
71 (2014); arXiv:1301.3480.

\bibitem{drake_swainson1990}
G.~W.~F.~Drake and R.~A.~Swainson, ``Bethe logarithms for hydrogenic
atoms,'' \textit{Phys.\ Rev.\ A}\ \textbf{41}, 1243 (1990).

\bibitem{camporesi_higuchi1996}
R.~Camporesi and A.~Higuchi, ``On the eigenfunctions of the Dirac
operator on spheres and real hyperbolic spaces,''
\textit{J.\ Geom.\ Phys.}\ \textbf{20}, 1 (1996).

\end{thebibliography}

\end{document}
